\documentclass[11pt]{article}
\usepackage[utf8]{inputenc}
\usepackage[a4paper,margin=2.6cm]{geometry}
\usepackage{amsmath,amssymb,amsthm,mathtools,booktabs,array}
\usepackage{graphicx}
\usepackage{enumitem}
\usepackage{fancyvrb}
\usepackage{fvextra}
% generated by gen_appA.py
\DeclareUnicodeCharacter{00A7}{\S{}}
\DeclareUnicodeCharacter{00B2}{\ensuremath{^2}}
\DeclareUnicodeCharacter{00B3}{\ensuremath{^3}}
\DeclareUnicodeCharacter{00B7}{\ensuremath{\cdot}}
\DeclareUnicodeCharacter{02E2}{\ensuremath{^s}}
\DeclareUnicodeCharacter{0302}{\^{}}
\DeclareUnicodeCharacter{0303}{\~{}}
\DeclareUnicodeCharacter{03A3}{\ensuremath{\Sigma}}
\DeclareUnicodeCharacter{03B2}{\ensuremath{\beta}}
\DeclareUnicodeCharacter{03B3}{\ensuremath{\gamma}}
\DeclareUnicodeCharacter{03B5}{\ensuremath{\varepsilon}}
\DeclareUnicodeCharacter{03B6}{\ensuremath{\zeta}}
\DeclareUnicodeCharacter{03BB}{\ensuremath{\lambda}}
\DeclareUnicodeCharacter{03BC}{\ensuremath{\mu}}
\DeclareUnicodeCharacter{03BD}{\ensuremath{\nu}}
\DeclareUnicodeCharacter{03C0}{\ensuremath{\pi}}
\DeclareUnicodeCharacter{03C1}{\ensuremath{\rho}}
\DeclareUnicodeCharacter{03C8}{\ensuremath{\psi}}
\DeclareUnicodeCharacter{1DA0}{\ensuremath{^f}}
\DeclareUnicodeCharacter{2013}{--}
\DeclareUnicodeCharacter{2014}{---}
\DeclareUnicodeCharacter{2026}{\ensuremath{\ldots}}
\DeclareUnicodeCharacter{2032}{\ensuremath{{}^{\prime}}}
\DeclareUnicodeCharacter{2074}{\ensuremath{^4}}
\DeclareUnicodeCharacter{2077}{\ensuremath{^7}}
\DeclareUnicodeCharacter{2078}{\ensuremath{^8}}
\DeclareUnicodeCharacter{2080}{\ensuremath{_0}}
\DeclareUnicodeCharacter{2081}{\ensuremath{_1}}
\DeclareUnicodeCharacter{2082}{\ensuremath{_2}}
\DeclareUnicodeCharacter{2085}{\ensuremath{_5}}
\DeclareUnicodeCharacter{2102}{\ensuremath{\mathbb{C}}}
\DeclareUnicodeCharacter{2115}{\ensuremath{\mathbb{N}}}
\DeclareUnicodeCharacter{211D}{\ensuremath{\mathbb{R}}}
\DeclareUnicodeCharacter{2192}{\ensuremath{\to}}
\DeclareUnicodeCharacter{21A6}{\ensuremath{\mapsto}}
\DeclareUnicodeCharacter{2200}{\ensuremath{\forall}}
\DeclareUnicodeCharacter{2203}{\ensuremath{\exists}}
\DeclareUnicodeCharacter{2208}{\ensuremath{\in}}
\DeclareUnicodeCharacter{2211}{\ensuremath{\Sigma}}
\DeclareUnicodeCharacter{2212}{\ensuremath{-}}
\DeclareUnicodeCharacter{221A}{\ensuremath{\surd}}
\DeclareUnicodeCharacter{221E}{\ensuremath{\infty}}
\DeclareUnicodeCharacter{2227}{\ensuremath{\wedge}}
\DeclareUnicodeCharacter{2229}{\ensuremath{\cap}}
\DeclareUnicodeCharacter{222A}{\ensuremath{\cup}}
\DeclareUnicodeCharacter{222B}{\ensuremath{\int}}
\DeclareUnicodeCharacter{2248}{\ensuremath{\approx}}
\DeclareUnicodeCharacter{2264}{\ensuremath{\leq}}
\DeclareUnicodeCharacter{2265}{\ensuremath{\geq}}
\DeclareUnicodeCharacter{22A4}{\ensuremath{\top}}
\DeclareUnicodeCharacter{22EF}{\ensuremath{\cdots}}
\DeclareUnicodeCharacter{27E8}{\ensuremath{\langle}}
\DeclareUnicodeCharacter{27E9}{\ensuremath{\rangle}}
\DeclareUnicodeCharacter{1D426}{\ensuremath{\mathbf{m}}}
% characters of the comparator files
\DeclareUnicodeCharacter{0304}{\={}}
\DeclareUnicodeCharacter{2075}{\ensuremath{^5}}
\DeclareUnicodeCharacter{207B}{\ensuremath{^-}}

\usepackage{xcolor}
\usepackage[colorlinks=true,linkcolor=blue!50!black,citecolor=green!40!black,urlcolor=blue!50!black]{hyperref}
\hypersetup{pdftitle={On the proportions of simple and of distinct zeros of the Riemann zeta function},pdfauthor={Oliver D'Souza},pdfkeywords={Riemann zeta function; simple zeros; distinct zeros; Lean 4; formal verification}}
% Macros for Paper A (notation follows Alpoge--Furman, arXiv:2608.13637v2)
\newtheorem{theorem}{Theorem}[section]
\newtheorem{lemma}[theorem]{Lemma}
\newtheorem{proposition}[theorem]{Proposition}
\newtheorem{corollary}[theorem]{Corollary}
\theoremstyle{definition}
\newtheorem{definition}[theorem]{Definition}
\theoremstyle{remark}
\newtheorem{remark}[theorem]{Remark}
\renewcommand{\Re}{\operatorname{Re}}
\renewcommand{\Im}{\operatorname{Im}}
\newcommand{\tr}{\operatorname{tr}}
\newcommand{\rank}{\operatorname{rank}}
\newcommand{\HS}[1]{\lVert #1\rVert_{\mathrm{HS}}}
\newcommand{\abs}[1]{\lvert #1\rvert}
\newcommand{\norm}[1]{\lVert #1\rVert}
\newcommand{\floor}[1]{\lfloor #1\rfloor}
\DeclareMathOperator{\sinc}{sinc}
\DeclareMathOperator{\ch}{ch}
\DeclareMathOperator{\mg}{mg}
\DeclareMathOperator{\Slack}{Slack}
\newcommand{\R}{\mathbb R}
\newcommand{\C}{\mathbb C}
\newcommand{\Z}{\mathbb Z}
\newcommand{\cO}{\mathcal O}
\newcommand{\cT}{\mathcal T}
\newcommand{\cR}{\mathcal R}
\newcommand{\cB}{\mathcal B}
\newcommand{\bm}{\mathbf m}
\newcommand{\eps}{\varepsilon}
% Release metadata: the address and the tag of the repository, each in ONE place.
\newcommand{\repourl}{\url{https://github.com/Oliverds321/simple-zeros}}
\newcommand{\repotag}{\texttt{v1.0}}
% Text that the author still has to confirm. To switch the marker off: \newcommand{\draftnote}[1]{#1}
\newcommand{\draftnote}[1]{#1\textbf{ [DRAFT, author to confirm]}}
\newcommand{\tsp}{\mathsf T}


\title{On the proportions of simple and of distinct zeros\\ of the Riemann zeta function}
\author{Oliver D'Souza\thanks{\copyright~2026 Oliver D'Souza. This paper is licensed under CC~BY~4.0.}}
\date{Release version, October 2026 (repository tag \repotag)}

\begin{document}
\maketitle

\begin{abstract}
We prove, without hypotheses on the zeros and with computer-assisted steps, that, as $T\to\infty$, at least $0.676102666$ of the zeros
$\rho$ of $\zeta(s)$ with $T<\Im\rho\le2T$, counted with multiplicity, are simple, and at least $0.838051333$ of them are distinct. The
largest previous unconditional bounds that we located in the refereed and arXiv literature are, truncated, $0.67250077$ for simple zeros, which is the bound for
simple zeros on the critical line, and $0.83625038$ for distinct zeros. Unrefereed computer-assisted claims for simple zeros on the line at
bandwidth one reach $0.67349239$; the three largest of them use the fixed-test-function pair-correlation formula, as Lamzouri does, rather
than the inputs of Alp\"oge and Furman used here. Two unrefereed works claim considerably more by other means, and we have not verified
them.
The argument extends the rank--trace method of Alp\"oge and Furman. With a smaller computer-certified local inequality ($K=5$ consecutive
zeros instead of $7$) it gives the slightly smaller pair $0.675158622$ and $0.837579311$. For this pair the whole proof, the local
certificate included, is checked by the kernel of Lean~4 with no hypotheses. For the $K=7$ pair and for the two statistics below, everything
except the local certificate (the analytic and algebraic steps, the majorant certificate and the closing arithmetic) is checked by the
kernel, and the local certificate enters as a single displayed hypothesis, verified outside Lean by interval arithmetic. Two further
statistics are recorded. At least $0.6733736895$ of the zeros are simple and on the critical line, by a method that others published first;
this is not a record, since for this statistic six unrefereed works announce larger values, from $0.673399$ to $0.67349239$, one of them
($0.6734$) as a hypothesis-free Lean theorem, whereas our formal statement carries a certificate hypothesis. At least $0.8879195$ of the
zeros are simple or on the line, which is Lamzouri's formula evaluated at $0.6733736895$. The results are lower bounds only and say nothing
about the Riemann hypothesis itself. None of them has been independently refereed, and the Lean development has not been rebuilt outside
this project.
\end{abstract}

\tableofcontents

\section{Introduction}\label{sec:intro}

\subsection{Results}\label{ssec:results}
Write a nontrivial zero of $\zeta(s)$ as $\rho=\beta+i\gamma$ and let $m_\rho\ge1$ be its multiplicity. As in Alp\"oge and Furman~\cite{AF}, for
$0\le T_1<T_2$ let
\begin{align*}
N(T_1,T_2)&:=\sum_{T_1<\gamma\le T_2}m_\rho &&\text{(zeros, counted with multiplicity)},\\
N^d(T_1,T_2)&:=\#\{\rho:\ T_1<\gamma\le T_2\} &&\text{(distinct zeros)},\\
N^s(T_1,T_2)&:=\#\{\rho:\ T_1<\gamma\le T_2,\ m_\rho=1\} &&\text{(simple zeros)},\\
N_0^s(T_1,T_2)&:=\#\{\rho:\ T_1<\gamma\le T_2,\ \beta=\tfrac12,\ m_\rho=1\} &&\text{(simple zeros on the critical line)},
\end{align*}
and let $N^{\mathrm{sc}}(T_1,T_2)$ be the number of zeros with $T_1<\gamma\le T_2$ that lie on the critical line, counted with multiplicity, plus
the number of those off the line that are simple. Thus $N_0^s\le N^s\le N^d\le N$ and $N_0^s\le N^{\mathrm{sc}}\le N$. All ratios below have
$N(T,2T)$, which counts multiplicity, in the denominator.

\begin{theorem}[Simple and distinct zeros]\label{thm:main}
As $T\to\infty$,
\[
N^s(T,2T)\ \ge\ (0.676102666-o(1))\,N(T,2T),\qquad N^d(T,2T)\ \ge\ (0.838051333-o(1))\,N(T,2T).
\]
With the smaller certificate of Section~\ref{sec:cert} \textup{(}$K=5$\textup{)} the same argument gives $0.675158622$ and $0.837579311$.
\end{theorem}

The constants are $\Sigma=(H+\frac{a_2}2-\nu)/(1-a_1+\frac{a_2}2)$ and $D=\frac12(1+\Sigma)$, where $H=H(\cos1.6s)=0.671981551000\ldots$ is the
Alp\"oge--Furman constant of the window $\cos(1.6s)$ and $(a_1,a_2,\nu)$ are the rational claims of a certified local inequality
(Theorem~\ref{thm:SigmaD}). For $K=7$, $\Sigma=0.67610266696413\ldots$ and $D=0.83805133348206\ldots$; for $K=5$, $\Sigma=0.67515862219923\ldots$
and $D=0.83757931109961\ldots$; the ninth digits above are truncated. In Lean~4~\cite{Lean4}, the $K=5$ pair is proved with no hypotheses and the
$K=7$ pair with its certificate as the one hypothesis (\S\ref{ssec:formal}).

\begin{theorem}[Simple zeros on the critical line]\label{thm:S}
As $T\to\infty$, $N_0^s(T,2T)\ge(0.6733736895-o(1))\,N(T,2T)$.
\end{theorem}

\begin{corollary}[Simple or on the line]\label{cor:SC}
As $T\to\infty$, $N^{\mathrm{sc}}(T,2T)\ge(0.8879195-o(1))\,N(T,2T)$.
\end{corollary}

\begin{table}[ht]\centering\small
\caption{Lower bounds for $\liminf_{T\to\infty}X(T,2T)/N(T,2T)$. ``Previous'' lists refereed or arXiv values, and for $N_0^s$ also the
range of unrefereed values; the unrefereed values are discussed in \S\ref{ssec:others}. The last column is explained in \S\ref{ssec:formal}.
Lamzouri's and Wang's values are stated for $0<\gamma\le T$; a lower bound for every dyadic interval $(T,2T]$ implies the same lower bound for
$0<\gamma\le T$. Wang's values are our truncations of the bounds $C_0+\delta_0$ and $C_1+\delta_0/2$ that it states (\S\ref{ssec:others});
Lamzouri's $0.8876200$ is its formula at the Montgomery--Taylor constant, evaluated by us (it states $88.76\%$).}\label{tab:results}
\begin{tabular}{@{}>{\raggedright\arraybackslash}p{2.4cm}>{\raggedright\arraybackslash}p{4.4cm}>{\raggedright\arraybackslash}p{3.5cm}>{\raggedright\arraybackslash}p{4.2cm}@{}}
\toprule
$X$ & previous & this paper & status of the proof in Lean \\
\midrule
$N^d$ (distinct) & $\frac56$ and $0.8362503$ \cite{AF}; $0.8362$ \cite{Lamzouri}; $0.83625038$ \cite{Wang24167} &
$\mathbf{0.838051333}$ ($K=7$)\newline $\mathbf{0.837579311}$ ($K=5$) &
$K=7$: certificate as hypothesis\newline $K=5$: no hypothesis (the certificate is replayed by the kernel)\\[3pt]
$N^s$ (simple) & $0.6725007$ \cite{AF,Lamzouri}; $0.67250077$ \cite{Wang24167}; all implied by $N_0^s\le N^s$ &
$\mathbf{0.676102666}$ ($K=7$)\newline $\mathbf{0.675158622}$ ($K=5$) & as for $N^d$\\[3pt]
$N^{\mathrm{sc}}$ (simple or on the line) & $0.8876200$ \cite{Lamzouri} & $0.8879195$ (Cor.~\ref{cor:SC}) & certificate as hypothesis\\[3pt]
$N_0^s$ (simple, on the line) & $0.6725007$ \cite{AF,Lamzouri}; $0.67250077$ \cite{Wang24167}; unrefereed: six values from $0.673399$ to
$0.67349239$, one announced as a hypothesis-free Lean theorem (\S\ref{ssec:others}) & $0.6733736895$ (Thm~\ref{thm:S}; not a record) & certificate as hypothesis\\
\bottomrule
\end{tabular}
\end{table}

For distinct zeros Theorem~\ref{thm:main} improves Wang's bound~\cite{Wang24167}, $0.83625038$ truncated (\S\ref{ssec:others}), by
$1.80\cdot10^{-3}$ ($1.32\cdot10^{-3}$ with $K=5$).
For simple zeros there is, to our knowledge, no earlier bound specific to the statistic: since a simple zero on the line is simple, every lower
bound for $N_0^s$ is one for $N^s$, and Theorem~\ref{thm:main} exceeds the largest refereed or arXiv value that we located, $0.67250077$~\cite{Wang24167}, by
$3.60\cdot10^{-3}$, and the largest unrefereed bandwidth-one claim for $N_0^s$ that we know of ($0.67349239$, \S\ref{ssec:others}) by
$2.61\cdot10^{-3}$; larger unrefereed claims obtained by other methods are discussed there.
Both values of Theorem~\ref{thm:main} lie below what is known under the Riemann hypothesis: $0.6792$~\cite{CGdL} and $19/27$~\cite{BHB} for
simple zeros, $0.8477$~\cite{CGdL} for distinct zeros.

Theorem~\ref{thm:main} is the result that this paper claims as new, as far as we have found in the public record: apart from the unrefereed
claims of Yang and Yang by other methods (\S\ref{ssec:others}), we found no earlier bound for $N^s$ or $N^d$ beyond the standard
consequences of a bound for $N_0^s$ (\S\ref{ssec:new}). We read $24$ related works at source and made three web searches; this says only
that we did not find such a bound, not that none exists. Theorem~\ref{thm:S} is not a record: for simple zeros on the line, six unrefereed
works announce larger values, from $0.673399$ to $0.67349239$, and one of them, Zhaoyi-Tian's $0.6734$, is announced as a hypothesis-free
Lean theorem, while the formal statement of Theorem~\ref{thm:S} carries a certificate hypothesis (\S\ref{ssec:others},
\S\ref{ssec:formal}).

\subsection{What is new}\label{ssec:new}
Alp\"oge and Furman~\cite{AF} restrict Weil's explicit formula to a finite family of modulated copies of a window and obtain a real symmetric
matrix $G$ whose trace and Hilbert--Schmidt norm are known asymptotically from the prime side, and whose zero side splits into a positive
semidefinite part $P$ from the zeros on the line and a part $Q$ from the off-line pairs with at most one positive eigenvalue per pair.
A rank--trace inequality then gives $N_0^s\ge(2-R(\psi)-o(1))N$ for every admissible window $\psi$.

The main results of this paper, Theorem~\ref{thm:main}, keep more of the matrix. For $N^s$ and $N^d$ a simple off-line pair counts twice in the
numerator, and this is exactly what its negative eigenvalue costs in the rank--trace accounting (the ``simple pair'' entry in Step~3 of the
proof of Theorem~\ref{thm:SigmaD} is $2-2a_1-c_*=0$). Off-line pairs and zeros of multiplicity at least three can therefore be removed at no loss
(Lemma~\ref{lem:removal}), and what remains is an inequality about the zeros on the line of multiplicity one or two, the on-line spectral
condition (Theorem~\ref{thm:OLL}). We prove it with four ingredients:
\begin{enumerate}[label=(\roman*)]
\item an exact identity for the slack of the rank--trace inequality on a set of on-line zeros (Lemma~\ref{lem:onslack}), which isolates the
Gram energy among \emph{all} of them, whatever their multiplicities;
\item a local inequality over all patterns of multiplicities $\bm\in\{1,2\}^K$ of $K$ consecutive zeros, certified by computer
(Definition~\ref{def:LIm}, Section~\ref{sec:cert}), and its aggregation along the ordered zeros (Proposition~\ref{prop:agg});
\item a localisation lemma for the eigenvalues above a threshold (Theorem~\ref{thm:loc}), which reduces the spectral condition to a set of
zeros with a close neighbour, where it is paid for by the zeros themselves (Lemma~\ref{lem:tight});
\item a majorant of Beurling--Selberg type for the window, with Fourier transform supported in $[-1,1]$, which bounds the largest eigenvalue on
the remaining separated set (Lemma~\ref{lem:room}).
\end{enumerate}
A split of the zeros into close pairs and a separated remainder, whose Gram norm is bounded by a sinc-square Fourier majorant, was published
earlier for simple zeros on the line, by \mbox{gfreund123} in pull request~\#888 to \mbox{GettysburgResearch/riemann}
(14~September~2026)~\cite{Gettysburg888} [V$\sim$]; ingredients (iii) and (iv) use a device of the same kind for a different spectral
condition, and the priority is theirs. Related devices appear, for $N_0^s$, in \mbox{trmdy}'s second campaign~\cite{trmdy} (off-line pairs priced as virtual on-line
doubles, 11~August~2026, left open there) and in \mbox{RakitinD}'s exact decomposition of the slack with off-line and multiple-zero
residuals~\cite{Rakitin} [V$\sim$]; neither is used there for $N^s$ or $N^d$.
In one line, for $X$ the number of simple zeros: with $\Slack_X:=X-2\tr G+\HS G^2$, the prime side gives
$X\ge(H(\psi)-o(1))N+\Slack_X$, exactly as in~\cite{AF}; the four ingredients show
\[
\Slack_X\ \ge\ a_1X+\tfrac{a_2}2(N-X)-\nu N-o(N),
\]
where $a_1,a_2,\nu$ are the certified constants; and solving the two inequalities for $X$ gives
$X\ge\Sigma N-o(N)$ with $\Sigma=(H+\frac{a_2}2-\nu)/(1-a_1+\frac{a_2}2)$. The bound of~\cite{AF} corresponds to $a_1=a_2=\nu=0$.


Apart from the unrefereed claims of Yang and Yang discussed in \S\ref{ssec:others}, we know of no public work that treats the proportions
of simple zeros or of distinct zeros beyond the standard consequences of a bound $S$ for $N_0^s$: the bound $S$ itself for $N^s$ (since
$N_0^s\le N^s$), and for $N^d$ the value $\frac12(1+S)$ that the chain of~\cite[Theorem~A(ii)]{AF} gives for $S=H(\psi)$ and that of
Theorem~\ref{thm:S} gives for its constant (Remark~\ref{rem:Sdist}; for Theorem~\ref{thm:S} this is $0.8366868$).

Theorem~\ref{thm:S} and Corollary~\ref{cor:SC} are secondary, and we state their status plainly in \S\ref{ssec:others}. Corollary~\ref{cor:SC} is
Lamzouri's formula~\cite{Lamzouri} for zeros that are simple or on the line, evaluated at the constant of Theorem~\ref{thm:S}; it is not an
independent advance. For Theorem~\ref{thm:S} the method appeared publicly before our work on it, and several unrefereed values exceed ours.

\subsection{What the results are not}\label{ssec:not}
The theorems are lower bounds only. They do not show that the remaining zeros are multiple or off the line; they are merely not reached by the
certificates. Theorem~\ref{thm:main} counts simple zeros \emph{whether or not} they lie on the critical line: a configuration in which a
positive proportion of the zeros are simple off-line pairs is consistent with it, and it gives no information on how many zeros are on the
line. It says nothing about the Riemann hypothesis. The proportion of zeros that are simple \emph{and} on the line is not improved by the route
of Theorem~\ref{thm:main} (Remark~\ref{rem:notS}); for that statistic this paper has only Theorem~\ref{thm:S}. The analytic inputs are those
of~\cite{AF}; in particular no zero-density estimate, zero-free region or mollifier is used. As observed in~\cite[\S1.4]{AF}, these inputs
are insensitive to $o(N)$ zeros off the line, and the local inequalities added here concern only the kernel $k_\psi$. The local inequalities behind the constants are
verified by computer; how, and at what level of trust, is stated in \S\ref{ssec:formal} and Section~\ref{sec:cert}.

\subsection{Relation to other work}\label{ssec:others}
The method of this paper is that of Alp\"oge and Furman~\cite{AF}, whose argument makes Montgomery's pair-correlation deduction unconditional.
Lamzouri~\cite{Lamzouri} gave a second proof of their bound by a Hilbert-space inequality, and obtained the only previous value, $0.8876200$, for
the proportion of zeros that are simple or on the line. Wang~\cite{Wang24167} raised the constants, ``by refining the method of Lamzouri'', to
$C_0+\delta_0$ for simple zeros on the line and $C_1+\delta_0/2$ for distinct zeros, with $C_0=\frac32-\frac1{\sqrt2}\cot\frac1{\sqrt2}=0.67250\ldots$,
$C_1=\frac12(C_0+1)=0.83625\ldots$ and $\delta_0=6.66624\ldots\cdot10^{-8}$, for $0<\gamma\le T$; truncated, these are $0.67250077$ and
$0.83625038$ (our evaluation). We have read its abstract and the statement of its main theorem [V$\sim$].
Santib\'a\~nez-Leal~\cite{SantibanezLeal} (24~September~2026, unrefereed) sharpens one step of Wang's argument to $0.6725007995$ and
$0.8362503997$ [V$\sim$]. Before 2026 the unconditional bounds for distinct zeros included Farmer's
$0.63952$~\cite{Farmer95} and Wu's $0.66036$~\cite{Wu12}, the latter cited as the record in~\cite{AF}; Lamzouri~\cite{Lamzouri} attributes a
value above $70\%$ to Ki and Lee~\cite{KiLee12}, which we have not been able to check.

\paragraph{Simple zeros on the line: others' work and what is ours.}
In August and September 2026 many unrefereed, computer-assisted claims above~\cite{AF} for $N_0^s/N$ appeared on GitHub and Zenodo. In this
paragraph and throughout the paper, statements about other work carry the tag [V$\sim$] when we read them only through the authors'
descriptions or repository files, and [M] when they are our inference; we have built or run none of these works. As far as their descriptions show
[V$\sim$], most of them and Theorem~\ref{thm:S} share one mechanism, published first, as far as we found, by \mbox{ainta}~\cite{ainta} on 11~August~2026: the stability
form of the rank--trace inequality, a local $K$-point functional with a length penalty, an offset-averaged block aggregation, and an interval
branch-and-bound. The step from the local inequality to the global bound differs between works (a piecewise-linear subaction in
Devine~\cite{Devine22066689}, a cap-one block assembly in Zhaoyi-Tian~\cite{ZhaoyiTian}, a Bellman coboundary in
tawanerguo-cn~\cite{TawanRepo}, a fusion polyhedron in \mbox{RakitinD}~\cite{Rakitin}, a separation-or-pair transfer in pull
request~\#888~\cite{Gettysburg888}); we have not compared them [M].
Position weights and windows other than a single cosine appeared the same day in \mbox{trmdy}~\cite{trmdy}, and also in
\mbox{sxuff}~\cite{sxuff} (11~August), \mbox{AMTOPA}~\cite{AMTOPA} (12~August) and Gebendorfer~\cite{Gebendorfer} (7~September); the aggregation function $\Phi_m$ of Lemma~\ref{lem:Phi}, with a
proof, appeared the same day in tawanerguo-cn~\cite{TawanRepo} and \mbox{sxuff}~\cite{sxuff}, and on 12~August in \mbox{trmdy} [V$\sim$]. Box
pruning by convexity certificates and tangent planes appears in the verifiers of \mbox{sxuff}, \mbox{AMTOPA} and \mbox{npip99}~\cite{npip99}, and interval-Hessian
basins in Zhaoyi-Tian's certificate; we have not checked whether our box bound differs from these in substance [M].
Six unrefereed values that are, as far as we can tell, bandwidth-one certificates [M] exceed Theorem~\ref{thm:S}, by $2.5\cdot10^{-5}$ to $1.18\cdot10^{-4}$:
\mbox{gfreund123}, in pull request~\#888 to the repository \mbox{GettysburgResearch/riemann}~\cite{Gettysburg888}, $0.67349239$, a draft pull request,
which states that independent analytic review is required and that ``no independently reviewed world-record claim is made''; it states its
bound for $0<\gamma\le T$ and builds on the window and one local inequality of Gebendorfer's draft [V$\sim$];
\mbox{RakitinD}~\cite{Rakitin}, $0.6734775921119$, described as a frozen computer-assisted result that builds on a comparator from
Gebendorfer's work, with, in its author's words, a ``partial Lean formalization with explicit remaining boundaries''; stated for
$0<\gamma\le T$ [V$\sim$, M];
Jonas Jakob Gebendorfer~\cite{Gebendorfer}, $0.6734760237$, called by its author an internally checked research draft; stated for
$0<\gamma\le T$ [V$\sim$];
\mbox{AMTOPA}~\cite{AMTOPA}, $0.673416490971$, described by its authors as an interval-certified research-draft record [V$\sim$];
Zhaoyi-Tian~\cite{ZhaoyiTian}, $0.6734$, announced as a Lean theorem with no hypotheses, checked by the kernel with the three standard axioms,
on the same Lean and Mathlib~\cite{Mathlib} versions as ours, with a kernel-replayed certificate of $3{,}203$ leaves, crediting \mbox{ainta} (the retained spectral defect and seven
consecutive zeros), \mbox{sxuff} and \mbox{AMTOPA}'s earlier repository (position-dependent and exact shifted pressure) and \mbox{AMTOPA} (the corrected window,
pair loads and pressure data) for ingredients it uses [V$\sim$]; we have not built it. By its description it is a formally verified bound
above Theorem~\ref{thm:S}. Zhaoyi-Tian announced $0.673$ as a hypothesis-free Lean theorem on 31~August and $0.6734$ on 3~September~2026;
\mbox{MichaelMobius}~\cite{Hurtado} announces a hypothesis-free Lean theorem for $0.6731175$ (latest commits 21--22~September~2026), and we have not
established which of these was complete first; an end-to-end Lean reduction of the bound $0.673195$ to a single local-certificate premise
was published by \mbox{npip99}~\cite{npip99} on 11--12~August~2026 [V$\sim$];
and Michael Devine~\cite{Devine22066689}, $0.673399$ (stated also for simple and for distinct zeros, as a consequence), whose aggregation by
a piecewise-linear subaction appears to differ from the block aggregation used here and by the others [V$\sim$, M]. Devine's later
record~\cite{Devine22287432} announces a computation-free improvement of the Montgomery--Taylor constant, with a Lean formalisation in which
the published pair-correlation and Riemann--von Mangoldt inputs are represented by explicit interfaces [V$\sim$].
Further claims lie between $0.67300$ and $0.67332$ (among them \mbox{ainta}'s $0.673008527927$ and \mbox{trmdy}'s
$0.6733127422$)\footnote{The others we know of [V$\sim$]: \mbox{yuhangshi888}~\cite{yuhangshi} ($0.6733169771$), \mbox{AMTOPA}'s earlier $0.673262375585$,
\mbox{sxuff}~\cite{sxuff} ($0.673205978423$), \mbox{npip99}~\cite{npip99} ($0.673195198901$), tawanerguo-cn~\cite{TawanRepo} ($0.6731929114731422$),
\mbox{MichaelMobius}~\cite{Hurtado} ($0.6731175265$, announced with a hypothesis-free Lean proof), \mbox{vstaln}~\cite{vstaln} ($0.67306445176$),
Zhongshan-Big-Jun~\cite{Zhongshan} ($0.6730664726$ of 15~August; its README also lists $0.6730536459$ of 14~August) and uwe-schwarz~\cite{uweschwarz} ($0.6730266625$).};
teal-sea/zeta-lab~\cite{tealsea} reports Lean proofs of small-$K$ bounds ($0.6728470197$), which, according to its README, the Palomar registry
rebuilt and replayed [V$\sim$].
What is ours for Theorem~\ref{thm:S}, as far as we have found, is the even polynomial window with exact rational constants as a window for
$\zeta$ (its $H$ is an exact rational; Alp\"oge and Furman used a quartic polynomial window for the zeros of $\xi'$~\cite[Remark~7.1]{AF}),
and the box bound LB3 in its exact form (\S\ref{ssec:bnb}), which we have not compared in substance with the pruning of other verifiers. The
aggregation function $\Phi_m$ of Lemma~\ref{lem:Phi}, which Theorem~\ref{thm:S} needs on its concave branch since its block energy
$A=1.1144$ exceeds $m/(m-1)$, is not ours: it was published with proofs by tawanerguo-cn~\cite{TawanRepo} and \mbox{sxuff}~\cite{sxuff} on
11~August~2026 and by \mbox{trmdy}~\cite{trmdy} on 12~August, each using it on its concave branch; \mbox{yuhangshi888}~\cite{yuhangshi} uses a supporting
plane of it, and Zhaoyi-Tian prove in Lean a cap-one variant [V$\sim$]. We have verified none of these proofs.

\paragraph{Claims by other methods.}
The ceiling $\approx0.682$ of~\cite[\S7.2]{AF} concerns bandwidth-one certificates that use only the first two trace moments
(Section~\ref{sec:remarks}). For $N^s$ the largest value that the bandwidth-one claims above would give is $0.67349239$; for $N^d$, if their chains give $\frac12(1+S)$ as
those of~\cite[Theorem~A(ii)]{AF} and of Theorem~\ref{thm:S} do (Remark~\ref{rem:Sdist}), it is $\frac12(1+0.67349239)=0.8367461$. Both lie
below Theorem~\ref{thm:main}.
Hongyi Yang and Shihua Yang claim much more, by methods that use trace moments beyond the first two [V$\sim$]. In Zenodo record~21975237~\cite{YY1} they claim $0.7962$ for
simple zeros on the line and $0.8981$ for distinct zeros, with ineffective constants, a result they declare ``at the certified-candidate
level, not as an established theorem''; according to the record, its identity and certificate layers are compiled and kernel-checked in
Lean~4, while its analytic layer is not yet formalised and awaits external peer review [V$\sim$]. In Zenodo record~22065921~\cite{YY2}, Shihua Yang and Hongyi Yang
claim that almost all zeros of $\zeta$ are simple and lie on the critical line, in the density sense, unconditionally and with ineffective
constants, with the explicit bounds $0.910460$ for simple zeros on the line and $0.955230$ for distinct zeros; the record describes the
analytic chain as ``a certified candidate pending external review'' and states that its certificates are kernel-checked in
Lean~4 [V$\sim$]. The record page of~\cite{YY2} shows its files under embargo,
with the stated reason that the version contains errors and needs revision [V$\sim$, read 28~September~2026]. We have not verified either
claim; if they are correct, they supersede Theorems~\ref{thm:main}--\ref{thm:S} and Corollary~\ref{cor:SC}.

\subsection{Formal verification}\label{ssec:formal}
The four results are formalised in Lean~4~\cite{Lean4}, in a library \texttt{ZetaS} built on the Lean artifact
\texttt{Zeta23} of~\cite{AF} (toolchain v4.33.0-rc2, Mathlib~\cite{Mathlib} revision 51e6992), against the zero-counting functions that~\cite{AF} define
directly from Mathlib's \texttt{riemannZeta} (Appendix~\ref{app:lean}). We say that a declaration is \emph{kernel-checked} if
\texttt{\#print axioms} reports only \texttt{propext}, \texttt{Classical.choice} and \texttt{Quot.sound}: no \texttt{sorry}, no added axiom,
no \texttt{native\_decide}. The formalisation is in the repository \repourl, at tag \repotag{} (Appendix~\ref{app:layout}). The
trust level of each result is as follows.
\begin{enumerate}[label=(\alph*)]
\item \emph{Kernel-checked with no hypotheses.} The two headline theorems for the $K=5$ pair ($0.675158622$, $0.837579311$;
Appendix~\ref{app:replay}), which are in the library \texttt{ZetaSReplay} of the repository, built separately from
\texttt{ZetaS}. Further, the inputs of Proposition~\ref{prop:AF} for the two windows used, at every bandwidth
$\lambda<1$ (from the artifact of~\cite{AF} and a window-general version of its pipeline); every lemma, proposition and theorem of
Sections~\ref{sec:setting}--\ref{sec:sc} in its abstract form (not the remarks, among them Remark~\ref{rem:Sdist}, and not the examples of
sharpness), with the qualifications stated after Lemmas~\ref{lem:residue}
and~\ref{lem:Phi}, where the formal version is weaker than the printed statement but covers every use: the linear algebra of
Section~\ref{sec:linalg}; the local-to-global steps of Sections~\ref{sec:first}--\ref{sec:simple} (aggregation, packing, transfer, tight
chains, the room lemma; the on-line spectral condition and the deduction of Theorem~\ref{thm:main} from its certificates only for the two
certificates used, Theorems~\ref{thm:OLL} and~\ref{thm:SigmaD} for arbitrary certificate data being a mechanical generalisation not yet
formalised; the deductions of Theorem~\ref{thm:S} and Corollary~\ref{cor:SC} from their certificate); the assembly from the
zeros of \texttt{riemannZeta} to the matrices of Section~\ref{sec:setting} at every bandwidth $\lambda<1$ (Poisson and Gram identities, positive
definiteness of the kernel, the truncation defect, the inertia of off-line pairs), together with the limit $\lambda\to1^-$; the closing
arithmetic of all constants; and the majorant certificate of Lemma~\ref{lem:room}, replayed by the kernel in $614$ cells (about $88$
seconds).
\item \emph{Kernel-checked with one displayed hypothesis, a numerical certificate.} In the library the six headline theorems each have one
hypothesis: \texttt{CertAM5} for the $K=5$ pair, \texttt{CertAM7} for the $K=7$ pair ($0.676102666$, $0.838051333$), and \texttt{CertS8} for
Theorem~\ref{thm:S} and Corollary~\ref{cor:SC}. Each of these is one \texttt{Prop} asserting that a local inequality holds for all gap
vectors (Appendix~\ref{app:lean}); the data that realise it are identified by sha256 in Table~\ref{tab:certs}, and they are verified
outside Lean by interval arithmetic (\S\ref{ssec:bnb}). Nothing else is assumed: in the library at commit \texttt{bb7223b}, \texttt{\#print axioms} for each
of the six headline theorems gives only the three standard axioms (Appendix~\ref{app:census}). The fourteen statements of record (the six headline
statements, their six cumulative forms on $0<\Im\rho\le T$ and two non-vacuity companions) are in the comparator package of the library
(Appendix~\ref{app:record}), and all fourteen are proved with the three standard axioms. For the $K=5$ pair the
hypothesis is discharged by the kernel replay, (a) and (c). The $K=7$ certificate and the certificate behind Theorem~\ref{thm:S} are too
large for kernel replay on the hardware used (\S\ref{ssec:replay}) and stay as displayed hypotheses.
\item \emph{The kernel replay of the $K=5$ certificate.} The kernel checker of the replay is version~3, which admits a tangent line only
when the tangent point lies in the span interval (without this guard a leaf check could accept a false claim, \S\ref{ssec:replay}). The
checker, its soundness nodes and the chain \texttt{ZetaS.CertV2.\allowbreak ChainV3.\allowbreak certAM5\_of\_checks}, which derives
\texttt{CertAM5} from the acceptance of the $20$ class certificates by the checker and the identity of their data with the data of
\texttt{CertAM5}, are in the library with the three standard axioms (Appendix~\ref{app:chain}). The full replay, against the library at
commit \texttt{bb7223b}, whose Lean files the release changes only in comments (Appendix~\ref{app:layout}), compiled all $1{,}478$ files
with no failure, and the kernel checked the $46{,}047$ leaves and $89{,}362$ point records in exact rational
arithmetic (\S\ref{ssec:replay}). Its top files prove \texttt{ZetaS.CertV2.\allowbreak certAM5\_replayed}${}:{}$\texttt{ZetaS.CertAM5} and,
from it, the two theorems of the $K=5$ pair without hypothesis (Appendix~\ref{app:replay}). Each of
the three depends only on \texttt{propext}, \texttt{Classical.choice} and \texttt{Quot.sound} and rests on no open declaration
(Appendix~\ref{app:census}). The
$1{,}478$ source files match their manifest file by file, and a scan of all of them finds no \texttt{sorry}, \texttt{native\_decide},
\texttt{axiom}, \texttt{unsafe}, \texttt{implemented\_by} or \texttt{extern}. The replay files, byte-identical to those
checked, and the two theorems are in the repository as the library \texttt{ZetaSReplay}, outside the default build targets;
\texttt{lake build ZetaSReplay} repeats the replay (\S\ref{ssec:reproduce}). The certificates themselves were
found by floating-point and interval computation; the kernel checked them.
\end{enumerate}

\paragraph{Earlier formal work.} Formal proofs of bounds of this kind were announced before ours [V$\sim$]: Zhaoyi-Tian's
hypothesis-free Lean theorems for $N_0^s$, at $0.673$ (31~August) and $0.6734$ (3~September~2026)~\cite{ZhaoyiTian}; a hypothesis-free Lean
theorem for $N_0^s$ at $0.6731175$ announced by \mbox{MichaelMobius}~\cite{Hurtado} (latest commits 21--22~September~2026; we could not
establish when its Lean theorem was first complete); \mbox{npip99}'s end-to-end Lean reduction of
$0.673195$ to one local-certificate premise~\cite{npip99} (11--12~August); Lean-proved small-$K$ bounds for $N_0^s$ in
teal-sea/zeta-lab~\cite{tealsea}, which, according to its README, the Palomar registry rebuilt; the finite deduction of \mbox{yuhangshi888}~\cite{yuhangshi} in Lean; Devine's
formalised computation-free improvement of the Montgomery--Taylor constant~\cite{Devine22287432}, with the pair-correlation and
Riemann--von Mangoldt inputs as interfaces; partial formalisations by \mbox{RakitinD}~\cite{Rakitin} and of the certificate layers of Yang and
Yang~\cite{YY1,YY2}; and a formalisation of Lamzouri's proof, with the Riemann--von Mangoldt formula and a pair-correlation asymptotic as
stated hypotheses, by \mbox{AxiomMath}~\cite{AxiomMath}. Our formal theorem for simple zeros on the line, Theorem~\ref{thm:S}, has a smaller constant
than the one Zhaoyi-Tian announces, and its certificate is a hypothesis. Its formal status is that of \mbox{npip99}'s earlier theorem. The Lean artifact of~\cite{AF}, on which our
library is built, already proves hypothesis-free bounds for $N^d$ ($\frac34$ and $\frac56$, and $0.75329\ldots$ and
$0.83625\ldots$ with the Montgomery--Taylor window) as well as for $N_0^s$ and for the distinct zeros on the line; apart from these and, under
its hypotheses, \mbox{AxiomMath}'s formalisation of Lamzouri's bounds, we know of no earlier formal result for $N^s$, $N^d$ or $N^{\mathrm{sc}}$
(Yang and Yang state that the identity and certificate layers of their claims, which include $N^d$, are kernel-checked in Lean~4).

\subsection{Plan}
Section~\ref{sec:setting} fixes the setting of~\cite{AF} and the inputs taken from it, at bandwidth $\lambda<1$. Section~\ref{sec:linalg}
collects the linear algebra. Section~\ref{sec:first} proves Theorem~\ref{thm:S}, Section~\ref{sec:simple} Theorem~\ref{thm:main}, and
Section~\ref{sec:sc} Corollary~\ref{cor:SC}. Section~\ref{sec:cert} describes the certificates and Section~\ref{sec:remarks} the limits of the
method. Appendix~\ref{app:lean} reproduces the Lean statements; Appendix~\ref{app:how} records how the work was done.

\section{Setting and the inputs from Alp\"oge--Furman}\label{sec:setting}

We follow the notation of~\cite[\S\S1--2]{AF}. Throughout, $\ell:=\log(T/2\pi)$ and $N:=N(T,2T)=\frac T{2\pi}(\ell+2\log2-1)+O(\ell)$.
For a zero $\rho$ put $\gamma_\rho:=(\rho-\frac12)/i\in\C$, which is real if and only if $\rho$ lies on the critical line. Off-line zeros come in
pairs $\pi=\{\rho,1-\bar\rho\}$ with $\gamma_{1-\bar\rho}=\overline{\gamma_\rho}$ and equal multiplicities.

\subsection{The window, the test family and the matrix}\label{ssec:family}
Fix a \emph{bandwidth} $\lambda\in(0,1]$ and put $L:=\lambda\ell$. A \emph{window} is an even function $\psi\in C^2([-\frac12,\frac12])$ with
$0<\psi\le1$. (The hypotheses of~\cite[\S2.2]{AF} are that $\psi$ is even, $C^2$ and positive; the bound $\psi\le1$ is used there tacitly,
through $0\le\varphi\le1$.) Fix a smooth nondecreasing function $\chi$ with $\chi=0$ on $(-\infty,0]$ and $\chi=1$ on $[1,\infty)$, and let $\chi_w(t):=\chi(t/w)$, the
ramp of width $w$, for $1\le w=o(L)$; \cite{AF} take a $C^\infty$ ramp of width $1$. Implied constants may depend on $\chi$ and $\psi$. Put, as
in~\cite[(2.7)]{AF},
\begin{equation}\label{eq:window}
\varphi(u):=\chi_w\big(\tfrac L2+u\big)\,\chi_w\big(\tfrac L2-u\big)\,\psi(u/L)^{1/2},
\end{equation}
so that $\varphi$ is real and even, $\operatorname{supp}\varphi=[-\frac L2,\frac L2]$ and $\varphi(\pm\frac L2)=0$. Let
$a:=L^{-1}\int\varphi^2=\int\psi+O(w/L)$ and $\widehat\varphi(z):=\int\varphi(u)e^{-iuz}\,du$, an even entire function, real on $\R$, with
$\widehat\varphi(\bar z)=\overline{\widehat\varphi(z)}$. The grid is $\alpha_k:=T+2\pi k/L$ for $0\le k<d:=\floor{LT/2\pi}$, and for a zero
$\rho$ we put
\[
v_\rho:=\big(\widehat\varphi(\gamma_\rho-\alpha_k)\big)_{0\le k<d}\in\C^d,\qquad u_\rho:=(aL^2)^{-1/2}v_\rho .
\]
Let $I:=[T,2T)$ and $I':=[T-\sqrt T,\,2T+\sqrt T)$. The matrix of~\cite[(2.10)]{AF} is
\begin{equation}\label{eq:G}
G:=\sum_{\Re\gamma_\rho\in I'}m_\rho\,u_\rho u_\rho^{\tsp},
\end{equation}
the sum running over distinct zeros, with the bilinear transpose. It is real symmetric: $v_{1-\bar\rho}=\overline{v_\rho}$ because the grid is
real and $\widehat\varphi(\bar z)=\overline{\widehat\varphi(z)}$, the multiplicities of $\rho$ and $1-\bar\rho$ agree, and both have the same
$\Re\gamma$, so that they enter~\eqref{eq:G} together. An on-line zero contributes $m_\rho u_\rho u_\rho^{\tsp}\succeq0$. An off-line pair with
$u_\rho=x+iy$ ($x,y\in\R^d$) contributes $m_\pi(u_\rho u_\rho^{\tsp}+\bar u_\rho\bar u_\rho^{\tsp})=2m_\pi(xx^{\tsp}-yy^{\tsp})$, which has at
most one positive and at most one negative eigenvalue.

\subsection{The kernel and the window constant}\label{ssec:kernel}
For a window $\psi$ put $v_\psi:=\psi/\!\int\psi$ and
\[
k_\psi(x):=\int_{-1/2}^{1/2}v_\psi(s)\cos(2\pi xs)\,ds ,
\]
an even, positive definite function with $k_\psi(0)=1$ and $\abs{k_\psi}\le1$. Let
\[
J(\psi):=\iint\abs{u-v}\,\psi(u)\psi(v)\,du\,dv,\qquad R_\lambda(\psi):=\frac{\int\psi^2+\lambda^2J(\psi)}{\lambda\,(\int\psi)^2},\qquad
H_\lambda(\psi):=2-R_\lambda(\psi),
\]
all integrals over $[-\frac12,\frac12]$, and $R:=R_1$, $H:=H_1$. Substituting $v=u+r$ gives the one-integral form
$J(\psi)=2\int_0^1r\,(\psi\star\psi)(r)\,dr$ with $(\psi\star\psi)(r):=\int\psi(u)\psi(u+r)\,du$; the prime side of~\cite{AF} produces this form,
and the two agree for every continuous $\psi$ (evenness is not needed). $R$ is the constant of~\cite[Lemma~5.6]{AF}; for $\psi=\cos(\sqrt2s)$,
$H=\frac32-\cot(1/\sqrt2)/\sqrt2=0.6725007036\ldots$, which is optimal among constants of the form $2-R(\psi)$~\cite{CCLM}. Every quantity
used below is homogeneous of degree $0$ in $\psi$, so a window may be rescaled by a positive constant without changing $k_\psi$, $R_\lambda$
or $H_\lambda$. For $\psi(s)=\cos(\alpha s)$,
\begin{equation}\label{eq:kcos}
k_\psi(x)=\frac{\sinc(\pi x-\frac\alpha2)+\sinc(\pi x+\frac\alpha2)}{2\sinc(\frac\alpha2)},\qquad\sinc y=\frac{\sin y}y\ (y\ne0),\quad\sinc0=1 .
\end{equation}

\subsection{The inputs}\label{ssec:inputs}
\begin{proposition}[Inputs from~\cite{AF}]\label{prop:AF}
Let $\psi$ be a window and $0<\lambda<1$. As $T\to\infty$:
\begin{enumerate}[label=\textup{(AF\arabic*)}]
\item \textup{(\cite[Lemma~2.1]{AF})} $\norm{u_\rho}^2\le1$ for every on-line zero.
\item \textup{(\cite[Lemma~3.1 and Proposition~4.1]{AF})} Each off-line pair contributes to $G$ a real symmetric matrix with at most one
positive and at most one negative eigenvalue.
\item \textup{(\cite[Proposition~4.2]{AF})} $\tr G=N(I')+O(T^{1/2}L^2)$, and $N(I')=N+O(\sqrt T\log T)$.
\item \textup{(\cite[Theorem~5.7]{AF} at $\lambda=1$; for $\lambda<1$ its proof, as indicated in~\cite[Remark~6.1]{AF}, which states the case
$\psi=\psi_0$; the general formula for $R_\lambda(\psi)$ is formally verified for the two windows used here)} $\HS{G}^2=(R_\lambda(\psi)+o(1))\,N$.
\end{enumerate}
\end{proposition}

The prime-side input behind (AF3)--(AF4) is the Montgomery--Vaughan mean value theorem for a Dirichlet polynomial of length
$(T/2\pi)^\lambda$, as in~\cite{AF}.
\cite{AF} prove (AF4) also at $\lambda=1$, with error $O(L^{-1})$; all arguments below go through verbatim there. The formal proof works at
$\lambda<1$, where its prime-side error term is $o(N)$, and then lets $\lambda\to1^-$ (\S\ref{ssec:bandwidth}); we follow it.

\subsection{Normalised ordinates and the Gram identity}\label{ssec:gram}
For an on-line zero put $x_\rho:=\gamma_\rho L/2\pi$ (the \emph{normalised ordinate}) and, for $k\in\Z$,
$U_\rho(k):=(aL^2)^{-1/2}\widehat\varphi(\gamma_\rho-\alpha_k)$, so that $u_\rho=(U_\rho(k))_{0\le k<d}$. Let
\[
\delta_\rho:=\sum_{k\notin[0,d)}U_\rho(k)^2,\qquad \widehat E_{\rho\rho'}:=\sum_{k\notin[0,d)}U_\rho(k)U_{\rho'}(k),\qquad
k_\varphi(\xi):=\frac{\int\varphi(Ls)^2e^{2\pi i\xi s}\,ds}{\int\varphi(Ls)^2\,ds}.
\]
Let $v_\varphi:=\varphi(L\cdot)^2/\!\int\varphi(L\cdot)^2$, $\delta_0:=\int v_\psi\,(1-\varphi(L\cdot)^2/\psi)$, and
\[
\eps_T:=\norm{v_\varphi-v_\psi}_{L^1},\qquad \eps_T':=\frac{\delta_0}{1-\delta_0}.
\]
Both tend to $0$ because $w=o(L)$. Finally let $\Lambda:=\abs{I'}L/2\pi$ be the normalised length of $I'$.

\begin{lemma}[Gram identity and residues]\label{lem:residue}
Let $T$ be so large that $a\ge\frac12\int\psi$, so that $\int\varphi^2>0$. For on-line zeros $\rho,\rho'$ with $\gamma_\rho,\gamma_{\rho'}\in I'$:
\begin{enumerate}[label=\textup{(\roman*)}]
\item $\sum_{k\in\Z}U_\rho(k)U_{\rho'}(k)=k_\varphi(x_\rho-x_{\rho'})$; hence $\norm{u_\rho}^2=1-\delta_\rho$ with $0\le\delta_\rho\le1$, and
$\langle u_\rho,u_{\rho'}\rangle=k_\varphi(x_\rho-x_{\rho'})-\widehat E_{\rho\rho'}$, where $\widehat E\succeq0$ and
$\abs{\widehat E_{\rho\rho'}}\le(\delta_\rho\delta_{\rho'})^{1/2}$.
\item $k_\varphi$ is even and positive definite with $k_\varphi(0)=1$, so $\abs{k_\varphi}\le1$; moreover $\sup_\R\abs{k_\varphi-k_\psi}\le\eps_T$
and $v_\varphi\le(1+\eps'_T)\,v_\psi$, with $\eps_T'\ge0$.
\item $\sum_{\rho}m_\rho\delta_\rho=o(N)$ \textup{(}indeed $\ll\sqrt T\log T$\textup{)}, the sum over on-line zeros with $\gamma_\rho\in I'$.
\item The consecutive gaps between the normalised ordinates of any set of zeros with $\Re\gamma\in I'$ add up to at most
$\Lambda=\lambda N(1+o(1))$; every multiplicity is $\ll\log T$.
\end{enumerate}
\end{lemma}

\begin{proof}
(i) With $f(u):=\varphi(u)e^{-iu(\gamma-T)}$, supported in $[-\frac L2,\frac L2]$ and vanishing at the ends, $\widehat\varphi(\gamma-\alpha_k)$ is
$L$ times a Fourier coefficient of $f$ on the period $L$; Parseval gives the identity~\cite[Lemma~2.1]{AF}, and the case $\rho=\rho'$ gives
$\sum_{k\in\Z}U_\rho(k)^2=1$. Subtracting the off-grid part gives the rest; $\widehat E$ is a Gram matrix, and Cauchy--Schwarz gives the bound on its
entries. (The normalisation needs $\int\varphi^2>0$: for the zero window the identity fails, $k_\varphi(0)$ being $0$.)
(ii) Every kernel matrix $(k_\varphi(x_i-x_j))$ is the Gram matrix of full-lattice vectors by~(i), so $k_\varphi$ is positive definite, and it is
even because $\varphi$ is. Next $v_\varphi=(1-\delta_0)^{-1}(\varphi(L\cdot)^2/\psi)\,v_\psi\le(1+\eps_T')v_\psi$, since the ramp only removes mass;
and $\abs{k_\varphi-k_\psi}\le\norm{v_\varphi-v_\psi}_1$.
(iii) For real $z$, two integrations by parts give $\abs{\widehat\varphi(z)}\le\abs z^{-2}\int\abs{\varphi''}$, and
$\abs{\widehat\varphi(z)}\le\int\abs\varphi\le L\sqrt a$; with $\xi=(\gamma_\rho-\alpha_k)L/2\pi$ this is
$\abs{U_\rho(k)}\le\min\{1,\lambda_\varphi\xi^{-2}\}$ with $\lambda_\varphi:=L\int\abs{\varphi''}/(4\pi^2\sqrt a)\ll L/w+1$ (here
$a\ge\frac12\int\psi$ and $\int\abs{\chi_w''}=w^{-1}\int\abs{\chi''}$ are used). The grid covers $[T,2T)$, and $\delta_\rho\le1$ always. The zeros with $\gamma_\rho\in I'\setminus[T,2T)$ number
$\ll\sqrt T\log T$, and those in $[T,2T)$ within distance $1$ (in $\gamma$) of $T$ or $2T$ number $\ll\log T$. The others lie in shells of unit
length in $\gamma$ at distance $j\ge1$ from the ends of $[T,2T)$, each with $\ll\log T\ll L$ zeros; for such a zero every off-grid $k$ has
$\abs\xi\ge jL/2\pi$, so $\delta_\rho\ll\lambda_\varphi^2(jL)^{-3}$, and these zeros contribute $\ll\sum_jL\lambda_\varphi^2(jL)^{-3}\ll1$.
(iv) is the Riemann--von Mangoldt formula and $N(t+1)-N(t)\ll\log t$.
\end{proof}

The formal proof has (ii) in a weaker form, the positive definiteness of $(1+\eps_T')k_\psi-k_\varphi$ in place of the pointwise bound
$v_\varphi\le(1+\eps_T')v_\psi$, and (iii) with $o(N)$ only; these are all that is used below.

At bandwidth $\lambda$ the normalised ordinates have mean spacing $\lambda$, not $1$. Every local inequality below is stated in these units,
where the kernel between two on-line zeros is $k_\varphi(x_\rho-x_{\rho'})$ at every $\lambda$; the certificates therefore apply at every
$\lambda$ without change, and $\Lambda\le N(1+o(1))$ is the only place where the density enters.

\subsection{The limit \texorpdfstring{$\lambda\to1^-$}{lambda to 1}}\label{ssec:bandwidth}
Each of our bounds is proved in the form: for every $\lambda\in(0,1)$ and every $\eps>0$, $X(T,2T)\ge(f(\lambda)-\eps)N$ for $T\ge T_0$, where
$f(\lambda)$ is obtained from the bandwidth-one constant by replacing $H(\psi)$ by $H_\lambda(\psi)$. Since $R_\lambda(\psi)\to R(\psi)$ as
$\lambda\to1^-$ and each constant is continuous (indeed affine) in $H$, letting $\lambda\to1^-$ gives the bandwidth-one constants. For instance the
constant $\Sigma$ of Theorem~\ref{thm:main} at $K=7$ equals $0.5950$, $0.66902$ and $0.675404$ at $\lambda=0.9$, $0.99$ and $0.999$.
(Equivalently, in units of the mean spacing the certificates transfer by the exact identity $F^{(\lambda)}_{\gamma,\lambda\mu}(h)=F_{\gamma,\mu}(\lambda h)$,
where $F^{(\lambda)}$ is built from $k_\psi(\lambda\,\cdot)$; so no slack is consumed. The formal proof does not need this, because it works in
normalised units throughout.)

\section{Linear algebra}\label{sec:linalg}

All matrices are finite. For a Hermitian matrix $X$, $\lambda_1(X)\ge\lambda_2(X)\ge\cdots$ are its eigenvalues, $n_\pm(X)$ the numbers of positive
and negative eigenvalues, $\HS X^2=\tr X^2$, and $\tr f(X)=\sum_if(\lambda_i(X))$ for a function $f$ on the spectrum.

\subsection{The stability rank--trace inequality}
Let
\begin{equation}\label{eq:Psi}
\Psi(t):=\begin{cases}(t-1)^2,&0\le t\le2,\\ 2t-3,&t\ge2,\end{cases}
\end{equation}
so that $\Psi\ge0$, $\Psi(0)=1$, $\Psi$ is $2$-Lipschitz on $[0,\infty)$, and $\Psi(t)=(t-1)^2-((t-2)_+)^2$. It is convex, being the supremum of
the tangent lines to $(t-1)^2$ at the points $s\le2$.

\begin{lemma}[Stability rank--trace inequality]\label{lem:stab}
Let $P,Q$ be Hermitian $d\times d$ matrices with $P=VV^*$, $V\in\C^{d\times r}$, and $n_+(Q)\le b$. Put $M:=V^*V\in\C^{r\times r}$. Then
\begin{equation}\label{eq:stab}
r\ \ge\ 2\tr P+4\tr Q-4b-\HS{P+Q}^2+\tr\Psi(M).
\end{equation}
\end{lemma}

Dropping $\tr\Psi(M)\ge0$ gives~\cite[Lemma~3.2]{AF}. Equality holds for $P=\Pi_1$, $Q=2\Pi_2$ with orthogonal projections $\Pi_1\perp\Pi_2$.
This form of the inequality, with the term $\tr\Psi(M)$, was published by \mbox{ainta}~\cite{ainta} on 11~August~2026; we know of no earlier
source [V$\sim$].

\begin{proof}
The nonzero eigenvalues of $VV^*$ and $V^*V$ agree with multiplicity, so $\sum_{i\le r}g(\lambda_i(M))=\sum_{j\le d}g(\lambda_j(P))$ for every $g$
with $g(0)=0$. Applied to $g=\Psi-1$ this gives $\tr\Psi(M)-r=\sum_{j\le d}(\Psi(p_j)-1)$, where $p_1\ge\dots\ge p_d\ge0$ are the eigenvalues of
$P$; no case distinction between $r\le d$ and $r>d$ is needed. Write $Q=Q_+-Q_-$ with $Q_\pm\succeq0$, $Q_+Q_-=0$; let $q_1,\dots,q_{b'}$
($b'\le b$) be the positive eigenvalues of $Q_+$ and $n_1\ge n_2\ge\dots\ge0$ those of $Q_-$. Since $\tr(PQ_+)\ge0$ and $\tr(Q_+Q_-)=0$,
\[
\HS{P+Q}^2\ \ge\ \HS{P-Q_-}^2+\sum_jq_j^2\ \ge\ \sum_i(p_i-n_i)^2+\sum_jq_j^2,
\]
the second step by von Neumann's trace inequality $\tr(PQ_-)\le\sum_ip_in_i$. Using $\tr P=\sum p_i$, $\tr Q=\sum q_j-\sum n_i$ and the identity
$1-2p+4n+(p-n)^2=(p-n-1)^2+2n$,
\[
d-\big[2\tr P+4\tr Q-4b-\HS{P+Q}^2\big]\ \ge\ \sum_{i\le d}\big[(p_i-n_i-1)^2+2n_i\big]+\sum_{j\le b'}(q_j-2)^2+4(b-b').
\]
For $p\ge0$ the minimum over $n\ge0$ of $(p-n-1)^2+2n$ is attained at $n=(p-2)_+$ and equals $\Psi(p)$. Hence the right side is at least
$\sum_{i\le d}\Psi(p_i)$, and subtracting $d-r=\sum_{i\le d}1-r$ gives~\eqref{eq:stab} by the first sentence.
\end{proof}

\begin{lemma}[Two-parameter form]\label{lem:Ast}
In the setting of Lemma~\ref{lem:stab}, let $0\le s\le t$ and
\[
\Psi_{s,t}(p):=\begin{cases}(p-s)^2,&0\le p\le t,\\ (t-s)^2+2(t-s)(p-t),&p\ge t.\end{cases}
\]
Then $s^2r+t^2b\ge2s\tr P+2t\tr Q-\HS{P+Q}^2+\tr\Psi_{s,t}(M)$. For $(s,t)=(1,2)$ this is Lemma~\ref{lem:stab}.
\end{lemma}

\begin{proof}
As above, by the identity $s^2-2sp+(p-n)^2+2tn=(p-n-s)^2+2(t-s)n$,
\[
s^2d+t^2b-\big[2s\tr P+2t\tr Q-\HS{P+Q}^2\big]\ \ge\ \sum_{i\le d}\big[(p_i-n_i-s)^2+2(t-s)n_i\big]+\sum_{j\le b'}(q_j-t)^2+t^2(b-b').
\]
For $p\ge0$ the minimum over $n\ge0$ of $(p-n-s)^2+2(t-s)n$ is attained at $n=(p-t)_+$ and equals $\Psi_{s,t}(p)$, so the right side is at
least $\sum_{i\le d}\Psi_{s,t}(p_i)$. Since $\Psi_{s,t}(0)=s^2$ (as $t\ge0$), the first sentence of the proof of Lemma~\ref{lem:stab},
applied to $g=\Psi_{s,t}-s^2$, gives $\tr\Psi_{s,t}(M)-s^2r=\sum_{i\le d}\Psi_{s,t}(p_i)-s^2d$; subtracting $s^2(d-r)$ from both sides of
the display gives the claim.
\end{proof}

$\Psi_{s,t}$ is convex (the supremum of the tangents to $(p-s)^2$ at the points $\le t$), $C^1$ and nonnegative, and $\Psi_{1,t}\ge\Psi$ for
$t\ge2$ and every real $p$. The one-parameter case $(s,t)=(c/2,c)$, which contains Lemma~\ref{lem:stab} ($c=2$), is proved in Lean in
Zhaoyi-Tian's development~\cite{ZhaoyiTian}, in the form $c\tr P+\sum_i\mathrm{gc}_c(\lambda_i)+2c\tr Q-c^2b\le\HS{P+Q}^2$ [V$\sim$; that this is the
case $(c/2,c)$ of Lemma~\ref{lem:Ast} is our reading, M]. Corollary~\ref{cor:SC} needs the point $(1,2+\sqrt2)$, which lies off that ray.

\subsection{Energy and the function \texorpdfstring{$\Phi_m$}{Phi\_m}}
\begin{lemma}[Aggregation function]\label{lem:Phi}
Let $G\succeq0$ be $m\times m$ with unit diagonal, $m\ge2$, and $E:=\HS{G-I}^2$. Then $\tr\Psi(G)\ge\Phi_m(E)$, where
\[
\Phi_m(E):=\begin{cases}E,& E\le\frac{m}{m-1},\\[2pt]
E-\Big(\sqrt{\tfrac{(m-1)E}{m}}-1\Big)^2,& E\ge\frac m{m-1}.\end{cases}
\]
$\Phi_m$ is increasing and concave on $[0,\infty)$, $\Phi_m(0)=0$ and $\Phi_m(E)\le E$; equality holds for $G=(1-t)I+tJ$, $0\le t\le1$.
\end{lemma}

\begin{proof}
Let $1+u_i$ be the eigenvalues of $G$; then $\sum u_i=0$, $\sum u_i^2=E$ and $u_i\ge-1$, and $\tr\Psi(G)=E-\sum_{i\in S}(u_i-1)^2$ with
$S:=\{i:u_i>1\}$. If $S=\emptyset$ we are done. Otherwise let $j:=\abs S$ (so $1\le j<m$), $\sigma:=\sum_Su_i>j$, $F:=\sum_S(u_i-1)^2$ and
$V:=1+\sqrt F$. The numbers $u_i-1$, $i\in S$, are positive with sum $\sigma-j$, so $\sqrt F\le\sigma-j$; also $\sum_Su_i^2=F+2(\sigma-j)+j$, and
the $m-j$ remaining $u_i$ have sum $-\sigma$, hence sum of squares $\ge\sigma^2/(m-j)$. So
\[
E\ \ge\ (V-1)^2+2(V-1)+j+\frac{(V+j-1)^2}{m-j}=V^2+(j-1)+\frac{(V+j-1)^2}{m-j}\ \ge\ V^2\frac m{m-1},
\]
with equality in the last step for $j=1$, and for $j\ge2$ because $V+j-1\ge V$ and $m-j<m-1$. Thus $V\le U:=\sqrt{E(m-1)/m}$. If
$E\le m/(m-1)$ then $U\le1\le V$ forces $F=0$; otherwise $F=(V-1)^2\le(U-1)^2$. The two branches have slope $1$ at $E=m/(m-1)$ and the second is
concave.
\end{proof}

The inequality $\tr\Psi(G)\ge\Phi_m(E)$ is kernel-checked (the formal version does not need $m\ge2$); the properties of $\Phi_m$ listed after it
in the lemma and the equality case are not formalised, and the formal proofs do not use them. The function $\Phi_m$ and this inequality
were published earlier by tawanerguo-cn~\cite{TawanRepo} and \mbox{sxuff}~\cite{sxuff} (11~August~2026) and by \mbox{trmdy}~\cite{trmdy}
(12~August) [V$\sim$].

We use Lemma~\ref{lem:Phi} on its concave branch in Theorem~\ref{thm:S}. A consequence used there: for $0\le E$ and $0<A$,
$\Phi_m(E)\ge\frac{\Phi_m(A)}A\min(E,A)$, by concavity and $\Phi_m(0)=0$.

\subsection{Pinching and monotonicity}
\begin{lemma}[Pinching]\label{lem:pinch}
Let $G\succeq0$ be real symmetric and let $f$ be convex on $[0,\infty)$. For every partition of the index set into blocks $b$,
$\tr f(G)\ge\sum_b\tr f(G_{bb})$.
\end{lemma}

\begin{proof}
For each block let $w_{b,1},w_{b,2},\dots$ be an orthonormal eigenbasis of $G_{bb}$, extended by zero outside $b$. These vectors form an
orthonormal basis, in particular a Parseval frame of unit vectors, and $\langle w_{b,j},Gw_{b,j}\rangle=\lambda_j(G_{bb})$. For a unit vector $w$,
$\langle w,Gw\rangle=\sum_kc_k^2\lambda_k(G)$ with $\sum_kc_k^2=1$, where $c_k$ are the coordinates of $w$ in an eigenbasis of $G$; by Jensen,
$f(\langle w,Gw\rangle)\le\sum_kc_k^2f(\lambda_k(G))$. Summing over the frame, and using $\sum_{b,j}c_{k}(w_{b,j})^2=1$ for each $k$,
$\sum_b\tr f(G_{bb})\le\tr f(G)$.
\end{proof}

No convexity of $X\mapsto\tr f(X)$ is needed. The lemma applies to $\Psi$, to $\Psi_{s,t}$ and to $(2-t)_+^2$.

\begin{lemma}[Monotonicity]\label{lem:weyl}
Let $A,B$ be positive semidefinite of the same size with $A-B\succeq0$, and let $f$ be $c$-Lipschitz on $[0,\infty)$. Then
$\abs{\tr f(A)-\tr f(B)}\le c\tr(A-B)$.
\end{lemma}

\begin{proof}
By Weyl's inequality $\lambda_i(A)\ge\lambda_i(B)$ for every $i$, so
$\abs{\tr f(A)-\tr f(B)}\le c\sum_i(\lambda_i(A)-\lambda_i(B))=c\tr(A-B)$.
\end{proof}

This is the form that is formalised and used (in Lemma~\ref{lem:transfer}); the same proof gives it for Hermitian $A,B$ and $f$ Lipschitz on
an interval containing both spectra.

\subsection{Pair removal}
\begin{lemma}[Pair removal]\label{lem:pairrem}
Let $G\succeq0$ and $Q$ be real symmetric of the same size, with $n_+(Q)\le b$ and $n_-(Q)\le p$. Then
\[
\HS{G+Q}^2-\HS G^2-2\tr Q\ \ge\ 2(\tr Q-2b)-\sum_{i\le p}\big(\lambda_i(G)-2\big)_+^2 .
\]
\end{lemma}

\begin{proof}
Write $Q=Q_+-Q_-$ with $Q_\pm\succeq0$, $Q_+Q_-=0$, $\tau_+:=\tr Q_+$, $\tau_-:=\tr Q_-$. Since $\tr(GQ_+)\ge0$ and $\rank Q_+\le b$,
$\HS{G+Q}^2\ge\HS G^2-2\tr(GQ_-)+\HS{Q_-}^2+\tau_+^2/b$. For every real $\tau_+$,
$\tau_+^2/b-2\tau_+=2(\tau_+-2b)+(\tau_+-2b)^2/b\ge2(\tau_+-2b)$ (if $b=0$ then $Q_+=0$ and this step is void), and $\tau_+=\tr Q+\tau_-$. So the left
side is at least $2(\tr Q-2b)+4\tau_-+\HS{Q_-}^2-2\tr(GQ_-)$. By von Neumann, $\tr(GQ_-)\le\sum_{i\le p}\lambda_i(G)n_i$, where $n_1\ge\dots\ge n_p\ge0$
are the eigenvalues of $Q_-$, and $n^2+4n-2\lambda n\ge-(\lambda-2)_+^2$ for all $n\ge0$.
\end{proof}

The inequality is sharp: $G=\operatorname{diag}(3.5,0,0)$, $Q=\operatorname{diag}(-1.5,2,0)$.

\subsection{Localisation of the large eigenvalues}
Fix $c_*\ge0$ and put $\lambda_c:=2+\sqrt{c_*}$,
\[
\kappa(\lambda):=\big[(\lambda-2)_+^2-c_*\big]_+,\qquad \ch(X):=\sum_i\kappa(\lambda_i(X))
\]
for real symmetric $X$. The \emph{charge} $\ch(X)$ sees only the eigenvalues above $\lambda_c$; $\kappa$ is nondecreasing.

\begin{theorem}[Localisation]\label{thm:loc}
Let $M=\begin{pmatrix}A&Y\\Y^{\tsp}&R\end{pmatrix}$ be real symmetric, $A\in\R^{n\times n}$, $R\in\R^{r\times r}$, with $\lambda_{\max}(R)\le\lambda_c$
\textup{(}no condition if $r=0$\textup{)}. Then
\[
\ch(M)\ \le\ \HS{A-2I}^2+2\HS Y^2-\tr(2I-A)_+^2 .
\]
\end{theorem}

\begin{proof}
Let $\widetilde M:=\begin{pmatrix}A&Y\\Y^{\tsp}&\lambda_cI_r\end{pmatrix}$. Then $\widetilde M-M=0\oplus(\lambda_cI-R)\succeq0$; this is the only use
of the hypothesis. By Weyl, $\lambda_l(M)\le\lambda_l(\widetilde M)=:\tilde\mu_l$ for every $l$, and since $\kappa$ is nondecreasing,
$\ch(M)\le\ch(\widetilde M)$. By Cauchy interlacing in counting form: the subspace $0\oplus\R^r$, on which the form of $\widetilde M$ is
$\lambda_c\norm x^2$, shows that $\widetilde M$ has at least $r$ eigenvalues $\ge\lambda_c$, so $\tilde\mu_l\ge\lambda_c$ for $l\le r$; and the
number of eigenvalues of $\widetilde M$ above any $\theta$ is at most that of $A$ plus $r$, so $\tilde\mu_{r+j}\le\alpha_j:=\lambda_j(A)$ for
$1\le j\le n$. Let $h(\mu):=(\mu-2)^2-\kappa(\mu)\ge0$. Then $h(\tilde\mu_l)=c_*$ for $l\le r$, and $h(\tilde\mu_{r+j})\ge(2-\alpha_j)_+^2$
(if $\tilde\mu_{r+j}\le2$ then $h=(2-\tilde\mu_{r+j})^2$ and $\tilde\mu_{r+j}\le\alpha_j$; otherwise $\alpha_j>2$). Since
$\sum_l(\tilde\mu_l-2)^2=\HS{\widetilde M-2I}^2=\HS{A-2I}^2+2\HS Y^2+r\,c_*$,
\[
\ch(\widetilde M)=\sum_l(\tilde\mu_l-2)^2-\sum_lh(\tilde\mu_l)\ \le\ \HS{A-2I}^2+2\HS Y^2-\sum_{j\le n}(2-\alpha_j)_+^2 .\qedhere
\]
\end{proof}

No positivity is used, and no Schur complement or integral over thresholds: the domination by $\widetilde M$ above suffices. The
hypothesis $\lambda_{\max}(R)\le\lambda_c$ cannot be dropped.

\section{The local inequality and simple zeros on the line}\label{sec:first}

This section proves Theorem~\ref{thm:S}. The mechanism is not ours (\S\ref{ssec:others}); we give it in full because
Section~\ref{sec:simple} builds on it and Corollary~\ref{cor:SC} uses its proof.

\subsection{The local inequality}
\begin{definition}[Local inequality]\label{def:LI}
Let $K\ge2$. \emph{Position weights} are numbers $\gamma_{s,i}\ge0$, $1\le s\le K-1$, $1\le i\le K-s$, with $\sum_i\gamma_{s,i}=2$ for every
$s$; \emph{penalties} are numbers $\mu_i\ge0$, $1\le i\le K-1$, with $\nu:=\sum_i\mu_i$. For an even function $k$ with $k(0)=1$ and
$c\in\R$ we say that $(\mathrm{LI})$ \emph{holds for $k$ with data $(\gamma,\mu,c)$} if
\begin{equation}\label{eq:LI}
F(g):=\sum_{i=1}^{K-1}\mu_ig_i+\sum_{s=1}^{K-1}\sum_{i=1}^{K-s}\gamma_{s,i}\,k(g_i+\dots+g_{i+s-1})^2\ \ge\ c\qquad\text{for all }g\in[0,\infty)^{K-1}.
\end{equation}
\end{definition}

The variables $g_i$ are the gaps between $K$ consecutive simple zeros on the line, in normalised units, and $\mu\cdot g$ charges the length that
such a run occupies. With $\gamma_{s,i}=2/(K-s)$ and $\mu_i=\mu$ this is the uniform functional of~\cite{ainta}. Since $k^2\ge0$, $F(g)\ge\mu_ig_i$,
so~\eqref{eq:LI} holds automatically outside a bounded box when every $\mu_i>0$; this is what makes it a finite computation (\S\ref{ssec:bnb}).
Coincident points ($g_i=0$) belong to the domain.

\begin{lemma}[Local to block]\label{lem:local}
Suppose $(\mathrm{LI})$ holds for $k$ with data $(\gamma,\mu,c)$. Let $y_1<\dots<y_m$ be real, $m\ge K$, with gaps $g_i:=y_{i+1}-y_i$. Then
\[
E:=\sum_{i\ne j}k(y_i-y_j)^2\ \ge\ c\,(m-K+1)-\sum_{t=1}^{m-K+1}\sum_{i=1}^{K-1}\mu_i\,g_{t+i-1}.
\]
\end{lemma}

\begin{proof}
Apply~\eqref{eq:LI} to each run $y_t,\dots,y_{t+K-1}$, $1\le t\le m-K+1$, and sum. A pair at index distance $s<K$ receives total weight at most
$\sum_i\gamma_{s,i}=2$, which is its contribution to $E$ (two ordered pairs); pairs at index distance $\ge K$ are dropped, since $k^2\ge0$.
\end{proof}

\begin{proposition}[Packing]\label{prop:pack}
Let $k$ be even and positive definite with $k(0)=1$, and suppose $(\mathrm{LI})$ holds for $k$ with data $(\gamma,\mu,c)$. Let
$y_1<\dots<y_S$ lie in an interval of length $\Lambda\ge0$ and $G:=(k(y_i-y_j))_{i,j}$. For $m\ge K$ put
\begin{equation}\label{eq:ABtau}
A:=c\,(m-K+1),\qquad B:=\Phi_m(A),\qquad \tau:=\frac{\nu\,(m-K+1)}m .
\end{equation}
If $A>0$, then $\tr\Psi(G)\ge\frac BmS-2B-\frac BA\tau\Lambda$.
\end{proposition}

The hypothesis $\Lambda\ge0$ matters only when $S=0$; without it the statement is false ($K=2$, $k\equiv1$, $c=2$, $m=2$, $S=0$, $\Lambda=-10$).

\begin{proof}
$G\succeq0$ has unit diagonal. Fix an offset $o\in\{0,\dots,m-1\}$ and cut the ordered points into consecutive blocks of $m$, starting after
the first $o$ points; there are at least $S/m-2$ full blocks. By Lemma~\ref{lem:pinch} (discarding the partial blocks, for which
$\tr\Psi\ge0$) and Lemma~\ref{lem:Phi}, $\tr\Psi(G)\ge\sum_{\text{full }b}\Phi_m(E_b)$. By Lemma~\ref{lem:local}, $E_b\ge A-P_b$ with
$P_b\ge0$ the penalty sum of the runs inside $b$, and $\Phi_m(E_b)\ge\frac BA\min(E_b,A)\ge B-\frac BAP_b$. Hence
$\tr\Psi(G)\ge B(S/m-2)-\frac BA\sum_bP_b$ for every $o$. Now average over $o$. A given gap occupies each relative position
$j\in\{0,\dots,m-1\}$ in its block for exactly one offset ($j=0$ meaning that it straddles two blocks, where it is not charged); at position
$j\ge1$ it is charged $\sum_t\mu_{j-t+1}$ over the runs $t$ of the block that contain it, and summing over $j$ gives
$\sum_{t=1}^{m-K+1}\sum_{i=1}^{K-1}\mu_i=(m-K+1)\nu$. So the average of $\sum_bP_b$ is at most $\tau\sum_ig_i\le\tau\Lambda$.
\end{proof}

\subsection{From the zeros to the local inequality}
Let $\mathrm{on}_1$ be the set of simple on-line zeros with $\gamma\in I'$, $s_1:=\#\mathrm{on}_1$, and let
$M^\circ:=(\langle u_\rho,u_{\rho'}\rangle)_{\rho,\rho'\in\mathrm{on}_1}$ be their Gram matrix.

\begin{lemma}[Transfer]\label{lem:transfer}
Let $\psi$ be a window and suppose $(\mathrm{LI})$ holds for $k_\psi$ with data $(\gamma,\mu,c)$, $c>0$. Let $0<c'<c$ and let $A',B'$ be as
in~\eqref{eq:ABtau} with $c'$ in place of $c$. Then for $T\ge T_0(c')$,
\[
\tr\Psi(M^\circ)\ \ge\ \frac{B'}m\,s_1-\frac{B'}{A'}\,\tau\,\Lambda-o(N).
\]
\end{lemma}

\begin{proof}
By Lemma~\ref{lem:residue}(i), $M^\circ=K^\circ-\widehat E^\circ$ with $K^\circ:=(k_\varphi(x_\rho-x_{\rho'}))$ and $\widehat E^\circ\succeq0$. Both
$M^\circ$ and $K^\circ$ are positive semidefinite, $\Psi$ is $2$-Lipschitz on $[0,\infty)$, and $\tr\widehat E^\circ=\sum\delta_\rho=o(N)$, so
Lemma~\ref{lem:weyl} gives $\abs{\tr\Psi(M^\circ)-\tr\Psi(K^\circ)}=o(N)$. Next $\abs{k_\varphi^2-k_\psi^2}\le2\eps_T$ by
Lemma~\ref{lem:residue}(ii), so the functional~\eqref{eq:LI} for $k_\varphi$ differs from that for $k_\psi$ by at most
$2\eps_T\sum_{s,i}\gamma_{s,i}=4(K-1)\eps_T$, and $(\mathrm{LI})$ holds for $k_\varphi$ with $c'$ once $T$ is large. The normalised ordinates of
$\mathrm{on}_1$ are distinct and lie in an interval of length $\Lambda$; apply Proposition~\ref{prop:pack} to $K^\circ$.
\end{proof}

\begin{theorem}[Stability plus packing]\label{thm:G}
Let $\psi$ be a window and suppose $(\mathrm{LI})$ holds for $k_\psi$ with data $(\gamma,\mu,c)$, $c>0$. Let $m\ge K$ with $1-B/m>0$
\textup{(}automatic when $A<m$, since $B\le A$\textup{)}. Then
\[
\liminf_{T\to\infty}\frac{N_0^s(T,2T)}{N(T,2T)}\ \ge\ \frac{H(\psi)-(B/A)\,\tau}{1-B/m}.
\]
\end{theorem}

\begin{proof}
Fix $\lambda\in(0,1)$. Apply Lemma~\ref{lem:stab} with $V:=[u_\rho]_{\rho\in\mathrm{on}_1}$ (real, since $\gamma_\rho$ is real), $P:=VV^{\tsp}$,
$r=s_1$, $M=V^{\tsp}V=M^\circ$, $Q:=G-P$ and $b:=s_2+p$, where $s_2$ is the number of distinct multiple on-line zeros and $p$ the number of
off-line pairs with $\Re\gamma\in I'$: each multiple on-line zero contributes a rank-one positive semidefinite term and each pair at most one
positive eigenvalue (AF2), so $n_+(Q)\le b$ by subadditivity of $n_+$. Since $\tr P+2b\le s_1+2s_2+2p\le N(I')$ by (AF1),
\[
2\tr P+4\tr Q-4b=4\tr G-2(\tr P+2b)\ \ge\ 4\tr G-2N(I'),
\]
and (AF3)--(AF4) give $s_1\ge(H_\lambda(\psi)-o(1))N+\tr\Psi(M^\circ)$. Insert Lemma~\ref{lem:transfer} with $\Lambda\le N(1+o(1))$:
$s_1(1-B'/m)\ge(H_\lambda(\psi)-(B'/A')\tau)N-o(N)$. Since $N_0^s(T,2T)\ge s_1-O(\sqrt T\log T)$, the bound holds with $H_\lambda$ and $c'$; let
$c'\uparrow c$ and $\lambda\to1^-$ (\S\ref{ssec:bandwidth}); the right side is continuous in both.
\end{proof}

The quantity $\tr\Psi(M^\circ)$ is the Gram energy among the simple on-line zeros that the argument of~\cite{AF} discards: if every eigenvalue
of $M^\circ$ is at most $2$, then $\tr\Psi(M^\circ)=\HS{M^\circ-I}^2$ up to the diagonal defects $\delta_\rho$. For $\psi=\cos(\alpha s)$ with
$0<\alpha<\pi$ the positive zeros of $k_\psi$ form a sum-free set (by~\eqref{eq:kcos} each of them satisfies $x\tan(\pi x)=c$ with
$c=\frac\alpha{2\pi}\tan\frac\alpha2>0$, and if $x$, $y$ and $x+y$ were all zeros, the addition formula for the tangent would give
$x^2+xy+y^2+c^2=0$), so three simple zeros are never pairwise orthogonal and the energy cannot
vanish locally (for the Montgomery--Taylor window this observation, with its proof, was published by \mbox{ainta}~\cite{ainta} on
11~August~2026, and it is the basis of Devine's computation-free improvement of the Montgomery--Taylor constant~\cite{Devine22287432}); it can vanish for the flat window, whose kernel vanishes on $\Z\setminus\{0\}$. No lower bound for this energy can come from the
prime side alone: the periodic configuration with $\frac46$ of the sites simple, $\frac16$ double and $\frac16$ empty~\cite[\S7.2]{AF} has the
flat-window trace moments and zero energy. The bound is therefore configuration-wise, which is why Theorem~\ref{thm:G} solves an implicit
linear inequality for $s_1$.

\subsection{Proof of Theorem~\ref{thm:S}}
Let
\[
\tilde\psi(s):=\frac{50000}{50037}\Big(1+\frac{127}{2500}\,u-\frac{8997}{10000}\,u^2+\frac{6187}{10000}\,u^3+\frac{3}{1250}\,u^4\Big),\qquad u=(2s)^2,
\]
on $[-\frac12,\frac12]$. It is even and $C^\infty$, with $0.7679\le\tilde\psi\le1$; it is not monotone on $[0,\frac12]$ (the bracket has maximum
$1.000731$ at $u=0.0291$, which is why it is normalised by $\frac{50000}{50037}$, and minimum $0.768517$ at $u=0.9357$). Its window constant is the
exact rational
\[
H(\tilde\psi)=\frac{26957030199857}{40103091391077}=0.6721933213820763\ldots,
\]
with $\int\tilde\psi$, $\int\tilde\psi^2$ and $J(\tilde\psi)$ exact rationals ($J$ through the kernel $\int\abs{s-s'}\tilde\psi(s')\,ds'$, a
polynomial of degree $10$ in $s$ on $\abs s\le\frac12$); this is $2.5\cdot10^{-4}$ below $H(\cos1.48s)$, and the gain comes from the packing
term. Take $K=8$, the exact rational position weights and penalties of the certificate file named in Table~\ref{tab:certs}, with
$\nu=\frac7{1700}$, and $c=\frac{199}{25000}=0.00796$. Then $(\mathrm{LI})$ holds for $k_{\tilde\psi}$ (Section~\ref{sec:cert}). Take $m=147$: then
$A=0.00796\cdot140=\frac{1393}{1250}=1.1144>\frac{147}{146}$, so $B=\Phi_{147}(A)=A-(\sqrt{14527/13125}-1)^2$ lies on the concave branch of
Lemma~\ref{lem:Phi}, and $\tau=\frac{7}{1700}\cdot\frac{140}{147}=\frac1{255}$. Theorem~\ref{thm:G} gives
\[
\frac{H(\tilde\psi)-(B/A)\tau}{1-B/m}=0.673373689541190\ldots\ >\ 0.6733736895 .
\]
In the formal proof the square root is enclosed in an interval of width $10^{-18}$, and the final inequality, linear in the enclosed quantity,
is checked in exact rational arithmetic; the margin is $4.1\cdot10^{-11}$. Since every quantity is homogeneous of degree $0$ in the window, the
formal proof may and does use the window $\frac45v$, with $v$ the bracket above: it satisfies the side conditions of the window class of the
Lean artifact of~\cite{AF} (which $\tilde\psi$, being non-monotone with total variation on $\R$, including the jumps at $\pm\frac12$, of
about $2.016$, does not; for $\frac45v$ it is about $1.614$), and it has the same $H$,
the same kernel and the same certificate. \qed

\medskip
For distinct zeros the same chain gives $\frac12(1+0.6733736895)=0.8366868$ (Remark~\ref{rem:Sdist}), superseded by
Theorem~\ref{thm:main}. Earlier certificates of the same kind gave $0.6732063$ (window $\cos1.48s$, $K=7$) and $0.6733246847$
(window $\tilde\psi$, $K=7$).

\begin{remark}[Credit]\label{rem:credit}
For this statistic the stability inequality with the term $\tr\Psi(M)$, the $K$-point functional with uniform weights $\gamma_{s,i}=2/(K-s)$
and the offset-averaged
aggregation were published by \mbox{ainta}~\cite{ainta}, position weights with span capacity $2$ and non-cosine windows by
\mbox{trmdy}~\cite{trmdy}, both on 11~August~2026, position-dependent penalties by \mbox{sxuff}~\cite{sxuff} (11~August), the exact count of penalties
over shifted blocks (our $\tau$; the identification is ours [M]) in \mbox{AMTOPA}'s work~\cite{AMTOPA} and in \mbox{trmdy} (12~August), eight-point
certificates by \mbox{AMTOPA} (12~August), and the aggregation function $\Phi_m$ by tawanerguo-cn~\cite{TawanRepo} and \mbox{sxuff} (11~August) and \mbox{trmdy}
(12~August), all before our work on it began (\S\ref{ssec:others}). What is ours here, as far as we have found, is the polynomial window with
exact rational constants as a window for $\zeta$ and the box bound LB3 of \S\ref{ssec:bnb}, which we have not compared in substance with the
pruning of other verifiers; Lemma~\ref{lem:Phi} was
published earlier by tawanerguo-cn, \mbox{sxuff} and \mbox{trmdy} (\S\ref{ssec:others}). Several unrefereed works claim higher values
(\S\ref{ssec:others}), among them a Lean theorem announced by Zhaoyi-Tian~\cite{ZhaoyiTian}.
\end{remark}

\section{Simple and distinct zeros}\label{sec:simple}

Throughout this section $\psi(s)=\cos(1.6s)$, a window in the sense of \S\ref{ssec:family} ($\cos0.8\le\psi\le1$), with
$H=H(\psi)=0.67198155100037074\ldots$ and $v_\psi(s)=\cos(8s/5)/(\frac54\sin\frac45)$. The kernel $k=k_\psi$ is given by~\eqref{eq:kcos} with
$\alpha=1.6$.

\subsection{The all-marks local inequality}
\begin{definition}[All-marks local inequality]\label{def:LIm}
Let $K\ge2$, let $\gamma_{s,i}$ and $\mu_i\ge0$ be position weights and penalties as in Definition~\ref{def:LI}, $\nu:=\sum_i\mu_i$, and let
$b_i(1),b_i(2)\in\R$ ($1\le i\le K$) be \emph{site credits}, with no sign condition; put $a_j:=\sum_ib_i(j)$. For a \emph{mark pattern}
$\bm\in\{1,2\}^K$ and $g\in[0,\infty)^{K-1}$ let
\[
F_{\bm}(g):=\sum_{s=1}^{K-1}\sum_{i=1}^{K-s}\gamma_{s,i}\,m_im_{i+s}\,k(g_i+\dots+g_{i+s-1})^2+\sum_{i=1}^{K-1}\mu_ig_i .
\]
We say that $(\mathrm{LI}_{\mathrm m})$ holds for $k$ with data $(\gamma,\mu,b)$ if $F_{\bm}(g)\ge C(\bm):=\sum_ib_i(m_i)$ for every $\bm$ and
every $g\in[0,\infty)^{K-1}$.
\end{definition}

The marks are multiplicities capped at $2$. With $\bm\equiv\mathbf1$ this is Definition~\ref{def:LI} with $c=a_1$.

\subsection{The slack of an on-line set}
Let $\cO$ be a finite set of distinct on-line zeros with $\gamma\in I'$ and multiplicities $m_\rho$. Put $d_\rho:=\norm{u_\rho}^2=1-\delta_\rho$,
\[
G_\cO:=\sum_{\rho\in\cO}m_\rho u_\rho u_\rho^{\tsp},\qquad M_\cO:=\big(\sqrt{m_\rho m_{\rho'}}\,\langle u_\rho,u_{\rho'}\rangle\big)_{\rho,\rho'\in\cO},
\qquad E(\cO):=\sum_{\rho\ne\rho'}(M_\cO)_{\rho\rho'}^2
\]
(ordered pairs), and $\Slack(\cO):=s_1(\cO)-2\tr G_\cO+\HS{G_\cO}^2$, where $s_j(\cO)$ is the number of sites of multiplicity $j$. The matrices
$G_\cO$ and $M_\cO$ have the same nonzero spectrum.

\begin{lemma}[On-line slack identity]\label{lem:onslack}
Exactly,
\begin{gather*}
\Slack(\cO)=\sum_{\rho\in\cO}\beta_\rho+E(\cO),\\
\beta_\rho:=\mathbf 1_{m_\rho=1}-2m_\rho d_\rho+m_\rho^2d_\rho^2\ \ge\ m_\rho(m_\rho-2)\mathbf1_{m_\rho\ge3}-2m_\rho(m_\rho-1)\,\delta_\rho .
\end{gather*}
\end{lemma}

\begin{proof}
$\tr G_\cO=\sum m_\rho d_\rho$ and $\HS{G_\cO}^2=\HS{M_\cO}^2=\sum_\rho m_\rho^2d_\rho^2+E(\cO)$. For $m=1$, $\beta=\delta^2\ge0$. For $m\ge2$,
$\beta=f(d)$ with $f(x)=m^2x^2-2mx$, and $f(1)-f(d)=m(1-d)(m(1+d)-2)\le2m(m-1)\delta$.
\end{proof}

With every $d_\rho=1$ and all zeros on the line, this says that the slack of the rank--trace inequality is exactly the Gram energy among all
on-line zeros plus $m(m-2)$ for each zero of multiplicity $m\ge3$; the stability term of Lemma~\ref{lem:stab} is the part of this energy
among simple zeros.

\begin{proposition}[Aggregation]\label{prop:agg}
Assume $(\mathrm{LI}_{\mathrm m})$ holds for $k_\psi$ with data $(\gamma,\mu,b)$. Let $\cO$ be as above; its Gram entries
$\langle u_\rho,u_{\rho'}\rangle$ are symmetric in $\rho,\rho'$. Let $s_{\ge2}(\cO)$ be the number of its sites with $m_\rho\ge2$ and
$B_b:=\sum_i(\abs{b_i(1)}+\abs{b_i(2)})$. Then
\[
E(\cO)\ \ge\ a_1s_1(\cO)+a_2s_{\ge2}(\cO)-\nu\Lambda-16K^2\Big(\eps_T\abs\cO+\sum_{\rho\in\cO}\delta_\rho\Big)-2KB_b .
\]
\end{proposition}

\begin{proof}
Order $\cO$ by ordinate, cap the marks, $\tilde m:=\min(m,2)$, and let $W$ run over the runs of $K$ consecutive sites, with gap vector $g_W$.
Put $F^\varphi(W):=\sum_{s,i}\gamma_{s,i}\tilde m_{W_i}\tilde m_{W_{i+s}}\langle u_{W_i},u_{W_{i+s}}\rangle^2+\sum_i\mu_ig_{W,i}$, and let
$F^\psi(W)=F_{\tilde\bm_W}(g_W)$ be the same with $k_\psi(x_{W_i}-x_{W_{i+s}})^2$ in place of $\langle u,u'\rangle^2$.
\emph{Upper bound.} A pair of sites at index distance $s<K$ occurs in the runs containing it with total weight at most $\sum_i\gamma_{s,i}=2$, and
$2\tilde m\tilde m'\langle u,u'\rangle^2\le mm'(\langle u,u'\rangle^2+\langle u',u\rangle^2)$ is the contribution of its two ordered pairs to
$E(\cO)$ --- this is where the symmetry is used; each gap receives at most $\nu$, and the gaps add up to at most $\Lambda$. So
$\sum_WF^\varphi(W)\le E(\cO)+\nu\Lambda$.
\emph{Transfer.} By Lemma~\ref{lem:residue}, $\abs{\langle u,u'\rangle-k_\psi}\le\eps_T+(\delta\delta')^{1/2}$ and both are at most $1$ in absolute
value, so $\abs{F^\varphi(W)-F^\psi(W)}\le8\sum_{s,i}\gamma_{s,i}(\eps_T+\frac12(\delta_{W_i}+\delta_{W_{i+s}}))$. A site lies in at most $K$ runs,
and $\sum_{s,i}\gamma_{s,i}=2(K-1)$; hence $\sum_W\abs{F^\varphi-F^\psi}\le16K^2(\eps_T\abs\cO+\sum\delta_\rho)$.
\emph{Lower bound.} $F^\psi(W)\ge C(\tilde\bm_W)$ for every run, and $\sum_WC(\tilde\bm_W)=a_1s_1+a_2s_{\ge2}$ up to the first and last $K-1$
sites, which miss positions; the defect is at most $2KB_b$.
\end{proof}

The symmetry cannot be dropped: for a general matrix $C$ in place of the Gram entries, and a kernel that is not even, the inequality fails
(with $K=2$, $\gamma_{1,1}=2$, $\mu_1=1$, $b\equiv1$, the kernel $k(t)=\sqrt{1-t/2}$ on $[0,2]$ and $0$ elsewhere, including $t<0$, and $100$ simple
sites at spacing $1.5$ with entries $C_{pq}=k(x_p-x_q)$, the right side of the inequality is $\frac{71}2$ while $E=\frac{99}4$; with $C$
symmetrised, $E=\frac{99}2$).

\subsection{Removing off-line pairs and heavy zeros}\label{ssec:removal}
Split the zeros with $\Re\gamma\in I'$ into the \emph{light} on-line zeros $\cO_L$ (multiplicity $1$ or $2$; $s_1$ and $s_2$ of them), with
$G_L:=G_{\cO_L}$ and $M_L:=M_{\cO_L}$; the \emph{heavy} on-line zeros (multiplicity $\ge3$; $s_H$ distinct ones, $N_H$ with multiplicity), with
$G_H:=\sum m_\rho u_\rho u_\rho^{\tsp}$; and the off-line pairs ($p$ of them, $p_1$ simple, $N_{\mathrm{off}}$ zeros with multiplicity), with
$G_{\mathrm{off}}:=G-G_L-G_H$. Put $Q:=G_H+G_{\mathrm{off}}$, $b_Q:=s_H+p$ and $N_Q:=N_H+N_{\mathrm{off}}\ge3s_H+2p$, so that
$N(I')=s_1+2s_2+N_Q$.

\begin{lemma}[Pair and heavy removal]\label{lem:removal}
$n_+(Q)\le b_Q$ and $n_-(Q)\le p$. For a count $X$ of zeros with $\Re\gamma\in I'$ let $X_L$ be its part carried by $\cO_L$,
$\Slack_X:=X-2\tr G+\HS G^2$ and $\Slack^L_X:=X_L-2\tr G_L+\HS{G_L}^2$. Then
\begin{gather*}
\Slack_X\ \ge\ \Slack^L_X+(X-X_L)+2(N_Q-2b_Q)-2E_T-\sum_{i\le p}\big(\lambda_i(M_L)-2\big)_+^2,\\
E_T:=(N(I')-\tr G)_+=o(N).
\end{gather*}
\end{lemma}

\begin{proof}
Each pair block has $n_\pm\le1$ (AF2), $G_H\succeq0$ has rank at most $s_H$, and $n_\pm$ are subadditive; $G_H$ adds nothing to $n_-$. By (AF1),
$\tr G_L\le s_1+2s_2$, so $\tr Q=\tr G-\tr G_L\ge N_Q-E_T$, and $E_T=o(N)$ by (AF3). Apply Lemma~\ref{lem:pairrem} with $G=G_L$, whose nonzero
eigenvalues are those of $M_L$; its right side is increasing in $\tr Q$.
\end{proof}

Only (AF1), (AF3) and the inertia of the pair blocks enter; no pair vector is ever evaluated, so no estimate at complex arguments is needed.

\subsection{The on-line spectral condition}
Given claims $a_1,a_2$ with $a_1<1$, let $c_*:=2-2a_1$, $\lambda_c:=2+\sqrt{c_*}$, and let $\kappa$, $\ch$ be as in Theorem~\ref{thm:loc}. Put
\[
\mg_L:=\Slack(\cO_L)-a_1s_1-a_2s_2+\nu\Lambda .
\]
The \emph{on-line spectral condition} is the statement $\ch(M_L)\le\mg_L+o(N)$. It is proved in Theorem~\ref{thm:OLL} from three
ingredients: room on separated sets, tight chains, and aggregation. The device of Lemmas~\ref{lem:room} and~\ref{lem:tight}, close pairs
paid for locally and a separated set whose Gram norm is bounded by a Fourier majorant, was published earlier, for simple zeros on the line
and with a sinc-square majorant, by \mbox{gfreund123} in pull request~\#888 to \mbox{GettysburgResearch/riemann} (14~September~2026)~\cite{Gettysburg888}
[V$\sim$]; the priority is theirs.

\begin{lemma}[Separated sites have room]\label{lem:room}
Let $\cB$ be an even, integrable function on $\R$ with $\cB\ge0$ on $\R$, $\cB\ge v_\psi$ on $[-\frac12,\frac12]$, and
$\widehat\cB(t):=\int\cB(s)\cos(2\pi ts)\,ds=0$ for $\abs t\ge1$; let $\beta_*\ge\sup_{3/4\le t\le1}\abs{\widehat\cB(t)}$. If $\cR\subset\cO_L$
has normalised ordinates pairwise at least $\frac34$ apart, then
$\lambda_{\max}(M_\cR)\le(1+\eps'_T)\,(2\widehat\cB(0)+4\beta_*)$.
\end{lemma}

\begin{proof}
For $y\in\R^\cR$, by Lemma~\ref{lem:residue}(i)--(ii) (dropping $-\widehat E\preceq0$),
\begin{align*}
\langle y,M_\cR y\rangle&\le\int v_\varphi\Big|\sum_ay_a\sqrt{m_a}\,e(x_as)\Big|^2ds\\
&\le(1+\eps'_T)\int_\R\cB\,\Big|\sum_ay_a\sqrt{m_a}\,e(x_as)\Big|^2ds=(1+\eps'_T)\langle y,\mathcal Ny\rangle,
\end{align*}
with $e(u)=e^{2\pi iu}$ and $\mathcal N_{ab}:=\sqrt{m_am_b}\,\widehat\cB(x_a-x_b)$. Since $\widehat\cB$ vanishes on $\abs t\ge1$ and the sites are
pairwise $\ge\frac34$ apart, each site has at most one other site within distance $<1$ on each side, at distance in $[\frac34,1)$. So each row of
$\mathcal N$ has diagonal entry $\le2\widehat\cB(0)$ and at most two off-diagonal entries, each of absolute value $\le2\beta_*$; Gershgorin gives
$\lambda_{\max}(\mathcal N)\le2\widehat\cB(0)+4\beta_*$.
\end{proof}

\paragraph{The majorant.} Let $\beta_0,\dots,\beta_{59}$ be sixty binary64 numbers, read as exact rationals (the certificate file of
Table~\ref{tab:certs}), and
\[
\cB(s):=\tfrac1{60}\sinc(\pi s/60)^2\,P(s),\qquad P(s):=\beta_0+2\sum_{j=1}^{59}\beta_j\cos(\pi js/30).
\]
From $\int\sinc(\pi u)^2\cos(2\pi\omega u)\,du=(1-\abs\omega)_+$ one gets $\widehat\cB(t)=\sum_{\abs j\le59}\beta_{\abs j}(1-\abs{60t-j})_+$, the
piecewise-linear interpolant of the $\beta_j$ at the nodes $j/60$, which vanishes for $\abs t\ge1$; so $\int\cB=\widehat\cB(0)=\beta_0=1.5556524027186915$,
and, since $\frac34=\frac{45}{60}$, $\sup_{[3/4,1]}\abs{\widehat\cB}=\max_{45\le j\le59}\abs{\beta_j}=\abs{\beta_{52}}=0.0237647699514\ldots=:\beta_*$. In
exact rational arithmetic $2\beta_0+4\beta_*=3.2063638852\ldots<\frac{33}{10}$. The two remaining conditions, $P\ge0$ on $[0,30]$ ($P$ is even and
$60$-periodic; its minimum is about $1.12\cdot10^{-3}$, near $s=3.50$) and $\cB\ge v_\psi$ on $[0,\frac12]$ (margin about $1.95\cdot10^{-5}$, near
$s=0.4994$), are finite computations, certified in Section~\ref{sec:cert}. Hence, as soon as $\eps'_T\le\frac1{40}$,
\begin{equation}\label{eq:room}
\lambda_{\max}(M_\cR)\le\tfrac{41}{40}\cdot\tfrac{33}{10}<3.4\le\lambda_c
\end{equation}
for the claims used below ($c_*\ge1.96$).

\begin{lemma}[Tight chains pay for themselves]\label{lem:tight}
Let $\cT\subset\cO_L$ be a union of $n_{\mathrm{ch}}$ maximal \emph{tight chains}: maximal runs of consecutive sites with consecutive gaps
$<\frac34$, each with at least two sites. Let $k_*\ge0$ with $k_\varphi(t)\ge k_*$ for $0\le t<\frac34$, and put
\[
\varpi(k):=\min\Big\{2k^2-2a_1,\ 4k^2-2a_2,\ \big(\tfrac12+(\tfrac14+2k^2)^{1/2}\big)^2-1-a_1-a_2\Big\}.
\]
Assume $a_1,a_2\ge0$ and $\varpi(k_*)\ge0$. Then
\[
\Delta_\cT:=\tr(2I-M_\cT)_+^2-(1+a_1)s_1(\cT)-a_2s_2(\cT)\ \ge\ n_{\mathrm{ch}}\big(\varpi(k_*)-\max(a_1,a_2)\big).
\]
\end{lemma}

\begin{proof}
$f(x):=(2-x)_+^2$ is convex and nonincreasing. Split each chain greedily into consecutive pairs and at most one singleton. By
Lemma~\ref{lem:pinch}, $\tr f(M_\cT)\ge\sum_w\tr f(M_w)$ over the blocks. By Lemma~\ref{lem:residue}(i),
$M_w=D_w^{1/2}(K^\varphi_w-\widehat E_w)D_w^{1/2}\preceq M'_w:=D_w^{1/2}K^\varphi_wD_w^{1/2}$ with $D_w=\operatorname{diag}(m)$ and $K^\varphi_w$ the
kernel matrix; as $f$ is nonincreasing, Weyl gives $\tr f(M_w)\ge\tr f(M'_w)$. For a pair at gap $<\frac34$ with $k:=k_\varphi(\text{gap})\in[k_*,1]$,
the eigenvalues of $M'_w$ are $1\pm k$, $2\pm2k$, $\frac32\pm(\frac14+2k^2)^{1/2}$ for marks $(1,1)$, $(2,2)$, $(1,2)$, so $\tr f(M'_w)$ minus the
claims of the block is $2k^2-2a_1$, $4k^2-2a_2$, $(\frac12+(\frac14+2k^2)^{1/2})^2-1-a_1-a_2$ respectively, each increasing in $k^2$, hence
$\ge\varpi(k_*)$. A singleton costs at least $-a_1$ (simple: $f(d)\ge1$) or $-a_2$ (double: $f(2d)\ge0$). A chain with $q\ge1$ pairs therefore
contributes at least $q\varpi(k_*)-\max(a_1,a_2)\ge\varpi(k_*)-\max(a_1,a_2)$ if it has a singleton (using $\varpi(k_*)\ge0$), and $q\varpi(k_*)\ge\varpi(k_*)-\max(a_1,a_2)$
if not (using $a_1,a_2\ge0$).
\end{proof}

Both hypotheses are needed. Take $k_*=0$, the kernel equal to $1$ at $0$ and $0$ elsewhere, and $\widehat E=0$. With the $K=7$ claims (so
$a_1,a_2>0$ but $\varpi(0)<0$) and one chain of six simple sites at gaps $\frac12$, the left side is $-6a_1=-0.10949\ldots$ and the right side
$-3a_2=-0.07008\ldots$. With $a_1=a_2=-\frac1{10}$ (so $\varpi(0)=\frac15\ge0$) and one chain of two simple sites, the left side is $\frac15$ and the right
side $\frac3{10}$. Since $k_\psi$ is nonincreasing on
$[0,1]$ ($k_\psi'(t)=-\int2\pi\sigma v_\psi(\sigma)\sin(2\pi t\sigma)\,d\sigma\le0$ there), Lemma~\ref{lem:residue}(ii) allows
$k_*=k_\psi(\frac34)-\eps_T$, with $k_\psi(\frac34)=0.3556941746\ldots$. At $\eps_T=0$ the three entries of $\varpi$ are $0.2165399517$, $0.4593506235$ and
$0.4206777378$ for the $K=7$ claims, and $\varpi-\max(a_1,a_2)=0.1931785717$ ($0.2118330717$ for $K=5$).

\begin{theorem}[On-line spectral condition]\label{thm:OLL}
Assume $(\mathrm{LI}_{\mathrm m})$ for $k_\psi$ with data $(\gamma,\mu,b)$, with $a_1,a_2\ge0$, $\lambda_c>\frac{33}{10}$ and
$\max(a_1,a_2)<\varpi(k_\psi(\frac34))$. Let $\cT$ be the set of light sites that have another light site within normalised distance $<\frac34$, and $\cR:=\cO_L\setminus\cT$.
Then, exactly,
\begin{equation}\label{eq:OLLid}
\mg_L-\Big[\HS{M_\cT-2I}^2+2\HS{M_{\cT\cR}}^2-\tr(2I-M_\cT)_+^2\Big]=\mg(\cR)+\Delta_\cT-2\sum_{\rho\in\cT}m_\rho\delta_\rho,
\end{equation}
with $\mg(\cR):=\Slack(\cR)-a_1s_1(\cR)-a_2s_2(\cR)+\nu\Lambda$; and for $T\ge T_0$,
\[
\ch(M_L)\ \le\ \mg_L+(16K^2+4)\sum_{\rho\in\cO_L}m_\rho\delta_\rho+16K^2\eps_T\abs{\cO_L}+2KB_b=\mg_L+o(N).
\]
\end{theorem}

\begin{proof}
\emph{Structure.} Any two sites of $\cR$ are at least $\frac34$ apart. Every site of $\cT$ has a site of $\cT$ within $\frac34$ that is adjacent to it in
the order of $\cT$ (a closer partner lies in between), so $\cT$ is the disjoint union of its maximal tight chains.
\emph{Localisation.} By Lemma~\ref{lem:room} and the majorant, $\lambda_{\max}(M_\cR)\le(1+\eps'_T)\frac{33}{10}<\lambda_c$ once $\eps'_T$ is small. After a permutation of indices, under which
$\tr f$ is invariant, $M_L=\begin{pmatrix}M_\cT&M_{\cT\cR}\\M_{\cR\cT}&M_\cR\end{pmatrix}$, and Theorem~\ref{thm:loc} bounds $\ch(M_L)$ by the bracket
in~\eqref{eq:OLLid}.
\emph{The identity.} $E(\cO_L)=E(\cT)+E(\cR)+2\HS{M_{\cT\cR}}^2$ (ordered pairs), the claims are additive, $\HS{M_\cT-2I}^2=\sum_{\rho\in\cT}(m_\rho d_\rho-2)^2+E(\cT)$,
and $(md-2)^2-\beta_\rho=\mathbf1_{m=1}+2m\delta$. Insert Lemma~\ref{lem:onslack} for $\Slack(\cO_L)$ and $\Slack(\cR)$.
\emph{The bound.} Take $k_*:=k_\psi(\frac34)-\eps_T$ in Lemma~\ref{lem:tight} (valid since $k_\psi$ is nonincreasing on $[0,1]$); for $\eps_T$
small, $\varpi(k_*)\ge\max(a_1,a_2)\ge0$ by the hypothesis and continuity, so $\Delta_\cT\ge0$. By Lemma~\ref{lem:onslack} ($\beta_\rho\ge-4\delta_\rho$
for $m_\rho\le2$) and Proposition~\ref{prop:agg} applied to $\cO=\cR$ (removing sites only lengthens gaps, and $(\mathrm{LI}_{\mathrm m})$ holds for all
$g\ge0$), $\mg(\cR)\ge-4\sum_\cR\delta_\rho-16K^2(\eps_T\abs\cR+\sum_\cR\delta_\rho)-2KB_b$. Finally $2\sum_\cT m_\rho\delta_\rho\le4\sum_\cT\delta_\rho$,
and the error is $o(N)$ by Lemma~\ref{lem:residue}(iii)--(iv) and $\eps_T\to0$.
\end{proof}

The truncation matrix $\widehat E$ enters only with the favourable sign and through $\abs{\widehat E_{\rho\rho'}}\le(\delta_\rho\delta_{\rho'})^{1/2}$, and
the charge is computed on the true matrix, so no perturbation bound for $M_L$ is needed. For the certificates used, the formal proof needs of
the numerical constants only $k_\psi(\frac34)\ge\frac14$ (proved by the fundamental theorem of calculus and $3.14<\pi<3.15$; already $\varpi(\frac14)>\max(a_1,a_2)$
for these claims), $\lambda_c\ge3.4$ and $2\beta_0+4\beta_*\le\frac{33}{10}$, not the more precise values quoted above.

\subsection{Proof of Theorem~\ref{thm:main}}
\begin{theorem}\label{thm:SigmaD}
Let $\psi=\cos(1.6s)$ and suppose $(\mathrm{LI}_{\mathrm m})$ holds for $k_\psi$ with data $(\gamma,\mu,b)$ whose claims $a_1,a_2,\nu$ satisfy:
$a_1,a_2,\nu\ge0$; $\lambda_c=2+\sqrt{2-2a_1}>\frac{33}{10}$; $\max(a_1,a_2)<\varpi(k_\psi(\frac34))$; $1-a_1+\frac{a_2}2>0$; and, for all integers $m\ge2$
and $m'\ge3$,
\begin{gather*}
(4-a_2)m-4-c_*\ge0,\quad 2m(1-a_2+a_1)-4a_1+2a_2-c_*\ge0,\\
(2-\tfrac{a_2}2)m'-4\ge0,\quad m'(1-a_2+a_1)-2-2a_1+a_2\ge0 .
\end{gather*}
Then
\begin{gather*}
\liminf_{T\to\infty}\frac{N^s(T,2T)}{N(T,2T)}\ge\Sigma:=\frac{H+\frac{a_2}2-\nu}{1-a_1+\frac{a_2}2},\\
\liminf_{T\to\infty}\frac{N^d(T,2T)}{N(T,2T)}\ge D:=\frac{1+H+a_2-a_1-\nu}{2-2a_1+a_2}=\frac{1+\Sigma}2 .
\end{gather*}
\end{theorem}

\begin{proof}
Fix $\lambda\in(0,1)$ and let $X$ be $s_1+2p_1$ (the simple zeros with $\Re\gamma\in I'$: a simple off-line zero and its partner are both simple) or
$s_1+s_2+s_H+2p$ (the distinct ones); they differ from $N^s(T,2T)$, $N^d(T,2T)$ by $O(\sqrt T\log T)$.

\emph{Step 1.} By definition $X=2\tr G-\HS G^2+\Slack_X$, so (AF3) and (AF4) give $X\ge(H_\lambda-o(1))N+\Slack_X$.

\emph{Step 2.} By Lemma~\ref{lem:removal}, the inequality $(\lambda-2)_+^2\le c_*+\kappa(\lambda)$, and Theorem~\ref{thm:OLL},
\[
\Slack_X\ \ge\ \big(\Slack^L_X-\Slack(\cO_L)\big)+a_1s_1+a_2s_2-\nu\Lambda+(X-X_L)+2(N_Q-2b_Q)-c_*p-o(N),
\]
where $\Slack^L_X-\Slack(\cO_L)=X_L-s_1$ is $0$ for simple zeros and $s_2$ for distinct zeros.

\emph{Step 3.} For simple zeros compare the right side with $a_1X+\frac{a_2}2(N(I')-X)-\nu\Lambda$; for distinct zeros with
$(2a_1-1-a_2)X+(1+a_2-a_1)N(I')-\nu\Lambda$. On the light zeros the two sides agree. Using $N(I')=s_1+2s_2+N_Q$, the remaining difference (available
minus required) is a sum over objects:
\begin{itemize}
\item a simple off-line pair: $2-2a_1-c_*=0$, for both counts;
\item an off-line pair of multiplicity $m\ge2$: $(4-a_2)m-4-c_*$, resp.\ $2m(1-a_2+a_1)-4a_1+2a_2-c_*$;
\item a heavy on-line zero of multiplicity $m\ge3$: $(2-\frac{a_2}2)m-4$, resp.\ $m(1-a_2+a_1)-2-2a_1+a_2$;
\end{itemize}
all nonnegative by hypothesis. So $\Slack_X\ge a_1X+\frac{a_2}2(N(I')-X)-\nu\Lambda-o(N)$, resp.\
$\Slack_X\ge(2a_1-1-a_2)X+(1+a_2-a_1)N(I')-\nu\Lambda-o(N)$.

\emph{Step 4.} Insert into Step~1 with $N(I')=N+o(N)$ and $\Lambda\le N(1+o(1))$ (and $\nu\ge0$):
$X(1-a_1+\frac{a_2}2)\ge(H_\lambda+\frac{a_2}2-\nu)N-o(N)$, resp.\ $X(2-2a_1+a_2)\ge(1+H_\lambda+a_2-a_1-\nu)N-o(N)$. Divide by the positive
denominators and let $\lambda\to1^-$ (\S\ref{ssec:bandwidth}).
\end{proof}

\begin{proof}[Proof of Theorem~\ref{thm:main}]
The $K=5$ and $K=7$ certificates of Table~\ref{tab:allmarks} satisfy $(\mathrm{LI}_{\mathrm m})$ for $k_\psi$ (Section~\ref{sec:cert}) and the conditions of
Theorem~\ref{thm:SigmaD} with wide margins: $\lambda_c\ge3.401$, $\varpi(k_\psi(\frac34))-\max(a_1,a_2)\ge0.193$, and every entry of the table in
Step~3 for a multiple off-line pair or a heavy on-line zero is at least $0.97$ (the entry for a simple off-line pair is exactly $0$, by the
choice of $c_*$). With $H=H(\cos1.6s)$, enclosed in an interval of width
$4\cdot10^{-21}$ (the formal proof uses an enclosure of width $10^{-16}$ from alternating Taylor series), the constants are those of
Table~\ref{tab:allmarks}.
\end{proof}

\begin{table}[ht]\centering\small
\caption{The all-marks certificates for $\psi=\cos1.6s$ ($\nu=(K-1)/500$). The digits of $\Sigma$ and $D$ are truncated. The $K=6$ row is
given for comparison; its certificate is not used and is not listed in Table~\ref{tab:certs}.}\label{tab:allmarks}
\begin{tabular}{@{}rllllll@{}}
\toprule
$K$ & $a_1$ & $a_2$ & $\nu$ & $\Sigma$ & $D$ & $\lambda_c$\\
\midrule
$5$ & $\frac{1280197}{10^8}$ & $\frac{48749}{3125000}$ & $\frac1{125}$ & $0.675158622199$ & $0.837579311099$ & $3.40513$\\[2pt]
$6$ & $\frac{195379}{12500000}$ & $\frac{976263}{50000000}$ & $\frac1{100}$ & $0.675709032131$ & $0.837854516065$ & $3.40312$\\[2pt]
$7$ & $\frac{1824837}{10^8}$ & $\frac{1168069}{5\cdot10^7}$ & $\frac3{250}$ & $0.676102666964$ & $0.838051333482$ & $3.40125$\\
\bottomrule
\end{tabular}
\end{table}

\begin{remark}[What is not claimed]\label{rem:notS}
The proportion of zeros that are simple \emph{and} on the line is not improved by this route. For it a simple off-line pair brings no bonus, the
threshold in the charge would be $2$ instead of $\lambda_c$, and Theorem~\ref{thm:loc} would need $\lambda_{\max}(M_\cR)\le2$, which fails already
for unit-spaced double zeros (two of them give $\lambda_{\max}=2(1+k_\psi(1))=2.1387$, and long runs approach $2/\!\int\psi=2.2304$).
\end{remark}

\section{Zeros that are simple or on the line}\label{sec:sc}

\begin{proof}[Proof of Corollary~\ref{cor:SC}]
Use the window $\tilde\psi$ and the data of the proof of Theorem~\ref{thm:S} throughout; the cancellation of $R$ below needs the same window in
(AF4) and in the packing term. Fix $\lambda\in(0,1)$ and let $t:=2+\sqrt2$, so that $t^2=4t-2$ and $t\ge2$. Let $N_{\mathrm{on}}$ be the number of
on-line zeros with $\gamma\in I'$ counted with multiplicity, $r\le N_{\mathrm{on}}$ the number of distinct ones, and $p_1$, $p_2$ the numbers of
simple and of multiple off-line pairs with $\Re\gamma\in I'$, so that $p=p_1+p_2$. The count to be bounded is $N^{\mathrm{sc}}(I'):=N_{\mathrm{on}}+2p_1$,
which differs from $N^{\mathrm{sc}}(T,2T)$ by $O(\sqrt T\log T)$.

Apply Lemma~\ref{lem:Ast} with $(s,t)=(1,2+\sqrt2)$, $P:=\sum_{\text{on-line}}m_\rho u_\rho u_\rho^{\tsp}=VV^{\tsp}$, where $V$ has the $r$ real columns
$\sqrt{m_\rho}\,u_\rho$, $M:=V^{\tsp}V$, and $Q:=G-P$, the sum of the off-line pair blocks, so that $n_+(Q)\le b:=p$ by (AF2). By (AF1),
$\tr P\le N_{\mathrm{on}}$, hence $2\tr P+2t\tr Q=2\tr G+2(t-1)\tr Q\ge2\tr G+2(t-1)(\tr G-N_{\mathrm{on}})$, and (AF3)--(AF4) give
\[
r+t^2p\ \ge\ (2t-R_\lambda(\tilde\psi))N-2(t-1)N_{\mathrm{on}}+\tr\Psi_{1,t}(M)-o(N).
\]
Since $N(I')-N_{\mathrm{on}}\ge2p_1+4p_2$, we have $p\le\frac14(N(I')-N_{\mathrm{on}}+2p_1)$; also $r\le N_{\mathrm{on}}$. Moving the
$N_{\mathrm{on}}$ terms to the left, their coefficient is $2t-1-t^2/4=t^2/4$, so
\[
\tfrac{t^2}4\,N^{\mathrm{sc}}(I')\ \ge\ \big(2t-R_\lambda(\tilde\psi)-\tfrac{t^2}4\big)N+\tr\Psi_{1,t}(M)-o(N).
\]
Pinch $M$ to the block $M^\circ$ of the simple on-line zeros (Lemma~\ref{lem:pinch}; $\Psi_{1,t}$ is convex) and use $\Psi_{1,t}\ge\Psi\ge0$:
$\tr\Psi_{1,t}(M)\ge\tr\Psi(M^\circ)$. For $c'<c$ let $A',B'$ be as in Lemma~\ref{lem:transfer}, $\beta':=B'/m$, $\tau'':=(B'/A')\tau$ and
$S_{c'}:=(H_\lambda(\tilde\psi)-\tau'')/(1-\beta')$. Lemma~\ref{lem:transfer} gives $\tr\Psi(M^\circ)\ge\beta's_1-\tau''N-o(N)$, and the proof of
Theorem~\ref{thm:G} gives $s_1\ge(S_{c'}-o(1))N$. Since $\beta'>0$ and $\beta'S_{c'}-\tau''=S_{c'}-H_\lambda(\tilde\psi)$,
\[
\tr\Psi(M^\circ)\ \ge\ \big(S_{c'}-H_\lambda(\tilde\psi)\big)N-o(N),\qquad H_\lambda(\tilde\psi)=2-R_\lambda(\tilde\psi).
\]
Substituting, $R_\lambda(\tilde\psi)$ cancels:
\[
N^{\mathrm{sc}}(I')\ \ge\ \frac{8t-t^2-8+4S_{c'}}{t^2}\,N-o(N)=\frac{2t-3+2S_{c'}}{2t-1}\,N-o(N).
\]
Let $T\to\infty$, then $c'\uparrow c$ and $\lambda\to1^-$; $S_{c'}$ tends to the constant $S^*=0.673373689541190\ldots$ of Theorem~\ref{thm:S}, and
$(1+2\sqrt2+2S^*)/(3+2\sqrt2)=0.887919569562\ldots>0.8879195$.
\end{proof}

\begin{remark}\label{rem:SC}
\begin{enumerate}[label=(\roman*)]
\item The corollary follows from the \emph{proof} of Theorem~\ref{thm:S} (the linear packing bound, the fixed-point identity and the limit
$c'\uparrow c$), not from the number $S^*$ alone; it uses no input beyond those of Theorem~\ref{thm:S} and Lemma~\ref{lem:Ast}.
\item The formula $(1+2\sqrt2+2S)/(3+2\sqrt2)$ is Lamzouri's~\cite{Lamzouri}, with Lamzouri's constant replaced by $S^*$. The value moves only through $S$,
with slope $2/(3+2\sqrt2)=0.343146$: $0.8876200$ at $S=H(\cos\sqrt2s)$ (Lamzouri's value) and $0.8879195$ at $S^*$. The whole gain,
$2.9\cdot10^{-4}$, comes from Theorem~\ref{thm:S}, so Corollary~\ref{cor:SC} is not an independent advance. At $S=0.67349239$, the largest
unrefereed bandwidth-one value of \S\ref{ssec:others}, the formula gives $0.8879603$.
\item The variant that counts a multiple on-line zero once is not claimed. Within the two-parameter family of Lemma~\ref{lem:Ast}, the choice
$s=1$, $t=2+\sqrt2$ is optimal for this argument, by a computation that we do not reproduce; nothing
above depends on it.
\end{enumerate}
\end{remark}

\begin{remark}[Distinct zeros from the chain of Theorem~\ref{thm:S}]\label{rem:Sdist}
The same argument gives
\[
\liminf_{T\to\infty}\frac{N^d(T,2T)}{N(T,2T)}\ \ge\ \tfrac12(1+S^*)=0.8366868\ldots,
\]
which Theorem~\ref{thm:main} supersedes. We
include it because it is the value that a bound of the kind of Theorem~\ref{thm:S} gives for distinct zeros. The remark is not formalised.
\emph{Proof.} Fix $\lambda\in(0,1)$ and $c'<c$, with $\beta'$, $\tau''$, $S_{c'}$ as above. Apply Lemma~\ref{lem:stab} with $V$ the matrix of the $r$
real columns $\sqrt{m_\rho}\,u_\rho$ over the distinct on-line zeros with $\gamma_\rho\in I'$, $P:=VV^{\tsp}$, $M:=V^{\tsp}V$, $Q:=G-P$ and $b:=p$
(AF2). Let $N_{\mathrm{off}}:=N(I')-N_{\mathrm{on}}\ge2p$. By (AF1) and (AF3), $\tr Q=\tr G-\tr P\ge N_{\mathrm{off}}-o(N)$, so
$2\tr P+4\tr Q=2\tr G+2\tr Q$ and (AF4) give
\[
r\ \ge\ H_\lambda(\tilde\psi)N+2N_{\mathrm{off}}-4p+\tr\Psi(M)-o(N).
\]
Pinch $M$ into the block $M^\circ$ of the simple on-line zeros and the $1\times1$ blocks $m_\rho\norm{u_\rho}^2=m_\rho(1-\delta_\rho)$ of the
multiple ones (Lemma~\ref{lem:pinch}). Since $\Psi$ is $2$-Lipschitz on $[0,\infty)$ and $\Psi(m)=2m-3$ for $m\ge2$,
$\Psi(m_\rho(1-\delta_\rho))\ge2m_\rho-3-2m_\rho\delta_\rho$, and $\sum_\rho m_\rho\delta_\rho=o(N)$; and since
$\tr\Psi(M^\circ)\ge(S_{c'}-H_\lambda(\tilde\psi))N-o(N)$ as shown above,
\[
N^d(I'):=r+2p\ \ge\ S_{c'}N+2N_{\mathrm{off}}-2p+\sum_{m_\rho\ge2}(2m_\rho-3)-o(N),
\]
the sum running over the multiple on-line zeros. Add the identity $N^d(I')=s_1+n_{\ge2}+2p$, where $n_{\ge2}$ is the number of multiple
on-line zeros, and subtract $N(I')=s_1+\sum_{m_\rho\ge2}m_\rho+N_{\mathrm{off}}$:
\[
2N^d(I')-N(I')\ \ge\ S_{c'}N+N_{\mathrm{off}}+\sum_{m_\rho\ge2}(m_\rho-2)-o(N)\ \ge\ S_{c'}N-o(N).
\]
Since $N(I')=N+o(N)$ and $N^d(T,2T)\ge N^d(I')-O(\sqrt T\log T)$, letting $T\to\infty$, then $c'\uparrow c$ and $\lambda\to1^-$ gives the claim.
With $\tr\Psi(M^\circ)\ge0$ in place of the packing bound the same computation gives $\frac12(1+H(\psi))$, the value
of~\cite[Theorem~A(ii)]{AF}.
\end{remark}

\section{The certificates}\label{sec:cert}

\subsection{What is certified}
Four finite statements enter the proofs. Three are local inequalities for all gap vectors in an orthant, one is the majorant.
\begin{itemize}
\item \texttt{CertS8}: $(\mathrm{LI})$ of Definition~\ref{def:LI} for $k_{\tilde\psi}$, $K=8$, with position weights and penalties summing to
$\nu=\frac7{1700}$, at $c=\frac{199}{25000}$ (Theorem~\ref{thm:S}, Corollary~\ref{cor:SC}).
\item \texttt{CertAM5} and \texttt{CertAM7}: $(\mathrm{LI}_{\mathrm m})$ of Definition~\ref{def:LIm} for $k_{\cos1.6s}$ with the claims of
Table~\ref{tab:allmarks} for $K=5$ and $K=7$ (Theorem~\ref{thm:main}); a $K=6$ certificate, shown in Table~\ref{tab:allmarks} for comparison, is not used.
\item \texttt{CertMajV2}: the majorant of \S\ref{sec:simple}: $P\ge0$ on $[0,30]$ and $\cB\ge v_\psi$ on $[0,\frac12]$, with $2\beta_0+4\beta_*<\frac{33}{10}$.
(The library also contains a first version, \texttt{CertMaj}, with the threshold $3.2063638853$, only $5.7\cdot10^{-11}$ above the
actual value; it is proved as well, and no theorem uses it.)
\end{itemize}
The data of the local inequalities (weights, penalties, credits) were chosen by linear programmes over finite sets of gap vectors found by
local minimisation, and the claims were then set slightly below the programme's values: for $K=5$ each claim lies $1.5\cdot10^{-5}$ below the programme's value,
and for $K=7$ every claim was lowered by $10^{-5}$ after the certifier had refuted the programme's own claims (\S\ref{ssec:bnb}). The
certificates, not the programmes, carry the proof. In the formalisation each is a single \texttt{Prop} (Appendix~\ref{app:lean}). The local ones are existential in the weights: they assert
\emph{some} data with the stated $\nu$ and credit sums $a_1,a_2$, since the constants depend on the data only through these numbers; the files
of Table~\ref{tab:certs} supply the witnesses. Gap vectors range over the whole orthant $[0,\infty)^{K-1}$; the reductions used by the
certifiers (to a box, and to classes of patterns) are theorems, not part of the statements.

\begin{table}[ht]\centering\small
\caption{Certificates, their data and their verification. Paths are relative to the folder
\texttt{supplementary/certificates/} of the repository (\S\ref{ssec:reproduce}). The seven sha256 values printed here are those of the
files in \texttt{supplementary/}; those of the three local certificates are also recorded in the doc-strings of their \texttt{Prop}s
(Appendix~\ref{app:defs}).}\label{tab:certs}
\footnotesize
\begin{tabular}{@{}p{1.3cm}>{\raggedright\arraybackslash}p{7.9cm}>{\raggedright\arraybackslash}p{5.8cm}@{}}
\toprule
name & data and sha256 & verification\\
\midrule
\texttt{CertS8} & \texttt{CertS8/\allowbreak w\_poly8A\_K8\_mu1700.json}\newline \texttt{257b340d9da4f343600a21589f16d891}\newline\texttt{5f3e082499d475e153c5453980671e78}\newline the window, \texttt{win\_poly8A.json} of the same folder (the unnormalised $v$):\newline \texttt{c4005f57a7920f78aa2df33b093c9c50}\newline\texttt{30efac2b8e39fc6cf7371b0c0f32dc14} &
interval branch-and-bound with LB3, about $4.5\cdot10^8$ boxes in $8$ shards; re-run from scratch ($4.51\cdot10^8$ boxes). In Lean: displayed hypothesis.\\[4pt]
\texttt{CertAM5} & \texttt{CertAM5/\allowbreak claims\_d15e6\_a1.6\_K5\_mu500.json}\newline \texttt{cf055dd2ee9a9e469b6510fe70b5fbc4}\newline\texttt{16b8a40515cc6746bc390d40d88777f2}\newline the $20$ files \texttt{xpat\_K5\_*.json} of the same folder, concatenated in sorted order:\newline \texttt{a5947c01d34ddcbc7366861b35e83942}\newline\texttt{099fcbf7ac71e13fcd432796882dd34e} &
$20$ classes; $1.18\cdot10^7$ boxes (LB1--LB2), $3.7\cdot10^5$ (LB3); re-run with identical counts. In Lean: replayed by the kernel ($46{,}047$ leaves, $89{,}362$ point records; \S\ref{ssec:replay}), so the $K=5$ pair has no hypothesis; the replay files are the library \texttt{ZetaSReplay} of the repository.\\[4pt]
\texttt{CertAM7} & \texttt{CertAM7/\allowbreak claims\_a1.6\_K7\_mu500\_L1e5.json}\newline \texttt{d2cf393019407bdc797c73147ec8e3dc}\newline\texttt{66ec41fe745600c84173f74ac66bfc5c}\newline the $72$ files \texttt{xpat\_K7\_*.json} of the same folder, concatenated in sorted order:\newline \texttt{33c4c8c5762bcebdf8f544ab34a137a6}\newline\texttt{a7ba4065d08f5caee0a51188e1a67514} &
$72$ classes covering $128$ patterns; about $4.8\cdot10^8$ boxes (LB3); complete run $72/72$ at these claims:
\texttt{logs/cert\_lb3\_rerun.txt} of the same folder, completeness checked in \texttt{logs/k7\_verify\_pexact2.txt}
(the log of the refutation described in \S\ref{ssec:bnb} is not part of \texttt{supplementary/}). In Lean: displayed hypothesis.\\[4pt]
\texttt{CertMajV2} & \texttt{CertMajV2/\allowbreak majorant\_delta1\_beta.npy} ($60$ binary64 numbers)\newline \texttt{fd909f79669e2fc1ba5fd4e44d01d600}\newline\texttt{7ccf5303b7ce62d4c5ee08f7ef024845} &
Arb at $106$ bits. In Lean: \emph{kernel-checked}, $614$ cells, about $88$ s; the hypothesis is discharged.\\
\bottomrule
\end{tabular}
\end{table}

\subsection{Interval branch-and-bound}\label{ssec:bnb}
\paragraph{Reductions.} Since $k^2\ge0$, $F_{\bm}(g)\ge\sum_i\mu_ig_i$, so a box on which $\sum_i\mu_i(c_i-h_i)$ already exceeds the claim is
discarded, and only a bounded box needs to be covered. The weights and credits of the all-marks certificates are reversal-symmetric, so
$F_{\bm}(g)$ is unchanged when both the pattern and the gap vector are reversed; this reduces the $2^K$ patterns to $20$, $36$ and $72$ classes for
$K=5,6,7$. Coincident points ($g_i=0$) belong to the domain; there $k(0)=1$ makes $F$ large.

\paragraph{Kernel bounds.} For a window $\psi\ge0$, $\abs{k_\psi^{(j)}}\le(2\pi)^j\int v_\psi\abs s^j$; for $\tilde\psi$ this gives
$\abs{k'''}\le M_3=\frac{1747823}{7775424}\pi^3<6.969843$. These are proved analytically and enter as constants. Values of $k,k',k''$ at box centres are enclosed in Arb (or in
\texttt{mpmath.iv} with explicit Taylor remainders).

\paragraph{Lower bounds on a box} $\mathbf c\pm\mathbf h$. For a span $s$ (a run of consecutive gaps) let $\mathrm{sc}_s$, $\mathrm{hs}_s$ be the
sums of the $c_i$ and of the $h_i$ over the span.
\begin{itemize}
\item[(LB1)] On the span interval, $\abs k\ge\abs{k(\mathrm{sc})}-\abs{k'(\mathrm{sc})}\,\mathrm{hs}-\frac12\sup\abs{k''}\,\mathrm{hs}^2$; clip at $0$,
square, sum with the weights, and add $\sum\mu_i(c_i-h_i)$.
\item[(LB2)] $F(\mathbf c)-\sum_i\abs{\partial_iF(\mathbf c)}h_i-\frac12\sum_s\gamma_sW_{2,s}\,\mathrm{hs}_s^2$, where $W_{2,s}$ bounds
$\abs{w''}=2\abs{k'^2+kk''}$ on the span interval; valid because the Hessian of a span term is $w''\,\mathbf v_s\mathbf v_s^{\tsp}$.
\item[(LB3)] By Taylor's theorem with integral remainder,
$F(\mathbf c+\boldsymbol\delta)\ge F(\mathbf c)+\nabla F(\mathbf c)\cdot\boldsymbol\delta+\frac12\sum_s\gamma_s\,\underline w''_s\,(\mathbf v_s\cdot\boldsymbol\delta)^2$
on the box, with $\underline w''_s$ a lower bound for $w''$ on the span interval from a Taylor model of $k,k',k''$ with $\abs{k'''}\le M_3$. Spans with
$\underline w''_s\le0$ are bounded below by $-\frac12\gamma_s\abs{\underline w''_s}\mathrm{hs}_s^2$. The remaining quadratic $q$ is convex, so
$q(\boldsymbol\delta)\ge q(\hat{\mathbf d})+\nabla q(\hat{\mathbf d})\cdot(\boldsymbol\delta-\hat{\mathbf d})$ for \emph{every} $\hat{\mathbf d}$ (a
floating-point quadratic programme chooses $\hat{\mathbf d}$, which affects only tightness), and this affine function is minimised over the box exactly.
\end{itemize}
A box is accepted when a lower bound, evaluated in interval arithmetic, exceeds the claim; otherwise it is bisected. The certificate is the
statement that the stack empties with every box domain-pruned or accepted. LB3 cuts the trees by factors of $30$ to $200$. Tangent planes
after a convexity certificate are used in the verifiers of \mbox{sxuff}, \mbox{AMTOPA} and \mbox{npip99} (\S\ref{ssec:others}); LB3 is our form of this step, and we
have not compared it with theirs in substance.

\paragraph{Arithmetic and trust.} Box edges are exact dyadics. The default engine uses float64 intervals widened outward by one unit in the last
place at every operation, which assumes IEEE-754 round-to-nearest; a scalar engine in Arb balls avoids this assumption at about fifteen times the
cost. The correctness of Arb is trusted. The closing arithmetic ($H$, $\Phi_m$, the constants) was done in Arb from exact rationals and is also
checked in Lean.

\paragraph{Validation and negative controls.} Claims set just above the located minimum must fail, and do, with a rigorous counterexample:
for $K=5$, class $11111$ at the claim raised by $4.5\cdot10^{-5}$; for $K=6$, class $211112$ at $0.0166770$ ($F=0.01667680$); for $K=7$ the
certifier refuted class $1121212$ at the claims first proposed by the linear programme that chose the weights ($F=0.02054932<0.02054941$), in a
basin that the programme's search had missed, after which every claim was lowered by $10^{-5}$ and all $72$ classes certified; at the
lowered claim plus $10^{-5}$ the certifier refutes class $2112112$. For the polynomial window, the same engine at $K=7$ refutes the claim $0.00736$
($F=0.00735942$), and its certificate at the claim $0.00732$ was reproduced box for box independently ($18{,}215{,}893$ boxes). The
$K=8$ claim $0.00796$ lies $0.56\%$ below the floating-point minimum $0.0080050$. An independent kernel stress test of LB3 on $2400$ boxes found no
violation. The $K=8$ and $K=7$ certificates were each re-run in full: independent runs, not independent code.

\subsection{Kernel replay inside Lean}\label{ssec:replay}
\paragraph{The majorant (done).} The checker is written in core Lean. For each cell it proves one enclosure of $(\sin\theta,\cos\theta)$, $\theta=\pi m/30$
at the cell midpoint $m$, and generates the other $58$ cosines by a fixed-point complex recurrence $z_{k+1}=\lfloor z_kz_1/2^{64}\rfloor$ whose error
grows only linearly ($\abs{z_k-e^{ik\theta}}\le k\cdot2^{-55}$, proved); the sums $P(m)$, $P'(m)$, $P''(m)$ are exact integer sums, and a third-order
bound with $\abs{P'''}\le645$ covers the cell. $P\ge0$ on $[0,30]$ needs $600$ cells, $\cB\ge v_\psi$ on $[0,\frac12]$ needs $14$. The replay takes
about $88$ seconds in $13$ files, at most $1.64$ GB each; the smallest margins are $4.1\cdot10^{-6}$ and $4.4\cdot10^{-6}$. Negative controls fail as
they should, in a python mirror ($62$ cells) and in the kernel, which proves that the checker returns \texttt{false} on $[3.25,3.75]$ and on
$[\frac14,\frac12]$ for deliberately wrong data. The analytic part (the Fourier pair $\sinc^2\leftrightarrow$ triangle, the transform of $\cB$, its
support, $\int\cB=\beta_0$, and $2\beta_0+4\abs{\beta_{52}}<\frac{33}{10}$ in exact rationals) is proved in Lean. So the hypothesis \texttt{CertMajV2}
is discharged.

\paragraph{The $K=5$ all-marks certificate (replayed by the kernel).} The branch-and-bound trees above cannot be replayed box by box in the kernel. A second
certificate with large leaves was built for the same data: $92{,}074$ nodes and $46{,}047$ leaves over the $20$ classes, of three kinds --- cap
leaves (the penalty alone exceeds the claim), convex leaves (a tangent plane at a floating-point minimiser plus an exact positive semidefiniteness
check) and linear-programming leaves (five affine minorants per span with rational dual multipliers; $4{,}660$, $2{,}343$ and $39{,}044$ leaves
respectively). Every leaf is accepted by the exact rational mirror of the corrected checker (version~3, described below; $46{,}047$ leaves,
no failure). In a differential test on mutated leaves the mirror was never more permissive than the Lean checker; in any case soundness
rests on the kernel replay, not on the mirror. The checker stores each kernel enclosure (the values
$k,k',k''$ at a rational point, or bounds on $w,w''$ over a cell of width $\frac1{64}$) as a record proved once by kernel evaluation; leaves only
combine stored rationals. Measured in the replay of class $11111$ ($1{,}725$ leaves): about $0.24$ s per point record, $0.025$ s per cell and
$0.74$ s per leaf; with at most $40$ leaves per file the leaf files peak at $3.43$ GB in the one-class pilot. The full replay ($1{,}478$ files: $1{,}475$ shards and
three common files; $46{,}047$ leaves, $89{,}362$ point records; manifest sha256 \texttt{360d9b22469b9dbf449899a2c6d0bd5b}\allowbreak
\texttt{e696380e826f137ef20da71ae96c1609}) was estimated from the pilot at about $15.6$ Lean-hours; it took $20.4$. At the minimiser of class $11111$ the margin is $1.49\cdot10^{-5}$; raising
the claim by $1.5\cdot10^{-5}$, to $8.9\cdot10^{-8}$ above the true minimum, makes exactly one leaf fail in the kernel, the one that contains
the minimiser, as it should.
The soundness theorem of the checker is split into $28$ nodes. An earlier version of one leaf check was unsound: the linear-programming leaf
took a tangent line at a point computed from the floating-point minimiser without checking that this point lies in the span interval, where
alone the second-derivative penalty is valid; a kernel-checked counterexample (one span, a claim of $\frac9{10}$ accepted where the functional
is about $0.446$) shows it. The majorant certificate uses only the interval and enclosure layer of the checker, which is proved.
Version~3 of the checker corrects the defect with one guard, which admits the tangent line at the floating-point point only when that point
lies in the span interval; no other definition used by the checker changed, and the $20$ class certificates pass it unchanged (through the exact mirror).
The whole soundness chain of version~3 is proved with the three standard axioms and is in the library (Appendix~\ref{app:chain}):
\texttt{ZetaS.CertV2.\allowbreak ChainV3.\allowbreak certAM5\_of\_checks} derives \texttt{CertAM5} from the acceptance of the $20$ class certificates by the checker
and the identity of their data with the data of \texttt{CertAM5}.

The replay files use the library's own definitions. Built from an empty build folder, the one-class pilot compiles $56$ files with no
failure in $1{,}976$ s, at a peak of $3.49$ GB. The full replay, against the library at commit \texttt{bb7223b}, whose Lean files the release
changes only in comments (Appendix~\ref{app:layout}), compiled all $1{,}478$ files, with no failure and no retry, in $2$~h~$50$~min of
wall time and $20.4$ Lean-hours of processor time; the lowest free memory was $2.48$ GB and the build products take $7.1$ GB. It ran in two
pools, $273$ light files sixteen at once and $1{,}203$ heavy files six at once, and then the two top files alone ($21$~s and $18$~s, peak
$4.65$ and $4.67$ GB). The kernel checked the $46{,}047$ leaves and $89{,}362$ point records in exact rational arithmetic. The final file applies
\texttt{certAM5\_of\_checks} with the identity of the data by \texttt{rfl} and proves \texttt{ZetaS.CertV2.\allowbreak certAM5\_replayed};
the two theorems of the $K=5$ pair without hypothesis follow (Appendix~\ref{app:replay}). These three theorems
depend only on \texttt{propext}, \texttt{Classical.choice} and \texttt{Quot.sound} and rest on no open declaration (Appendix~\ref{app:census}). The $1{,}478$ source files
match their manifest file by file, and a scan of all of them finds no \texttt{sorry}, \texttt{native\_decide}, \texttt{axiom},
\texttt{unsafe}, \texttt{implemented\_by} or \texttt{extern}. The $1{,}478$ files, byte-identical to those checked, are the
library \texttt{ZetaSReplay} of the repository, and the two theorems are in its module \texttt{ReplayHeadlines}; \texttt{lake build
ZetaSReplay} repeats the replay (\S\ref{ssec:reproduce}).

\paragraph{The $K=7$ and $K=8$ certificates.} Large leaves cut the $K=7$ trees only by a factor of about four: an estimated $6\cdot10^7$ leaves and
$15$ GB of certificate, because the functional has a lattice of near-minimal basins whose number grows geometrically with $K$ (the branch-and-bound
trees grow by a factor of $20$ to $30$ per step in $K$). This is beyond the $32$ GB machine used. These certificates therefore remain displayed
hypotheses, identified by the sha256 of Table~\ref{tab:certs} and verified outside Lean as described in \S\ref{ssec:bnb}.

\subsection{How to reproduce}\label{ssec:reproduce}
\sloppy
{\raggedright Everything is in the repository \repourl, at tag \repotag{} (Appendix~\ref{app:layout}). The Lean libraries are at its root, the replay
files of \S\ref{ssec:replay} are the library \texttt{ZetaSReplay}, and the folder \texttt{supplementary/} holds the data of
Table~\ref{tab:certs}, the certifiers and their logs, and the generator of the replay files, with a guide (\texttt{supplementary/README.md}),
the commands that run them (\texttt{supplementary/VERIFY.md}) and the sha256 of every file
(\texttt{supplementary/CHECKSUMS.sha256}). Commands are run from the root of the repository.\par}
\begin{itemize}\raggedright
\item \emph{The statements and their axioms.} The toolchain is pinned in \texttt{lean-toolchain}, and
\texttt{lake exe cache get} fetches the prebuilt Mathlib. Then
\begin{Verbatim}[fontsize=\small,xleftmargin=1em]
lake build ZetaS
lake build Solution.FollowUpZeta
lake env lean comparator/PrintAxioms/FollowUpZeta.lean
\end{Verbatim}
Every line printed by the last command must read \texttt{\ldots{} depends on axioms: [propext, Classical.choice, Quot.sound]}. This checks
the fourteen statements of record (Appendix~\ref{app:record}), in which \texttt{CertAM5}, \texttt{CertAM7} and \texttt{CertS8} are
hypotheses. The files \texttt{comparator/README\_followup.md} and \texttt{comparator/config-followup.json} (the fourteen theorem names and
the permitted axioms) describe the same check with the comparator tool used by the repository of~\cite{AF}.
\item \emph{The kernel replay of \texttt{CertAM5}} (about $20$ Lean-hours, and several GB of memory per file, \S\ref{ssec:replay}):
\begin{Verbatim}[fontsize=\small,xleftmargin=1em]
lake build ZetaSReplay
lake env lean comparator/PrintAxioms/FollowUpReplay.lean
\end{Verbatim}
The last command prints the axioms of \texttt{ZetaS.Top.zeta\_simple\_K5\_uncond}, \texttt{ZetaS.Top.zeta\_distinct\_K5\_uncond} and
\texttt{ZetaS.CertV2.certAM5\_replayed}; again every line must read as above.
\item \emph{The local certificates outside Lean.} In \texttt{supplementary/certificates/}: \texttt{CertS8/} (the certifier
\texttt{bnb\_lb3.py} with \texttt{polywin.py}; the driver logs of the run and of the re-run, which end with \texttt{ALL SHARDS OK}),
\texttt{CertAM5/} (\texttt{bnb\_interval\_w.py}, the $K=5$ run, with its driver \texttt{runcert.sh}) and \texttt{CertAM7/}
(\texttt{bnb\_lb3\_w.py}, the LB3 engine with the all-marks check, with its driver \texttt{runlb3.sh}, and the exact constants
\texttt{pexact2.py}), each with its data and logs; and \texttt{common/}, with the input checks \texttt{verify\_inputs.py} (weights
nonnegative, span sums $2$, penalties summing to $\nu$, claims exactly $\sum_ib_i(m_i)$, all patterns covered) and \texttt{d819\_exact.py}.
\item \emph{The majorant:} \texttt{supplementary/certificates/CertMajV2/r4\_majorant\_arb.py} (Arb, $106$ bits); in Lean the files
\texttt{MajCheck.lean}, \texttt{MajSoundNPos.lean}, \texttt{MajSoundNMaj.lean} and the $13$ replay shards, all in \texttt{ZetaS/Majorant/}
and built with \texttt{ZetaS}.
\item \emph{The replay files from the certificate:} \texttt{supplementary/replay-K5/} holds the certificate with large leaves
(\texttt{v2\_K5\_*.json}, one file per class, and \texttt{v2common\_K5\_4.json}) and the generator \texttt{gen\_replay.py}, which writes the
replay files from it; \texttt{ZetaSReplay/MANIFEST.sha256} gives the sha256 of each replay file.
\end{itemize}

\section{Remarks on the limits of the method}\label{sec:remarks}

\paragraph{Ceilings.} Alp\"oge and Furman~\cite[\S7.2]{AF} call a \emph{bandwidth-one certificate} for $N_0^s/N$ a function of the
configuration that depends on it only through the first two trace moments against test functions of Fourier support in $[-1,1]$ and the
on-line/off-line partition, and is valid configuration by configuration, and they describe a ceiling of about $0.682$ for such certificates.
What is proved there is a quantitative bound with qualifications. For a constant $c_0$ and a weight $r$ on $[0,1]$ ($r\in C^1$, with $r'$
differentiable off a countable set and $r''$ integrable) such that $c_0+\sum_{j=1}^{256}(S_j/256)\,r(j/256)\le p$, where
$(S_j)$ is the sampled form factor of an explicit $256$-periodic law, the Lean theorem \texttt{Zeta23.\allowbreak PairCeiling.\allowbreak ceiling\_law256}
gives $c_0+\int_0^1r(x)\,x\,dx\le p+0.824\abs{r(1)}+2.55\cdot10^{-6}\big(\abs{r'(1)}+\int_0^1\abs{r''}\big)$, under the displayed
hypothesis \texttt{EnclOK} that $(S_j)$ lies in enclosures obtained outside Lean by interval arithmetic. For a certificate of this form that
is valid configuration by configuration, averaging over the law (a step done on paper) gives the hypothesis with $p=p_0\le0.6818287$, the
proportion of simple points of the law, and the bound is then below $0.6819$ if $r(1)=0$ and $\abs{r'(1)}+\int_0^1\abs{r''}<8.2$, in
particular for the two windows of~\cite[(2.6)]{AF}. By the $0.682$ ceiling we mean this bound; a bound of $0.682$ for every bandwidth-one
certificate is not proved in~\cite{AF} or here. Theorem~\ref{thm:S} lies below $0.682$. Within that class, the route of Section~\ref{sec:first}, which keeps only the energy among
\emph{simple} zeros, is capped much lower: periodic adversaries with gaps alternating near $1.03$ and $1.96$ carry an energy of about $0.0019$ per
simple zero, which caps cosine windows near $0.6737$ ($0.673666$ at $\cos1.48s$ with exact energies, $0.673697$ for the best cosine; Devine~\cite{Devine22066689}
states that with the fixed optimal Montgomery--Taylor kernel the pure Gram/rank machinery cannot force a bound above $0.6736$ [V$\sim$]; \mbox{trmdy}~\cite{trmdy} states on 11~August~2026 that certificates reading only pair
energies of simple zeros cannot exceed $0.674826$, and Zhongshan-Big-Jun~\cite{Zhongshan} give pressure-class ceilings near $0.6731$
[V$\sim$]), and the
same adversaries give at most $0.673776$ for the window $\tilde\psi$; the caps of the chain as $K$ grows level off near $0.67350$ for $\tilde\psi$.
These values are floating-point computations, not theorems; they explain why the bandwidth-one values of \S\ref{ssec:others} and
Theorem~\ref{thm:S} are all close together. (The repository teal-sea/zeta-lab states that the pressure-certificate family built on the Montgomery--Taylor window saturates at
$\sup\Phi_n\le0.675142509660254$ [V$\sim$]; we have not compared that family with the class of our computations.)

The route of Section~\ref{sec:simple} is not subject to the simple-only cap, because the binding adversaries of that cap use on-line double
zeros, whose energy it keeps. Its own limits, again in floating point: the linear programme that chooses the weights gives $0.6761153$ at
$K=7$ ($\mu=\frac1{500}$) and $0.6763110$ at $K=8$ ($\mu=\frac1{600}$), and the all-marks functional with a single window and exact energies is
capped at $0.681447$ (at $\psi=\cos2s$). A $K=8$ certificate was not attempted: the trees grow by a factor of $20$ to $30$ per step in $K$, an
estimated $1$ to $1.4\cdot10^{10}$ boxes over $136$ classes. Flat-topped windows lowered the $K=5$ value, so $\cos1.6s$ was kept. Whether an analogue of
the $0.682$ ceiling holds for bandwidth-one certificates of $N^s/N$ we have not determined.

\paragraph{What would be needed to go further.} For $N^s$ and $N^d$ the constants move with the local certificate. The $K=5$ certificate with large leaves was small
enough for the kernel to replay (\S\ref{ssec:replay}); larger $K$, better windows and a certificate at $K\ge6$ small enough for kernel
replay would need a better leaf (one whole basin per leaf) rather than more boxes. Beyond bandwidth one, or beyond
the first two trace moments, the inputs of~\cite{AF} are no longer enough.

For $N_0^s$ the route of Section~\ref{sec:simple} would give the same constants, $0.675158622$ with the $K=5$ certificate and $0.676102666$ with
$K=7$, if one inequality held for the zeros of $\zeta$: that the slack of the rank--trace inequality is at least
$a_1s_1+\frac{a_2}2(N-s_1)-\nu N-o(N)$. Lemma~\ref{lem:onslack} and Proposition~\ref{prop:agg} prove it whenever all zeros with $\Re\gamma\in I'$
lie on the line, so the obstruction is the off-line pairs near the line, which for $N_0^s$ bring no bonus (Remark~\ref{rem:notS}; \mbox{trmdy}~\cite{trmdy} named the pricing of off-line pairs as the gate
beyond the simple-only route on 11~August~2026 [V$\sim$]). We have
reduced this inequality to families of local block inequalities and proved several of them on explicit classes of configurations, but it is
not proved, and none of that work is part of this paper.


\appendix
{\rightskip=0pt plus 4em\relax % ragged right in Appendix A only (long Lean names)
% GENERATED by gen_appA_r4.py from appA_template_r4.tex (Lean text from zeta-23-lean-release, release-v1 at 8c15295801687f63bd3a8c5c96822800565303cf). Do not edit; edit the template and re-run.
\section{The Lean statements}\label{app:lean}
\sloppy

The Lean text below is copied verbatim, by a script that reads the files, from the release of the formalisation: the repository
\repourl{} at tag \repotag{} (\S\ref{app:layout}). Only line breaks inside long lines are added by the typesetting.

\emph{What is claimed} is stated in three files, which import nothing but Mathlib and one another: \texttt{comparator/ChallengeDeps.lean} of the Lean artifact
of~\cite{AF}, which defines the zeros of \texttt{riemannZeta} and the counting functions (\S\ref{app:counting}); the definition layer of this
paper, \texttt{comparator/ChallengeDeps/FollowUpZeta.lean}, with the windows, the local inequalities and the three certificate
\texttt{Prop}s (\S\ref{app:defs}); and the fourteen statements of record, \texttt{comparator/Challenge/FollowUpZeta.lean}
(\S\ref{app:record}). A reader who wants to know what is claimed needs only these three files. The statements of record carry
\texttt{sorry} by design: they are the challenge side of the comparator set-up of the repository of~\cite{AF}. The file
\texttt{comparator/Solution/FollowUpZeta.lean} proves exactly these fourteen statements from the library \texttt{ZetaS}, whose definitions
are copies of those of \S\ref{app:defs} in the namespace \texttt{ZetaS}; the conversions between the two are proofs, not trusted text. The
library's theorems are in \S\ref{app:headline}, the chain for the $K=5$ certificate in \S\ref{app:chain}, the replay's theorems in
\S\ref{app:replay} and the majorant in \S\ref{app:majorant}.

The doc-strings are part of the Lean text and are reproduced unchanged. They are development notes, not part of what is claimed, and
not all current: they refer to an earlier draft of this paper by its labels (for instance \texttt{thm:zeta-mainw2}
is Theorem~\ref{thm:S}, \texttt{thm:sigd-Sigma} and \texttt{thm:sigd-D} are Theorem~\ref{thm:main}, \texttt{cor:oll-SC} is
Corollary~\ref{cor:SC}, \texttt{lem:sigd-env} is Lemma~\ref{lem:room}) and to its line numbers, they name work packages, folders and
dates of the development, and ``level A'' in them means kernel-checked in the sense of \S\ref{ssec:formal}.
The data files that the doc-strings name are in \texttt{supplementary/certificates/}, with the same sha256 (Table~\ref{tab:certs}). The
header comment of \S4 of the definition layer says that the certificates were established in Arb ball arithmetic; as \S\ref{ssec:bnb} states,
the box bounds were computed in outward-rounded float64 intervals, with the kernel values enclosed in Arb. Only the Lean code is claimed.

\subsection{The counting functions}\label{app:counting}
From \texttt{comparator/ChallengeDeps.lean} of the repository of~\cite{AF}, unchanged from the tree as shipped (the functions used below are
\texttt{Ncount}, \texttt{Ndist}, \texttt{N0simple}, \texttt{Nsimple} and \texttt{zerosIn}). The phrase ``Not used by the challenge
statements'' in two of these doc-strings refers to the challenge statements of~\cite{AF}; \texttt{Nsimple} is used by those of
\S\ref{app:record}.

\begin{Verbatim}[fontsize=\footnotesize,breaklines,frame=single,framesep=2mm]
/-! ## 1. Nontrivial zeros of ζ and the counting functions (Theorems A–D) -/

/-- ρ is a nontrivial zero of the Riemann zeta function: ζ(ρ) = 0 with 0 < Re ρ < 1 (the open
critical strip). Every zero in the strip is "nontrivial" in the usual sense (it is neither a trivial
zero −2(n+1), of real part ≤ −2, nor the pole s = 1); the converse — that every zero other than the
trivial ones lies in the open strip — is classical and is not needed to state the theorems. -/
def IsNontrivialZero (ρ : ℂ) : Prop := riemannZeta ρ = 0 ∧ 0 < ρ.re ∧ ρ.re < 1

/-- m_ρ, the multiplicity of ρ: the order of vanishing of ζ at ρ, via Mathlib's `analyticOrderAt`
(ℕ∞-valued; `toNat` sends ⊤ to 0, so `1 ≤ zeroMult ρ` encodes both "ζ is not locally identically
zero at ρ" and "ρ is a zero"). -/
def zeroMult (ρ : ℂ) : ℕ := (analyticOrderAt riemannZeta ρ).toNat

/-- The nontrivial zeros with ordinate in the window (T₁, T₂]: {ρ | ζ(ρ) = 0, 0 < Re ρ < 1,
T₁ < Im ρ ≤ T₂}. (Positive-ordinate window, not |Im ρ|.) -/
def zerosIn (T₁ T₂ : ℝ) : Set ℂ := {ρ | IsNontrivialZero ρ ∧ T₁ < ρ.im ∧ ρ.im ≤ T₂}

/-- N(T₁,T₂): the number of nontrivial zeros with T₁ < Im ρ ≤ T₂, counted WITH multiplicity
(`∑ᶠ` is Mathlib's finite sum; that the window contains finitely many zeros is proved on the
solution side, not assumed). -/
def Ncount (T₁ T₂ : ℝ) : ℕ := ∑ᶠ ρ ∈ zerosIn T₁ T₂, zeroMult ρ

/-- N_d(T₁,T₂): the number of nontrivial zeros with T₁ < Im ρ ≤ T₂, each distinct zero counted
once. -/
def Ndist (T₁ T₂ : ℝ) : ℕ := (zerosIn T₁ T₂).ncard

/-- N₀(T₁,T₂): the zeros ON the critical line Re ρ = 1/2 with T₁ < Im ρ ≤ T₂, with multiplicity.
(Not used by the challenge statements; kept so that the block matches Zeta23/Statement.lean.) -/
def N0 (T₁ T₂ : ℝ) : ℕ := ∑ᶠ ρ ∈ zerosIn T₁ T₂ ∩ {ρ | ρ.re = 1 / 2}, zeroMult ρ

/-- N₀*(T₁,T₂): the number of DISTINCT zeros ON the critical line Re ρ = 1/2 with
T₁ < Im ρ ≤ T₂. This is what Theorems A and D bound from below. -/
def N0star (T₁ T₂ : ℝ) : ℕ := (zerosIn T₁ T₂ ∩ {ρ | ρ.re = 1 / 2}).ncard

/-- N₀ˢ(T₁,T₂): the number of SIMPLE zeros (multiplicity exactly 1) ON the critical line with
T₁ < Im ρ ≤ T₂. -/
def N0simple (T₁ T₂ : ℝ) : ℕ := (zerosIn T₁ T₂ ∩ {ρ | ρ.re = 1 / 2} ∩ {ρ | zeroMult ρ = 1}).ncard

/-- Nˢ(T₁,T₂): the number of simple zeros with T₁ < Im ρ ≤ T₂ (anywhere in the strip).
(Not used by the challenge statements; kept so that the block matches Zeta23/Statement.lean.) -/
def Nsimple (T₁ T₂ : ℝ) : ℕ := (zerosIn T₁ T₂ ∩ {ρ | zeroMult ρ = 1}).ncard
\end{Verbatim}

\subsection{The definitions of this paper}\label{app:defs}
The definition layer \texttt{comparator/ChallengeDeps/FollowUpZeta.lean} (namespace \texttt{FollowUpZeta}), from its first line of code to
its end; the file's header comment, which lists its contents, is omitted.

\begin{Verbatim}[fontsize=\footnotesize,breaklines,frame=single,framesep=2mm]
import ChallengeDeps

open scoped BigOperators
open Finset

noncomputable section

namespace FollowUpZeta

/-! ## §1 Zeros that are simple or on the critical line -/

/-- N^sc(T₁,T₂) (cor:oll-SC): the nontrivial zeros with T₁ < Im ρ ≤ T₂ that lie ON the critical line, counted WITH
multiplicity, plus those OFF the line that are simple (both members ρ, 1 − ρ̄ of each simple off-line pair; multiple
off-line zeros are not counted). An off-line zero in the set is simple, so its weight `zeroMult ρ` is 1. -/
def Nsc (T₁ T₂ : ℝ) : ℕ :=
  ∑ᶠ ρ ∈ zerosIn T₁ T₂ ∩ ({ρ | ρ.re = 1 / 2} ∪ {ρ | zeroMult ρ = 1}), zeroMult ρ

/-! ## §2 Windows and the normalised kernel -/

/-- The polynomial window ψ̃ of thm:zeta-mainw2:
ψ̃(s) = (50000/50037)(1 + (127/2500)u − (8997/10000)u² + (6187/10000)u³ + (3/1250)u⁴), u = (2s)²,
on [−1/2, 1/2] (values outside are never read). -/
def psiPoly8A (s : ℝ) : ℝ :=
  50000 / 50037 * (1 + 127 / 2500 * (2 * s) ^ 2 - 8997 / 10000 * ((2 * s) ^ 2) ^ 2
    + 6187 / 10000 * ((2 * s) ^ 2) ^ 3 + 3 / 1250 * ((2 * s) ^ 2) ^ 4)

/-- The window ψ(s) = cos(1.6 s) of thm:sigd-Sigma and thm:sigd-D. -/
def psiCos16 (s : ℝ) : ℝ := Real.cos (8 / 5 * s)

/-- The normalised Fourier transform of a window: k_ψ(x) = ∫_{−1/2}^{1/2} ψ(s) cos(2πxs) ds / ∫_{−1/2}^{1/2} ψ. -/
def kPsi (ψ : ℝ → ℝ) (x : ℝ) : ℝ :=
  (∫ s in (-(1 / 2 : ℝ))..(1 / 2), ψ s * Real.cos (2 * Real.pi * x * s))
    / ∫ s in (-(1 / 2 : ℝ))..(1 / 2), ψ s

/-! ## §3 The local inequalities -/

/-- Position weights for K consecutive points (lem:zeta-Dprime), 0-based: pair weights `γ s i` for spans
1 ≤ s ≤ K − 1 and positions 0 ≤ i < K − s, nonnegative and summing to 2 for every span; gap penalties `μ l ≥ 0`,
l : Fin (K − 1). Values of `γ` outside the index range are never read. -/
structure LocalWeights (K : ℕ) where
  γ : ℕ → ℕ → ℝ
  μ : Fin (K - 1) → ℝ
  two_le : 2 ≤ K
  γ_nonneg : ∀ s ∈ Icc 1 (K - 1), ∀ i ∈ range (K - s), 0 ≤ γ s i
  γ_sum : ∀ s ∈ Icc 1 (K - 1), ∑ i ∈ range (K - s), γ s i = 2
  μ_nonneg : ∀ l, 0 ≤ μ l

/-- ν := Σ_i μ_i. -/
def LocalWeights.nu {K : ℕ} (W : LocalWeights K) : ℝ := ∑ l, W.μ l

/-- The span g_i + ⋯ + g_{i+s−1} of a gap vector (0-based; indices ≥ K − 1 contribute nothing). -/
def gapSpan {K : ℕ} (g : Fin (K - 1) → ℝ) (i s : ℕ) : ℝ :=
  ∑ l ∈ univ.filter (fun l : Fin (K - 1) => i ≤ l.val ∧ l.val < i + s), g l

/-- F_{γ,μ}(g) = Σ_i μ_i g_i + Σ_{s=1}^{K−1} Σ_i γ_{s,i} k(g_i + ⋯ + g_{i+s−1})² (eq:zeta-LIw). -/
def localF (k : ℝ → ℝ) {K : ℕ} (W : LocalWeights K) (g : Fin (K - 1) → ℝ) : ℝ :=
  (∑ l, W.μ l * g l)
    + ∑ s ∈ Icc 1 (K - 1), ∑ i ∈ range (K - s), W.γ s i * k (gapSpan g i s) ^ 2

/-- (LI′) of lem:zeta-Dprime: F_{γ,μ}(g) ≥ c for every gap vector g ∈ [0,∞)^{K−1}. -/
def LocalCert (k : ℝ → ℝ) {K : ℕ} (W : LocalWeights K) (c : ℝ) : Prop :=
  ∀ g : Fin (K - 1) → ℝ, (∀ l, 0 ≤ g l) → c ≤ localF k W g

/-- All-marks data (def:zeta-LIm): position weights as in `LocalWeights` and site credits `b i j` for site i : Fin K and
capped mark j : Fin 2 (j = 0 ↦ mark 1, j = 1 ↦ mark 2); no sign condition. -/
structure MarkWeights (K : ℕ) extends LocalWeights K where
  b : Fin K → Fin 2 → ℝ

/-- a_j := Σ_i b_i(j); `a 0 = a₁`, `a 1 = a₂`. -/
def MarkWeights.a {K : ℕ} (W : MarkWeights K) (j : Fin 2) : ℝ := ∑ i, W.b i j

/-- The numerical mark m_i ∈ {1, 2} at 0-based position i of a pattern (0 outside [0, K)). -/
def markVal {K : ℕ} (m : Fin K → Fin 2) (i : ℕ) : ℝ :=
  if h : i < K then ((m ⟨i, h⟩ : ℕ) : ℝ) + 1 else 0

/-- F_𝐦(g) = Σ_{s,i} γ_{s,i} m_i m_{i+s} k(g_i + ⋯ + g_{i+s−1})² + Σ_i μ_i g_i (def:zeta-LIm). -/
def localFm (k : ℝ → ℝ) {K : ℕ} (W : MarkWeights K) (m : Fin K → Fin 2) (g : Fin (K - 1) → ℝ) : ℝ :=
  (∑ s ∈ Icc 1 (K - 1), ∑ i ∈ range (K - s),
      W.γ s i * markVal m i * markVal m (i + s) * k (gapSpan g i s) ^ 2)
    + ∑ l, W.μ l * g l

/-- (LI_m) of def:zeta-LIm: for EVERY mark pattern 𝐦 ∈ {1,2}^K and every g ∈ [0,∞)^{K−1}, F_𝐦(g) ≥ Σ_i b_i(m_i). -/
def LocalCertAM (k : ℝ → ℝ) {K : ℕ} (W : MarkWeights K) : Prop :=
  ∀ m : Fin K → Fin 2, ∀ g : Fin (K - 1) → ℝ, (∀ l, 0 ≤ g l) → ∑ i, W.b i (m i) ≤ localFm k W m g

/-! ## §4 The certificate hypotheses

Each Prop asserts that SOME weights with the stated constants satisfy the local inequality on the whole orthant. The
weights that realise it are exact rationals in the data files named below (identified by sha256); the inequality for those
weights was established outside Lean by interval branch-and-bound (Arb ball arithmetic, every box accepted only when a
rigorous lower bound exceeds the claim; the unbounded part of the orthant is discarded exactly, since
F ≥ (min μ)·Σ g there). Paths are relative to the development's project folder; the files named here are in
`supplementary/certificates/` of the repository, some logs under new names (map: `supplementary/README.md`). -/

/-- **Certificate S8** (thm:zeta-mainw2, rem:zeta-P8, lem:zeta-Dprime): K = 8, window ψ̃, weights with ν = 7/1700 for
which (LI′) holds at the claim c = 0.00796.
Data: `round2/d6_5_numerics/w_poly8A_K8_mu1700.json`, sha256
`257b340d9da4f343600a21589f16d8915f3e082499d475e153c5453980671e78`; window `round2/d6_5_numerics/win_poly8A.json`,
sha256 `c4005f57a7920f78aa2df33b093c9c5030efac2b8e39fc6cf7371b0c0f32dc14` (the unnormalised v = (50037/50000)ψ̃; the kernel
is homogeneous of degree 0). Certifier `round2/d6_5_numerics/bnb_lb3.py` with `polywin.py`; runs
`round2/d6_5_numerics/runs/P8/DRIVER.log` and the independent re-run `runs/P8b/DRIVER.log`, both ending
`ALL SHARDS OK: certificate claim=0.00796 holds` (about 4.5·10⁸ boxes each). -/
def CertS8 : Prop :=
  ∃ W : LocalWeights 8, W.nu = 7 / 1700 ∧ LocalCert (kPsi psiPoly8A) W (199 / 25000)

/-- **Certificate AM5** (thm:zeta-allmarks, rem:sigd-cert, the K = 5 row): all-marks data with a₁ = 1280197/10⁸,
a₂ = 48749/3125000, ν = 1/125 for which (LI_m) holds for k_ψ, ψ = cos 1.6s.
Data: `round1/X2_numerics/claims_a1.6_K5_mu500.json` (byte-identical to `claims_d15e6_a1.6_K5_mu500.json`), sha256
`cf055dd2ee9a9e469b6510fe70b5fbc416b8a40515cc6746bc390d40d88777f2`; the 20 class files `xpat_K5_*.json` of the same
folder, sha256 of their concatenation in sorted file-name order
`a5947c01d34ddcbc7366861b35e83942099fcbf7ac71e13fcd432796882dd34e`. Certifier `round1/X2_numerics/bnb_interval_w.py`;
logs `cert_K5_w_a16_mu500_d15e6.txt` and `cert_K5_w_a16_mu500_d15e6_b.txt` (20/20 reversal classes ok, covering the 32
patterns; the weights are reversal-symmetric). A replay of this certificate by the Lean kernel was completed on
28 September 2026: `ZetaS.CertV2.certAM5_replayed` proves the library's copy `ZetaS.CertAM5`, and the hypothesis-free
K = 5 theorems are in the module `ReplayHeadlines` of the library `ZetaSReplay`. In the statements of this topic the
Prop remains a displayed hypothesis, like the other two. -/
def CertAM5 : Prop :=
  ∃ W : MarkWeights 5, W.a 0 = 1280197 / 10 ^ 8 ∧ W.a 1 = 48749 / 3125000 ∧ W.nu = 1 / 125 ∧
    LocalCertAM (kPsi psiCos16) W

/-- **Certificate AM7** (thm:sigd-Sigma, thm:sigd-D, rem:sigd-cert, the K = 7 row): all-marks data with
a₁ = 1824837/10⁸, a₂ = 1168069/(5·10⁷), ν = 3/250 for which (LI_m) holds for k_ψ, ψ = cos 1.6s.
Data: `round2/d8_9_numerics/k7m500/rerun/claims_a1.6_K7_mu500_L1e5.json` (the final claims, every claim of the LP lowered by
exactly 10⁻⁵), sha256 `d2cf393019407bdc797c73147ec8e3dc66ec41fe745600c84173f74ac66bfc5c`; the 72 class files
`xpat_K7_*.json` of the same folder, sha256 of their concatenation in sorted file-name order
`33c4c8c5762bcebdf8f544ab34a137a6a7ba4065d08f5caee0a51188e1a67514`. Certifier `round2/d8_9_numerics/bnb_lb3_w.py`
(driver `runlb3.sh`); certificate log `round2/d8_9_numerics/k7m500/rerun/cert_lb3_rerun.txt` (72/72 reversal classes ok,
each at exactly the final claim); completeness (the 72 classes cover all 128 patterns, each certified at a claim ≥ the
required one) `round2/d8_21_numerics/k7_verify_pexact2.txt`. The weights are reversal-symmetric. -/
def CertAM7 : Prop :=
  ∃ W : MarkWeights 7, W.a 0 = 1824837 / 10 ^ 8 ∧ W.a 1 = 1168069 / (5 * 10 ^ 7) ∧ W.nu = 3 / 250 ∧
    LocalCertAM (kPsi psiCos16) W

end FollowUpZeta
\end{Verbatim}

\subsection{The fourteen statements of record}\label{app:record}
The statement file \texttt{comparator/Challenge/FollowUpZeta.lean}, from its first line of code to its end (the header comment is omitted).
The six headline statements are those of Theorems~\ref{thm:main} and~\ref{thm:S} and Corollary~\ref{cor:SC}, each with its single
certificate hypothesis; the six \texttt{\_cumulative} forms are the same bounds on the windows $0<\Im\rho\le T$; the last two statements
show that the counts are genuine (every window holds finitely many zeros) and that the denominator tends to infinity.

\begin{Verbatim}[fontsize=\footnotesize,breaklines,frame=single,framesep=2mm]
import ChallengeDeps.FollowUpZeta

open FollowUpZeta

noncomputable section

/-- **Simple zeros (K = 5 certificate)**, thm:sigd-Sigma with the K = 5 row of rem:sigd-cert: at least 0.675158622 of
the nontrivial zeros with T < Im ρ ≤ 2T (counted with multiplicity) are simple, for all large T. -/
theorem zeta_simple_K5 (hcert : CertAM5) :
    ∀ ε > 0, ∃ T₀ : ℝ, ∀ T ≥ T₀,
      ((675158622 : ℝ) / 10 ^ 9 - ε) * (Ncount T (2 * T) : ℝ) ≤ Nsimple T (2 * T) := by
  sorry

/-- `zeta_simple_K5` on the cumulative windows 0 < Im ρ ≤ T (same constant). -/
theorem zeta_simple_K5_cumulative (hcert : CertAM5) :
    ∀ ε > 0, ∃ T₀ : ℝ, ∀ T ≥ T₀,
      ((675158622 : ℝ) / 10 ^ 9 - ε) * (Ncount 0 T : ℝ) ≤ Nsimple 0 T := by
  sorry

/-- **Distinct zeros (K = 5 certificate)**, thm:sigd-D with the K = 5 row of rem:sigd-cert: at least 0.837579311. -/
theorem zeta_distinct_K5 (hcert : CertAM5) :
    ∀ ε > 0, ∃ T₀ : ℝ, ∀ T ≥ T₀,
      ((837579311 : ℝ) / 10 ^ 9 - ε) * (Ncount T (2 * T) : ℝ) ≤ Ndist T (2 * T) := by
  sorry

/-- `zeta_distinct_K5` on the cumulative windows 0 < Im ρ ≤ T (same constant). -/
theorem zeta_distinct_K5_cumulative (hcert : CertAM5) :
    ∀ ε > 0, ∃ T₀ : ℝ, ∀ T ≥ T₀,
      ((837579311 : ℝ) / 10 ^ 9 - ε) * (Ncount 0 T : ℝ) ≤ Ndist 0 T := by
  sorry

/-- **Simple zeros**, thm:sigd-Sigma (K = 7 certificate): at least 0.676102666. -/
theorem zeta_simple_K7 (hcert : CertAM7) :
    ∀ ε > 0, ∃ T₀ : ℝ, ∀ T ≥ T₀,
      ((676102666 : ℝ) / 10 ^ 9 - ε) * (Ncount T (2 * T) : ℝ) ≤ Nsimple T (2 * T) := by
  sorry

/-- `zeta_simple_K7` on the cumulative windows 0 < Im ρ ≤ T (same constant). -/
theorem zeta_simple_K7_cumulative (hcert : CertAM7) :
    ∀ ε > 0, ∃ T₀ : ℝ, ∀ T ≥ T₀,
      ((676102666 : ℝ) / 10 ^ 9 - ε) * (Ncount 0 T : ℝ) ≤ Nsimple 0 T := by
  sorry

/-- **Distinct zeros**, thm:sigd-D (K = 7 certificate): at least 0.838051333. -/
theorem zeta_distinct_K7 (hcert : CertAM7) :
    ∀ ε > 0, ∃ T₀ : ℝ, ∀ T ≥ T₀,
      ((838051333 : ℝ) / 10 ^ 9 - ε) * (Ncount T (2 * T) : ℝ) ≤ Ndist T (2 * T) := by
  sorry

/-- `zeta_distinct_K7` on the cumulative windows 0 < Im ρ ≤ T (same constant). -/
theorem zeta_distinct_K7_cumulative (hcert : CertAM7) :
    ∀ ε > 0, ∃ T₀ : ℝ, ∀ T ≥ T₀,
      ((838051333 : ℝ) / 10 ^ 9 - ε) * (Ncount 0 T : ℝ) ≤ Ndist 0 T := by
  sorry

/-- **Simple zeros on the critical line**, thm:zeta-mainw2 (window ψ̃, K = 8): at least 0.6733736895. -/
theorem zeta_simple_on_line (hcert : CertS8) :
    ∀ ε > 0, ∃ T₀ : ℝ, ∀ T ≥ T₀,
      ((6733736895 : ℝ) / 10 ^ 10 - ε) * (Ncount T (2 * T) : ℝ) ≤ N0simple T (2 * T) := by
  sorry

/-- `zeta_simple_on_line` on the cumulative windows 0 < Im ρ ≤ T (same constant). -/
theorem zeta_simple_on_line_cumulative (hcert : CertS8) :
    ∀ ε > 0, ∃ T₀ : ℝ, ∀ T ≥ T₀,
      ((6733736895 : ℝ) / 10 ^ 10 - ε) * (Ncount 0 T : ℝ) ≤ N0simple 0 T := by
  sorry

/-- **Zeros that are simple or on the critical line**, cor:oll-SC (on-line zeros with multiplicity plus simple
off-line zeros): at least 0.8879195. -/
theorem zeta_simple_or_critical (hcert : CertS8) :
    ∀ ε > 0, ∃ T₀ : ℝ, ∀ T ≥ T₀,
      ((8879195 : ℝ) / 10 ^ 7 - ε) * (Ncount T (2 * T) : ℝ) ≤ Nsc T (2 * T) := by
  sorry

/-- `zeta_simple_or_critical` on the cumulative windows 0 < Im ρ ≤ T (same constant). -/
theorem zeta_simple_or_critical_cumulative (hcert : CertS8) :
    ∀ ε > 0, ∃ T₀ : ℝ, ∀ T ≥ T₀,
      ((8879195 : ℝ) / 10 ^ 7 - ε) * (Ncount 0 T : ℝ) ≤ Nsc 0 T := by
  sorry

/-- Non-vacuity companion: every window T₁ < Im ρ ≤ T₂ holds finitely many nontrivial zeros, so `Ncount`, `Nsc`
(finite sums) and `Nsimple`, `Ndist`, `N0simple` (`Set.ncard`) are genuine counts (Lean's `finsum` and `ncard` would
return 0 on an infinite set). -/
theorem zeta_zerosIn_finite (T₁ T₂ : ℝ) :
    (zerosIn T₁ T₂).Finite := by
  sorry

/-- Non-vacuity companion: the denominator N(T, 2T) tends to infinity. -/
theorem zeta_tendsto_Ncount :
    Filter.Tendsto (fun T : ℝ => (Ncount T (2 * T) : ℝ)) Filter.atTop Filter.atTop := by
  sorry

end
\end{Verbatim}

\subsection{The library's headline theorems}\label{app:headline}
The four certificate-only theorems for simple and distinct zeros (\texttt{ZetaS/Top/TopFinal.lean}). They apply the forms with the majorant
threshold $\frac{33}{10}$ to the proved majorant certificate \texttt{MajSound.certMajV2\_proved} (\S\ref{app:majorant}), so that their only
hypothesis is the local certificate:

\begin{Verbatim}[fontsize=\footnotesize,breaklines,frame=single,framesep=2mm]
namespace ZetaS
namespace Top

/-- `zeta_simple_K5_V2` with the majorant certificate proved (`MajSound.certMajV2_proved`). -/
theorem zeta_simple_K5_final (hcert : CertAM5) :
    ∀ ε > 0, ∃ T₀ : ℝ, ∀ T ≥ T₀,
      ((675158622 : ℝ) / 10 ^ 9 - ε) * (Ncount T (2 * T) : ℝ) ≤ Nsimple T (2 * T) :=
  zeta_simple_K5_V2 hcert MajSound.certMajV2_proved

/-- `zeta_distinct_K5_V2` with the majorant certificate proved (`MajSound.certMajV2_proved`). -/
theorem zeta_distinct_K5_final (hcert : CertAM5) :
    ∀ ε > 0, ∃ T₀ : ℝ, ∀ T ≥ T₀,
      ((837579311 : ℝ) / 10 ^ 9 - ε) * (Ncount T (2 * T) : ℝ) ≤ Ndist T (2 * T) :=
  zeta_distinct_K5_V2 hcert MajSound.certMajV2_proved

/-- `zeta_simple_K7_V2` with the majorant certificate proved (`MajSound.certMajV2_proved`). -/
theorem zeta_simple_K7_final (hcert : CertAM7) :
    ∀ ε > 0, ∃ T₀ : ℝ, ∀ T ≥ T₀,
      ((676102666 : ℝ) / 10 ^ 9 - ε) * (Ncount T (2 * T) : ℝ) ≤ Nsimple T (2 * T) :=
  zeta_simple_K7_V2 hcert MajSound.certMajV2_proved

/-- `zeta_distinct_K7_V2` with the majorant certificate proved (`MajSound.certMajV2_proved`). -/
theorem zeta_distinct_K7_final (hcert : CertAM7) :
    ∀ ε > 0, ∃ T₀ : ℝ, ∀ T ≥ T₀,
      ((838051333 : ℝ) / 10 ^ 9 - ε) * (Ncount T (2 * T) : ℝ) ≤ Ndist T (2 * T) :=
  zeta_distinct_K7_V2 hcert MajSound.certMajV2_proved

end Top
end ZetaS
\end{Verbatim}

The two theorems for simple zeros on the line and for zeros that are simple or on the line (\texttt{ZetaS/Top/TopSSC.lean}, namespace \texttt{ZetaS.Top});
statements only, each followed in the file by a tactic proof of about ten lines:

\begin{Verbatim}[fontsize=\footnotesize,breaklines,frame=single,framesep=2mm]
/-- **S — simple zeros on the critical line** (the architect's `zeta_simple_on_line`, thm:zeta-mainw2). -/
theorem zeta_simple_on_line (hcert : CertS8) :
    ∀ ε > 0, ∃ T₀ : ℝ, ∀ T ≥ T₀,
      ((6733736895 : ℝ) / 10 ^ 10 - ε) * (Ncount T (2 * T) : ℝ) ≤ N0simple T (2 * T) := by
\end{Verbatim}

\begin{Verbatim}[fontsize=\footnotesize,breaklines,frame=single,framesep=2mm]
/-- **SC — simple or critical zeros** (the architect's `zeta_simple_or_critical`, cor:oll-SC). -/
theorem zeta_simple_or_critical (hcert : CertS8) :
    ∀ ε > 0, ∃ T₀ : ℝ, ∀ T ≥ T₀,
      ((8879195 : ℝ) / 10 ^ 7 - ε) * (Ncount T (2 * T) : ℝ) ≤ Nsc T (2 * T) := by
\end{Verbatim}

The statements of these six theorems are, character for character, those of the six dyadic statements of record in \S\ref{app:record},
with the definitions of the namespace \texttt{ZetaS} in place of those of \texttt{FollowUpZeta}; \texttt{comparator/Solution/FollowUpZeta.lean}
derives the fourteen statements of record from them. The library's definition of \texttt{CertAM5} (\texttt{ZetaS/Top/TopDefs.lean}) is:

\begin{Verbatim}[fontsize=\footnotesize,breaklines,frame=single,framesep=2mm]
/-- **Certificate AM5**: all-marks data with `a₁ = 1280197/10⁸`, `a₂ = 48749/3125000`, `ν = 1/125` for which (LI_m)
holds for `k_ψ`, `ψ = cos 1.6s` (X2's `round1/X2_numerics/xpat_K5_*.json`, `claims_d15e6_a1.6_K5_mu500.json`). -/
def CertAM5 : Prop :=
  ∃ W : MarkWeights 5, W.a 0 = 1280197 / 10 ^ 8 ∧ W.a 1 = 48749 / 3125000 ∧ W.nu = 1 / 125 ∧
    LocalCertAM (kPsi psiCos16) W
\end{Verbatim}

\subsection{The chain for the \texorpdfstring{$K=5$}{K=5} certificate}\label{app:chain}
From \texttt{ZetaS/Cert/ChainV3.lean}: the soundness theorem of the corrected checker and the theorem that twenty checked certificates, one
per canonical class, give \texttt{CertAM5} (the file ends with a diagnostic \texttt{\#eval}, omitted):

\begin{Verbatim}[fontsize=\footnotesize,breaklines,frame=single,framesep=2mm]
namespace ZetaS.CertV2.ChainV3

open ZetaS.CertV2

/-- **The certificate theorem** (node C21, L5_1's proof), against the V3 checker, all dependencies proved. -/
theorem cert_check_sound (C : Cert) (h : C.check = true) : C.data.Holds := Cert.check_sound' C h

/-- **Twenty checked certificates, one per canonical class, give `CertAM5`** (V3 checker, level A). -/
theorem certAM5_of_checks (C : List Cert) (hdata : C.map Cert.data = classes5.map classData)
    (hchk : ∀ c ∈ C, c.check = true) : CertAM5 :=
  certAM5_of_checks_of_sound cert_check_sound C hdata hchk

end ZetaS.CertV2.ChainV3
\end{Verbatim}

\subsection{The replay's theorems: the \texorpdfstring{$K=5$}{K=5} pair without hypothesis}\label{app:replay}
These are in the library \texttt{ZetaSReplay} of the repository, which holds the $1{,}478$ files of the replay (\S\ref{ssec:replay}) in one
flat folder, byte-identical to the files that the kernel checked, and is not among the default build targets; \texttt{lake build
ZetaSReplay} repeats the replay (\S\ref{ssec:reproduce}). The replay compiled all $1{,}478$ files with no failure, in $2$~h~$50$~min of
wall time and $20.4$ Lean-hours, against the library at commit
\texttt{bb7223b}, whose Lean files the release changes only in comments (\S\ref{app:layout}). From \texttt{ZetaSReplay/ReplayAM5.lean}, the
statement of \texttt{certAM5\_replayed} and the first line of its proof (the proof term continues with the twenty theorems
\texttt{R<class>.cert\_ok} proved by the kernel; omitted):

\begin{Verbatim}[fontsize=\footnotesize,breaklines,frame=single,framesep=2mm]
namespace ZetaS.CertV2

theorem certAM5_replayed : ZetaS.CertAM5 :=
  ChainV3.certAM5_of_checks [R11111.cert, R11112.cert, R11121.cert, R11122.cert, R11211.cert, R11212.cert, R11221.cert, R11222.cert, R12112.cert, R12121.cert, R12122.cert, R12212.cert, R12221.cert, R12222.cert, R21112.cert, R21122.cert, R21212.cert, R21222.cert, R22122.cert, R22222.cert] rfl
\end{Verbatim}

From \texttt{ZetaSReplay/ReplayHeadlines.lean}, the two theorems of the $K=5$ pair without hypothesis; their statements are those of
\texttt{ZetaS.Top.zeta\_simple\_K5\_final} and \texttt{zeta\_distinct\_K5\_final} (\S\ref{app:headline}) without the hypothesis
\texttt{hcert}:

\begin{Verbatim}[fontsize=\footnotesize,breaklines,frame=single,framesep=2mm]
namespace ZetaS
namespace Top

/-- **Simple zeros, K = 5, hypothesis-free** (thm:sigd-Sigma with the K = 5 certificate). -/
theorem zeta_simple_K5_uncond :
    ∀ ε > 0, ∃ T₀ : ℝ, ∀ T ≥ T₀,
      ((675158622 : ℝ) / 10 ^ 9 - ε) * (Ncount T (2 * T) : ℝ) ≤ Nsimple T (2 * T) :=
  zeta_simple_K5_final ZetaS.CertV2.certAM5_replayed

/-- **Distinct zeros, K = 5, hypothesis-free** (thm:sigd-D with the K = 5 certificate). -/
theorem zeta_distinct_K5_uncond :
    ∀ ε > 0, ∃ T₀ : ℝ, ∀ T ≥ T₀,
      ((837579311 : ℝ) / 10 ^ 9 - ε) * (Ncount T (2 * T) : ℝ) ≤ Ndist T (2 * T) :=
  zeta_distinct_K5_final ZetaS.CertV2.certAM5_replayed

end Top
end ZetaS
\end{Verbatim}

\subsection{The majorant}\label{app:majorant}
The majorant is not a hypothesis of any headline theorem: it is proved. Its \texttt{Prop} (\texttt{ZetaS/ChallengeZetaS.lean}, namespace
\texttt{ZetaS}), the instance used (\texttt{ZetaS/Majorant/CertMajV2.lean}) and the theorem (\texttt{ZetaS/Majorant/CertMajV2Proof.lean},
namespace \texttt{ZetaS.MajSound}):

\begin{Verbatim}[fontsize=\footnotesize,breaklines,frame=single,framesep=2mm]
/-- **`MajorantCert ψ λ`** — the certified input of lem:sigd-env (l.1283–1296), in the form the proof uses:
an even integrable `B ≥ 0` on `ℝ` with `B ≥ v_ψ := ψ/∫ψ` on `[−1/2,1/2]`, whose Fourier transform
`B̂(t) = ∫ B(s) cos(2πts) ds` vanishes for `|t| ≥ 1`, and with `2 B̂(0) + 4 sup_{[3/4,1]} |B̂| < λ`.
The draft's instance: `B̂` piecewise linear with nodes `j/60`, `β₀ = 1.5556524027186915`, `β* = |β₅₂|`,
and `λ = λ_c = 2 + √(2 − 2a₁)`. -/
def MajorantCert (ψ : ℝ → ℝ) (lam : ℝ) : Prop :=
  ∃ B : ℝ → ℝ, MeasureTheory.Integrable B ∧ (∀ s, B (-s) = B s) ∧ (∀ s, 0 ≤ B s) ∧
    (∀ s ∈ Set.Icc (-(1 / 2 : ℝ)) (1 / 2), ψ s / (∫ t in (-(1 / 2 : ℝ))..(1 / 2), ψ t) ≤ B s) ∧
    (∀ t : ℝ, 1 ≤ |t| → ∫ s, B s * Real.cos (2 * Real.pi * t * s) = 0) ∧
    ∃ βstar : ℝ, (∀ t ∈ Set.Icc (3 / 4 : ℝ) 1, |∫ s, B s * Real.cos (2 * Real.pi * t * s)| ≤ βstar) ∧
      2 * (∫ s, B s) + 4 * βstar < lam
\end{Verbatim}

\begin{Verbatim}[fontsize=\footnotesize,breaklines,frame=single,framesep=2mm]
/-- **CertMajV2** — `MajorantCert psiCos16 (33/10)`. -/
def CertMajV2 : Prop := MajorantCert psiCos16 (33 / 10)
\end{Verbatim}

\begin{Verbatim}[fontsize=\footnotesize,breaklines,frame=single,framesep=2mm]
/-- **CertMajV2**, instantiated with the C05/C06 node statements. -/
theorem certMajV2_proved : CertMajV2 := certMajV2_of_CNodes cnodes_of_statements
\end{Verbatim}

\texttt{ZetaS/Majorant/CertMajProof.lean} also proves, as \texttt{MajSound.certMaj\_proved}, the first version \texttt{CertMaj} of
\texttt{ZetaS/ChallengeZetaS.lean}, with the threshold $3.2063638853$; no theorem uses it.

\subsection{Axiom census and \texttt{sorry}}\label{app:census}
At commit \texttt{bb7223b}, \texttt{\#print axioms} gives exactly \texttt{[propext, Classical.choice, Quot.sound]} for
each of the six theorems of \S\ref{app:headline}, for \texttt{ZetaS.CertV2.ChainV3.certAM5\_of\_checks} and
\texttt{ZetaS.CertV2.ChainV3.cert\_check\_sound}, and for all fourteen theorems of
\texttt{comparator/Solution/FollowUpZeta.lean}. For the replay's three theorems of \S\ref{app:replay} it gives the same three axioms, and
no open declaration lies beneath them. In the repository,
\texttt{comparator/PrintAxioms/FollowUpZeta.lean} prints the axioms of the fourteen theorems of record and
\texttt{comparator/PrintAxioms/FollowUpReplay.lean} those of the three theorems of \S\ref{app:replay} (\S\ref{ssec:reproduce}); the logs and
the \texttt{\#print axioms} outputs of the builds of the release are in \texttt{audit/followup/}.

The library \texttt{ZetaS} has $206$ files and $22{,}751$ lines. It contains $3$ \texttt{sorry}, all in the open node
\texttt{ZetaS/Skeleton/A8g\_SigmaDistGeneral.lean}: Theorems~\ref{thm:OLL} and~\ref{thm:SigmaD} for arbitrary certificate data. That node
is used only by the forms with a general majorant threshold, \texttt{ZetaS.Top.zeta\_simple\_K5\_of\_majorant} and its three companions
(\texttt{ZetaS/Top/TopSolutionGeneral.lean}), and by no headline theorem; the import closure of the fourteen theorems of record contains
no \texttt{sorry}. There is no \texttt{axiom}, \texttt{admit} or \texttt{native\_decide} in the library.

\subsection{Repository layout}\label{app:layout}
The formalisation is published in a modified copy of the Lean repository of~\cite{AF} (\texttt{zeta-23-lean}): \repourl, at tag
\repotag. The commit of the tag is not printed in this paper; it is recorded in the file \texttt{RELEASE\_NOTES.md} of the repository and in
the Zenodo record of the release. The repository is a snapshot of the release without its development history: the commits named in this paper
(such as \texttt{bb7223b}) are commits of that history. The parts that concern this paper:
\begin{itemize}
\item \texttt{ZetaS/} with its root module \texttt{ZetaS.lean}: the library of this paper, outside the default build targets (its folders are
listed below); \texttt{ZetaSCertShims/}: three forwarding modules through which the checker files are imported by their bare names, so
that they are kept byte for byte as checked;
\item \texttt{ZetaSReplay/}: the $1{,}478$ files of the kernel replay of \texttt{CertAM5} ($1{,}475$ generated shards and the three common
modules \texttt{ReplayCommon}, \texttt{ReplayAM5} and \texttt{ReplayHeadlines}), in one flat folder, as a library outside the default build
targets (\S\ref{ssec:replay}, \S\ref{app:replay});
\item \texttt{comparator/}: the comparator set-up of~\cite{AF}, with seven new files for this paper: \texttt{ChallengeDeps/FollowUpZeta.lean} and
\texttt{Challenge/FollowUpZeta.lean} (quoted above), \texttt{Solution/FollowUpZeta.lean}, \texttt{PrintAxioms/FollowUpZeta.lean} and
\texttt{PrintAxioms/FollowUpReplay.lean}, \texttt{config-followup.json} and \texttt{README\_followup.md};
\item \texttt{supplementary/}: the data, certifiers and logs of Table~\ref{tab:certs} and the generator of the replay files
(\S\ref{ssec:reproduce});
\item \texttt{papers/zeta/}: the sources and the PDF of this paper; \texttt{audit/followup/}: the logs and \texttt{\#print axioms} outputs of
the builds of the release.
\end{itemize}
The repository also contains the library \texttt{ZetaShell/} of a companion paper on families of Dirichlet $L$-functions, whose sources are in
\texttt{papers/family/}, and the library \texttt{ZetaQ/} added to the repository earlier (see \texttt{NOTICE}); \texttt{ZetaS} imports neither.
The statement files of~\cite{AF} (\texttt{Zeta23/Statement.lean} and the existing files of \texttt{comparator/}) are unchanged. Further,
the repository contains changes to $25$ files of the library \texttt{Zeta23} ($3{,}122$ lines added, $121$ deleted, between the upstream base
commit and \texttt{bb7223b}): a weaker parameter condition
(\texttt{Params.ValidQ} in place of \texttt{Params.Valid}), with the proofs adapted, and new modules for Gevrey ramps and a free buffer
(\texttt{Tail/GevreyTail}, \texttt{Taper/GevreyProduct}, \texttt{WindowD}). The axiom census above is for the tree with these changes.
The Lean files of \texttt{Zeta23}, \texttt{ZetaS}, \texttt{ZetaSCertShims} and \texttt{comparator/} in the release are those of commit
\texttt{bb7223b}, at which the census of \S\ref{app:census} was taken and against which the replay ran, except for comments: for the release,
the copyright headers of the four new comparator files of commit \texttt{bb7223b} and stale doc-strings in
\texttt{comparator/ChallengeDeps/FollowUpZeta.lean} and \texttt{ZetaS/ChallengeZetaS.lean} were corrected. Two files were added to
\texttt{comparator/PrintAxioms/}, \texttt{FollowUpReplay.lean} (\S\ref{app:census}) and \texttt{FollowUpShell.lean}, for the companion paper,
and \texttt{comparator/README\_followup.md} was brought up to date. The comparator package \texttt{FollowUpFamily} of the companion paper
(\texttt{ChallengeDeps/}, \texttt{Challenge/}, \texttt{Solution/} and \texttt{PrintAxioms/FollowUpFamily.lean}, \texttt{config-followup-family.json}
and \texttt{README\_followup\_family.md}) was also added.
The folders of \texttt{ZetaS}:
\begin{itemize}
\item \texttt{ZetaS/Interfaces.lean}, \texttt{InterfacesV2.lean}: the interfaces between the zero side and the local arguments (Gram data,
frame families at bandwidth $\lambda<1$); \texttt{ZetaS/ChallengeZetaS.lean}: the library's copies of the definitions of \S\ref{app:defs},
together with \texttt{MajorantCert} and \texttt{CertMaj};
\item \texttt{LinAlg/}: Section~\ref{sec:linalg} (stability inequality and its two-parameter form, $\Phi_m$, pinching, monotonicity, pair
removal, localisation, inertia);
\item \texttt{Window/}: the two windows, their exact moments and constants, and the window-general pipeline at bandwidth $\lambda<1$;
\texttt{Bandwidth/}: the limit $\lambda\to1^-$;
\item \texttt{ZeroSide/}: the Poisson and Gram identities, positive definiteness, truncation defect, inertia of pair blocks, and the assembly
from the zeros of \texttt{riemannZeta};
\item \texttt{SigmaDist/}: Section~\ref{sec:simple}; \texttt{KSide/}: Sections~\ref{sec:first} and~\ref{sec:sc}; \texttt{Top/}: the numerics of
the constants, the robustness conditions, the headline theorems (\texttt{TopFinal.lean}, \texttt{TopSSC.lean}) and the forms with a general
majorant threshold (\texttt{TopSolutionGeneral.lean});
\item \texttt{Majorant/}: the majorant, its analytic part, its cell checker, the $13$ replay shards and the two majorant theorems;
\texttt{Cert/}: the corrected checker for the $K=5$ certificate (\texttt{CheckerBase.lean}, \texttt{CheckerLeaves.lean},
\texttt{CheckerCoreV3.lean}, with kernel-checked negative controls in \texttt{CheckerTestsV3.lean}), the specification
\texttt{CertSpecAM5.lean}, the soundness nodes and the chain \texttt{ChainV3.lean};
\item \texttt{Skeleton/}: re-export files for the proved nodes, and the one open node \texttt{A8g\_SigmaDistGeneral.lean}.
\end{itemize}
The new Lean files of \texttt{comparator/} carry the author's copyright notice under the Apache License~2.0, the licence of the repository
of~\cite{AF}; the files that come from~\cite{AF} keep their own notices (see \texttt{NOTICE}).

\par}
\section{How the work was done}\label{app:how}

\emph{Disclosure of generative-AI use.} This paper and its Lean development were produced substantially by a generative AI system, Claude
(Anthropic), working under the direction of the named author, who takes full responsibility for its contents. No AI system is listed as an author.

The local inequalities of Section~\ref{sec:cert} were found by linear programming over gap vectors located by floating-point search and
verified by interval branch-and-bound; the certificates were re-run, some box for box, and their inputs and closing constants were checked
again with independently written code. What is checked by the Lean kernel, and under which hypotheses, is stated in \S\ref{ssec:formal}. The
Lean libraries, the certificate data with their certifiers and logs, the generator of the replay files, and the build logs and
\texttt{\#print axioms} outputs of the release are in the repository \repourl{} at tag \repotag{} (Appendix~\ref{app:layout}).

\emph{Scope of the verification.} Theorems~\ref{thm:main} and~\ref{thm:S} and Corollary~\ref{cor:SC} are Lean theorems that depend only on
the three standard axioms. For the $K=5$ pair of Theorem~\ref{thm:main} the local certificate is replayed in the Lean kernel, so these two
statements have no hypothesis. For the $K=7$ pair, Theorem~\ref{thm:S} and Corollary~\ref{cor:SC} the local certificate is a displayed
hypothesis: a finite inequality checked outside Lean by interval arithmetic. The paper has not been refereed, and the Lean development has
not been rebuilt outside this project.


\bibliographystyle{plain}
\bibliography{bib}
\end{document}
